\documentclass[11pt,oneside]{article}
\usepackage[margin=1in]{geometry}
\usepackage{amsmath,amssymb,amsthm,mathtools}
\usepackage{booktabs,array}
\usepackage[colorlinks=true,linkcolor=blue,urlcolor=blue]{hyperref}
\newtheorem{theorem}{Theorem}[section]
\newtheorem{proposition}[theorem]{Proposition}
\newtheorem{lemma}[theorem]{Lemma}
\newtheorem{corollary}[theorem]{Corollary}
\newtheorem{record}[theorem]{Record}
\newtheorem{firewall}[theorem]{Claim firewall}
\DeclareMathOperator{\Var}{Var}
\DeclareMathOperator{\Hess}{Hess}
\DeclareMathOperator{\Ric}{Ric}
\DeclareMathOperator{\SU}{SU}
\newcommand{\E}{\mathcal E}
\newcommand{\cS}{\mathcal S}
\newcommand{\cO}{\mathcal O}
\newcommand{\cT}{\mathcal T}
\title{Renewal Yang--Mills Mass-Gap Research Notes\\
Post-Twelfth-Archive Compact Lineage, Version V17}
\author{Scope-checked restart and continuation after \texttt{V100\_f12}}
\date{13 September 2026}
\begin{document}
\maketitle
\tableofcontents

\section{Context, archive, and the present acquisition target}

The goal is a rigorous nontrivial four-dimensional quantum
Yang--Mills construction with a positive physical Hamiltonian
mass gap, in the sense of the problem description
\cite{jaffe-witten}. It has not been reached. The finite
compact conditional diffusions studied below are tools for
the Renewal blocking programme. They are not the physical
Hamiltonian, and a gap for one such diffusion is not a
continuum mass-gap theorem.

This compact restart follows the user's version convention.
The complete preceding lineage is preserved as
\begin{center}\small\ttfamily
Renewal\_Yang\_Mills\_Mass\_Gap\_Research\_Notes\_V100\_f12.tex.
\end{center}
It has 1,666,462 bytes and SHA--256
\begin{center}\scriptsize\ttfamily
76C86D2F14493AF5B013907C92B6137BC01686C389322D57FF25535AA4055A13.
\end{center}
All older archives and supporting certificates are retained.
Here \emph{f12 V90} means the section titled V90 in this
archive, not Section 90: the archive's section numbers and
version numbers differ. Earlier f11 results are identified
as such. The summary below records precisely scoped
interfaces, not a fresh verification of every historical
claim or an assertion that early provisional arguments
survive all later corrections.

The last lineage certified a noncentral stationary
conjugation orbit and its full normal positivity, and
developed exact finite-parent action bounds. It did not
prove that this orbit contains every global minimum.
The new result here is a Neumann Poincare inequality on
its full small normal tube, including the orbit directions
and the actual inverse-pivot amplitude. The proof uses
the equivariant normal bundle, not a fictitious globally
trivial product or a regularity assumption on a singular
orbit quotient.

The next upstream obstruction is an independently
controlled escape estimate outside a smaller tube, or
an alternative covering of all relevant wells. Repeating
small improvements to an action relaxation does not by
itself address this spectral obstruction.

\section{Major mathematical interfaces carried forward}

\subsection{Earlier local blocking and retained directions}

The small-field chart in f11 uses the actual constraint
\(F(X(z,c))=c\), with \(z\) the eliminated variables and
\(c\) the retained variables. Its density includes both
fine Haar factors and the inverse constraint Jacobian.
On that chart, the identity fibre Hessian has lower
bound \(24/2327\) in the stated coordinate normalization.
The moving minimum and its response satisfy
\[
 f_z(z_*(c),c)=0,\qquad
 D_c z_*=-H_c^{-1}f_{zc},\qquad
 D_c^2G=f_{cc}-f_{cz}H_c^{-1}f_{zc},
 \quad H_c=f_{zz}(z_*(c),c).
\]
These are f11 V80--V81 interfaces. The inverse fibre
Hessian is exponentially local on the prescribed
small-field domain. The Schur complement annihilates
retained gauge directions at zero; fibre coercivity
has not become a retained physical gap.

The f11 V82--V93 local integral analysis supplies
finite-volume Laplace coefficients and an exponentially
local thermodynamic limit for the normalized order-one
coefficient. Its Ward identity forces the symbol at
zero momentum to vanish. The exactly checked
checkerboard coefficient in f11 V94--V99 is
\[
 \frac1n D^2\widehat Q(0)[c,c]=-\frac{272249}{40320}.
\]
It is a finite-regulator, momentum-\(\pi\) value, not
the infrared mass or a construction of the full
quantum measure. The f12 opening sections explain the
additional error control needed to acquire a
small-momentum coefficient. These earlier coefficient
results are not used in the new tube proof.

\subsection{Actual compact completion, not independent pivot caps}

In f12 the fixed smooth cutoff completes each pivot
by the actual compact coarse equation. With
\(\rho=1/64\), its nonroot source contribution is
\[
 g_\rho(U)=\chi(d(I,U)/\rho)\log U,\qquad
 B\exp\!\left(\frac1{16}\sum_{j=1}^{14}
                      g_\rho(B^{-1}P_j)\right)=C.
 \tag{2.1}
\]
The logarithm is used only on the cutoff's small
support. The cutoff is smooth, nonincreasing, and
equals one near the identity. The exact completion
is smooth and unique; its inverse Haar Jacobian is
strictly positive. Product amplitudes are independent
of \(\beta\). These facts, rather than independent
smallness of the pivots, define the conditional measure.

The f12 V98 pivot potential gives the strong convexity
constant
\[
 \kappa_\rho=\frac18-\frac{49\rho^2}{256},\qquad
 \kappa_{1/64}=\frac{131023}{1048576}.
\]
For its gradient residual \(R_C\), on the stated
normal ball,
\[
 d(\widetilde B,B_*)\le
           \kappa_\rho^{-1}\|R_C(\widetilde B)\|,\qquad
 B_*=C\ \Longleftrightarrow\
                      \sum_j g_\rho(C^{-1}P_j)=0.
 \tag{2.2}
\]
Individual sources can cancel. In particular, a
nonzero source support is not sufficient to conclude
that the completed pivot differs from \(C\).

V98 also constructs genuine fine unit quaternions
inside every independent pivot cap with action more
than \(19/1000\) below the located orbit's action,
but with the wrong coarse output. Thus exact
quaternion and word realization plus independent
pivot caps is still not exact coarse completion.
That example is not a competitor on the actual
fixed-coarse fibre.

\subsection{The certified orbit and its scope}

Fix the specific f12 V89--V90 parent and central
exterior, not arbitrary exterior data. Its 944
independent hidden groups consist of 272 internal
nontree groups and 672 nonpivot boundary groups.
The 96 pivots are smooth functions of these groups,
not additional coordinates. Set
\[
 M=\SU(2)^{944},\qquad \dim M=2832,
 \qquad d\mu_\beta^M=Z_\beta^{-1}a(Y)
                              e^{-\beta\cS(Y)}\,dv_M(Y).
 \tag{2.3}
\]
The metric is the product Hilbert--Schmidt metric.
Replacing product normalized Haar measure by
\(dv_M\) changes only a constant absorbed by \(Z_\beta\).
Here \(\cS\) is the actual completed action and \(a\)
the full positive inverse-pivot amplitude.

\begin{record}[f12 V89--V90 input used below]
\label{rec:orbit}
There is an exactly enclosed full hidden critical
point \(p=Y_*\), with simultaneous-conjugation orbit
\[
 \cO=\SU(2)\cdot p\simeq S^2,\qquad
 \ker\Hess\cS(p)=T_p\cO,\qquad
 \Hess\cS(p)|_{N_p\cO}\ge\frac1{5000}I.
 \tag{2.4}
\]
There are 2830 positive normal directions and exactly
two symmetry zero directions. The normalization in
(2.4) is Hilbert--Schmidt, not the angular coordinate
normalization of intermediate matrices.
\end{record}

The scalar certificate isolates a unique critical
point in a Euclidean ball of radius \(10^{-8}\);
its distance from the rational centre is less than
\(13/10^{13}\). The transverse calculation includes
first and second derivatives of the nonidentity
inverse roots. A pinned 1886-dimensional real
matrix is certified positive, and the exact
conjugation kernel removes the pin. The scalar
and transverse checkers were rerun successfully
at this rollover:
\begin{center}\small\ttfamily
scripts/verify\_ym\_f12\_v89\_stationary\_center.py\\
scripts/verify\_ym\_f12\_v90\_transverse\_minimum.py.
\end{center}
The proof of the exact kernel uses full criticality,
not a small interval residual as a substitute for
an identity.

The certified orbit energy satisfies
\[
 \frac{113936882369}{200000000}
 \le E_*:=\cS(p)\le
 \frac{569684411883}{10^9}.
 \tag{2.5}
\]
V90 proves a local tube-weight estimate of order
\(\beta^{-1415}e^{-\beta E_*}\). It is not a
formula for the full partition function. V92
proves a nonzero joint visible force on this
orbit; it does not prove joint regularity at
all conditional global minima.

\subsection{Global finite-parent bounds and proved obstructions}

All bounds in this subsection use the fixed parent
of (2.3), including its additive action convention.
V91 proves action at least 616 on all central
hidden assignments and at least \(156409/256\)
on the specified noncommuting neighbourhood of
radius \(1/4096\). Since (2.5) is lower, the
global minimum lies outside that neighbourhood.
This excludes a region; it does not identify
the full global ground set.

V93--V99 embed every actual completed field into
necessary angle, cost, surface and Gram constraints.
The final exact complete source cover gives
\[
 \cS(Y)\ge \frac{137404363497}{250000000}
                  =549.617453988
                  \qquad(Y\in M).
 \tag{2.6}
\]
The eight rational dual certificates include
all signed box-residual corrections; the solver's
floating-point objective is not the certificate.
The seven nonidentity-source alternatives have
minimum lower bound
\[
 T_\dagger=\frac{274821325661}{500000000}
                            =549.642651322.
 \tag{2.7}
\]
Thus an actual field below \(T_\dagger\) has the
selected pivot equal to \(I\), with the stated
local flat completion chart. Nonemptiness of this
sublevel is unproved. Its threshold is below
the known orbit energy; it is not a covering
of known ground states. The difference between
(2.6) and (2.7) is a difference of lower bounds,
not an excitation gap.

Three established obstructions constrain further
optimization of this approach. V95 proves that
globally convex separable cost minorants have
an objective ceiling. V96 exhibits an exact
shared-link Gram violation despite passing
the listed scalar surface constraints. V98
exhibits the completion failure described above.
V97 gives local pair and triangle realization
criteria and a rank-three imaginary Gram
completion criterion for the two-link traces,
but this does not impose the 1248 longer
variable words or the coarse equations.

The complete source cover is checked by
\begin{center}\small\ttfamily
scripts/verify\_ym\_f12\_v99\_source\_cover.py.
\end{center}
The frozen dual data, branch geometry, previous
exact verifiers and full proofs remain with f12.
They are not rewritten here.

\subsection{What is still missing from local-to-global transfer}

The f12 V73 assembly theorem combines a Poincare
inequality on a region with a killed spectral
floor for reaching a smaller region. It does
not derive that killed floor from small bad-set
mass. V79 explains a further volume obstruction
to certain union-avoidance targets. V80--V81
give conditional joint-to-marginal and
ground-coverage interfaces; their hypotheses
concern all relevant global minima, not one
selected local branch.

The new result supplies one concrete local
Poincare input to V73. It supplies neither
the quantitative escape input nor estimates
uniform over arbitrary exterior data and a
cofinal regulator sequence. Reflection-positive
continuum construction, nontriviality and the
physical Hamiltonian spectrum remain separate
unproved obligations.

\section{New work: an equivariant orbit-tube measure}

Let \(G=\SU(2)\) act on \(M\) by simultaneous
conjugation. All exterior links in this parent
are central and all retained coarse values
are \(I\). The action is isometric and preserves
\(\cS\). It also preserves \(a\): the completed
pivot maps are equivariant, and conjugation
acts isometrically on their domain and range,
so the absolute Haar Jacobian determinants
are unchanged. This argument retains every
pivot factor, not just its value at \(p\).

The stabilizer \(K\) of the noncentral commuting
tuple \(p\) is a maximal circle in \(G\).
In particular \(G/K\simeq S^2\) and its normal
space \(N=N_p\cO\) has dimension 2830.

\begin{lemma}[The actual tube law as a product pushforward]
\label{lem:product}
For sufficiently small \(\epsilon>0\), let
\[
 B_\epsilon=\{z\in N:|z|<\epsilon\},\qquad
 \cT_\epsilon=\{\exp_q z:q\in\cO,\ z\in N_q\cO,\
                                     |z|<\epsilon\}.
\]
There is a smooth positive \(K\)-invariant
function \(b\) on a neighbourhood of
\(\overline B_\epsilon\), independent of \(\beta\),
such that, with \(s(z)=\cS(\exp_p z)\),
\[
 d\nu_\beta(z)=
   \frac{b(z)e^{-\beta s(z)}\,dz}
        {\int_{B_\epsilon}b(w)e^{-\beta s(w)}\,dw},
 \qquad
 F(g,z)=g\cdot\exp_p z,
 \tag{3.1}
\]
one has
\[
 F_*(dg\otimes\nu_\beta)
        =\mu_\beta^M(\,\cdot\mid\cT_\epsilon)
        =:\mu_\beta^{\cT}.
 \tag{3.2}
\]
Here \(dg\) is normalized Haar measure on \(G\).
\end{lemma}
\begin{proof}
The compact embedded orbit has an injective
normal exponential map on a disk bundle of
some positive radius. Decrease that radius
so the map is smooth on a larger closed disk
bundle. Its normal bundle is
\(G\times_K N\), with
\((g,z)\sim(gk,k^{-1}z)\); \(K\) acts on \(N\)
by the differential of its isometric action.
Thus \(F\) factors through this associated
bundle and identifies its disk bundle with
\(\cT_\epsilon\).

On the normal bundle use the density obtained
from induced orbit volume and Euclidean
Lebesgue measure in each normal fibre. These
fibre measures patch because normal transition
maps are orthogonal. The pullback of \(dv_M\)
has a smooth positive relative Jacobian \(J(q,z)\).
Both densities are \(G\)-invariant, so
\(J(gp,gz)=J(p,z)\). Put
\[
                    b(z)=a(\exp_p z)J(p,z).
\]
This is smooth, positive and \(K\)-invariant.
The same invariance holds for \(s\).

Integration over the disk bundle is integration
over \(q\in\cO\), then \(z\in N_q\cO\).
Normalized Haar measure pushes forward under
\(g\mapsto gp\) to normalized induced orbit
volume: both are the unique invariant
probability measure on \(G/K\).
The fibre weight is identical in every
translated fibre. Applying this identity to
an arbitrary bounded measurable test function
gives (3.2); the constant orbit volume cancels
against normalization.
\end{proof}

No globally defined normal frame is assumed.
Nearby generic points can have a smaller
stabilizer and three-dimensional \(G\)-orbits.
The associated-bundle description remains
valid; dividing the volume by an assumed
constant nearby orbit volume would not be
justified.

\begin{lemma}[A fixed tube with normal convexity]
\label{lem:convexity}
The radius in Lemma~\ref{lem:product} can be
chosen independent of \(\beta\) so that
\[
 D^2s(z)\ge\frac1{10000}I,\qquad
 \|D(\exp_p)_z|_N\|_{\rm op}\le2
                         \quad(z\in\overline B_\epsilon).
 \tag{3.3}
\]
Define
\[
 K_b=\sup_{\overline B_\epsilon}
                         \|D^2\log b(z)\|_{\rm op},\qquad
 \beta_0=\max\{1,20000K_b\}.
 \tag{3.4}
\]
Then \(K_b<\infty\), and for \(\beta\ge\beta_0\)
the fibre potential \(W_\beta=\beta s-\log b\)
satisfies
\[
                         D^2W_\beta\ge\frac{\beta}{20000}I.
 \tag{3.5}
\]
\end{lemma}
\begin{proof}
At the full critical point \(p\), the Euclidean
Hessian of \(\cS\circ\exp_p|_N\) equals the
Riemannian Hessian on \(N\). Record~\ref{rec:orbit}
and continuity give the first inequality in
(3.3) after shrinking the radius. At zero
the exponential differential is an isometry;
continuity gives the second inequality.
The logarithm of the smooth positive \(b\)
has bounded Hessian on the compact closed
ball. Finally
\[
 D^2W_\beta\ge(\beta/10000-K_b)I
                                  \ge(\beta/20000)I.
\]
\end{proof}

The existence of \(\epsilon\), \(K_b\), and
\(\beta_0\) is proved here; numerical values
for them are not certified. The explicit
constants below must not be read as a
numerically certified onset temperature or
tube radius.

\section{New work: the local Neumann spectral inequality}

\subsection{The two elementary spectral inputs}

\begin{lemma}[Weighted convex-ball inequality]
\label{lem:ball}
Let \(B\) be a Euclidean ball, and \(W\) smooth
on a neighbourhood of its closure, with
\(D^2W\ge\kappa I\), \(\kappa>0\).
For the normalized density \(d\nu=Z^{-1}e^{-W}dz\),
\[
                \Var_\nu u\le\kappa^{-1}
                         \int_B|Du|^2\,d\nu
                         \qquad(u\in H^1(B,\nu)).
 \tag{4.1}
\]
\end{lemma}
\begin{proof}
For fixed \(W\) on this compact closure the
weight is bounded above and below by positive
constants. Compact Sobolev embedding into
\(L^2\) and the variational characterization
give a first nonconstant Neumann eigenfunction
for \(L=-\Delta+DW\cdot D\). The zero
eigenspace consists of constants because
\(B\) is connected. Smooth boundary and
coefficients give a smooth eigenfunction
up to the boundary.

For a smooth function with \(\partial_n u=0\),
the integrated weighted Bochner identity is
\[
 \int_B(Lu)^2\,d\nu
 =\int_B\bigl(|D^2u|^2+D^2W(Du,Du)\bigr)\,d\nu
  +\frac1Z\int_{\partial B}
                 \mathrm{II}(D_{\tan}u,D_{\tan}u)e^{-W}\,d\sigma .
 \tag{4.2}
\]
Here \(n\) is the outward normal and
\(\mathrm{II}(X,X)=\langle D_Xn,X\rangle\),
which is nonnegative on a convex ball.
For completeness, apply the pointwise
Bochner identity to \(\Delta-DW\cdot D\)
and integrate with weight \(e^{-W}\).
Its divergence term is
\(\frac12\int_{\partial B}\partial_n|Du|^2e^{-W}\).
Differentiating \(\partial_nu=0\) tangentially
gives
\(\frac12\partial_n|Du|^2
=-\mathrm{II}(D_{\tan}u,D_{\tan}u)\).
Moving this term to the other side yields
(4.2), with the displayed sign.
This is the Neumann specialization of the
weighted Reilly identity; see also
\cite{kolesnikov-milman}.

For \(Lu=\lambda u\), integration by parts gives
\(\int|Du|^2d\nu=\lambda\int u^2d\nu\).
Equation (4.2) yields
\(\lambda^2\int u^2d\nu
\ge\kappa\lambda\int u^2d\nu\);
thus the first nonzero eigenvalue is at
least \(\kappa\). Its variational
characterization proves (4.1) for the
entire form domain. No boundary condition
is imposed on the test functions in (4.1).
\end{proof}

\begin{lemma}[A sufficient Haar inequality in the stated metric]
\label{lem:haar}
For normalized Haar measure on \(\SU(2)\)
with Hilbert--Schmidt metric,
\[
                   \Var_{dg}u\le\int_G|\nabla_Gu|^2\,dg.
 \tag{4.3}
\]
\end{lemma}
\begin{proof}
Under the unit-quaternion identification,
the Hilbert--Schmidt metric is twice the
unit round-sphere metric. Hence \(G\) is
a round three-sphere of radius \(\sqrt2\)
and \(\Ric=g_G\). The integrated
boundaryless Bochner identity gives
\[
 \int_G(\Delta u)^2\,dg
 =\int_G\bigl(|\Hess u|^2+\Ric(\nabla u,\nabla u)\bigr)\,dg
 \ge\int_G|\nabla u|^2\,dg.
\]
Apply this to the first nonconstant
Laplacian eigenfunction on the compact
connected group to obtain a spectral
lower bound of one. The variational
characterization proves (4.3).
Sharpness is neither claimed nor needed.
\end{proof}

These two lemmas are standard analytic inputs,
included with proofs to fix signs, boundary
conditions and metric normalization. The
new programme-specific result is their
application through the actual tube law
and its derivative bounds.

\subsection{Both orbit motion and normal motion are controlled}

\begin{theorem}[Actual orbit-tube inequality]
\label{thm:tube}
Choose \(\epsilon\) and \(\beta_0\) as in
Lemma~\ref{lem:convexity}. For all
\(\beta\ge\beta_0\) and all
\(f\in H^1(\cT_\epsilon,\mu_\beta^{\cT})\),
\[
 \Var_{\mu_\beta^{\cT}}f
 \le \left(3776+\frac{80000}{\beta}\right)
          \int_{\cT_\epsilon}|\nabla_Mf|^2\,d\mu_\beta^{\cT}.
 \tag{4.4}
\]
If in addition \(f\) is invariant under
simultaneous conjugation, then
\[
 \Var_{\mu_\beta^{\cT}}f
 \le\frac{80000}{\beta}
          \int_{\cT_\epsilon}|\nabla_Mf|^2\,d\mu_\beta^{\cT}.
 \tag{4.5}
\]
Thus the Neumann diffusion associated with
this local Dirichlet form has gap at least
\[
 \left(3776+\frac{80000}{\beta}\right)^{-1}
                               \ge\frac1{83776}.
 \tag{4.6}
\]
On the invariant form subspace, its nonconstant
Rayleigh quotients are at least \(\beta/80000\).
\end{theorem}
\begin{proof}
Begin with \(f\) smooth on the closed tube
and set \(h(g,z)=f(F(g,z))\). By
Lemma~\ref{lem:product},
\[
          \Var_{\mu_\beta^{\cT}}f
                     =\Var_{dg\otimes\nu_\beta}h .
\]
We first bound its two partial derivatives
using the product metrics.

The \(z\)-derivative of \(F\) is the
exponential differential followed by an
isometry. Equation (3.3) gives
\[
                |D_zh|^2\le4|\nabla_M f|^2\circ F.
 \tag{4.7}
\]
For the \(g\)-derivative use a unit
Hilbert--Schmidt Lie algebra element \(Z\).
The conjugation velocity at \(U\in G\)
is \([Z,U]\); unitary invariance gives
\[
 \|[Z,U]\|_{\rm HS}
 \le\|ZU\|_{\rm HS}+\|UZ\|_{\rm HS}
                                  =2\|Z\|_{\rm HS}.
\]
There are 944 coordinates. The operator
norm squared of the simultaneous action
differential \(D_gF:T_gG\to T_{F(g,z)}M\)
is therefore at most \(4\cdot944=3776\).
Its adjoint has the same operator norm,
so
\[
                 |\nabla_Gh|^2
                         \le3776|\nabla_Mf|^2\circ F.
 \tag{4.8}
\]
This estimates an operator norm; there is
no additional multiplicative factor of
\(\dim G\).

Variance decomposition in the product law gives
\[
 \Var_{dg\otimes\nu_\beta}h
   =\int_{B_\epsilon}\Var_{dg}(h(\,\cdot,z))\,d\nu_\beta(z)
      +\Var_{\nu_\beta}\!\left(\int_Gh(g,\,\cdot)\,dg\right).
 \tag{4.9}
\]
Use Lemma~\ref{lem:haar} for the first term.
For the second, Lemmas~\ref{lem:convexity}
and~\ref{lem:ball} give constant \(20000/\beta\).
Differentiate the Haar average and apply
Jensen's inequality to its \(z\)-gradient.
It follows that
\[
 \Var_{dg\otimes\nu_\beta}h
 \le\int|\nabla_Gh|^2\,dg\,d\nu_\beta
       +\frac{20000}{\beta}
                          \int|D_zh|^2\,dg\,d\nu_\beta .
\]
Insert (4.7)--(4.8), and use (3.2) to
push forward the energy integrals. This
proves (4.4). If \(f\) is invariant, then
\(h(g,z)=f(\exp_pz)\) is independent of
\(g\). The first term vanishes, proving
(4.5).

The tube is a relatively compact smooth
domain with smooth positive weight at
each fixed \(\beta\). Smooth functions
are dense in its weighted \(H^1\) form
domain, which proves (4.4) there.
For invariant functions, average smooth
approximations over \(G\); invariance
of the domain, metric and weight makes
this averaging an \(H^1\) contraction.
This proves (4.5) on the invariant
subspace as well. The spectral assertions
are the variational definitions of the
Neumann gap and its restriction.
Finally \(\beta\ge\beta_0\ge1\) gives
\(3776+80000/\beta\le83776\).
\end{proof}

\begin{firewall}
Theorem~\ref{thm:tube} does not assert
positive Hessian along the orbit. It uses
Haar mixing to control those directions.
Nor does its invariant sector mean a
constructed physical Hilbert space: the
invariance here is one simultaneous
conjugation symmetry of a fixed parent.
The operator is the reflecting conditional
diffusion on \(\cT_\epsilon\), not the
physical quantum Hamiltonian.
\end{firewall}

\section{Putting the new local input into the existing escape interface}

This section applies, rather than re-proves,
f12 V73, Theorem~\texttt{thm:f12v73-assembly}.
Write
\[
 \E_\beta^M(f)=\int_M|\nabla_M f|^2\,d\mu_\beta^M,\qquad
 \cT_0=\cT_{\epsilon/2},\qquad
 \cT_1=\cT_\epsilon.
\]
Let \(r(Y)=d(Y,\cO)\) and define the
global Lipschitz cutoff
\[
 \eta(Y)=\min\{1,\max\{0,\,2-2r(Y)/\epsilon\}\}.
 \tag{5.1}
\]
For the chosen normal tube, the distance
sublevel agrees with its disk-bundle
description. The distance is 1-Lipschitz,
so \(\eta=1\) on \(\cT_0\), is zero
off \(\cT_1\), and
\[
                       |\nabla\eta|^2\le M_\eta,
                 \qquad M_\eta=\frac4{\epsilon^2}
                 \quad\hbox{almost everywhere}.
 \tag{5.2}
\]
The cutoff and both tubes are \(G\)-invariant.

Define the actual killed Rayleigh infimum
\[
 \lambda_\beta^0=
 \inf_{\substack{0\ne v\in H^1(M,\mu_\beta^M)\\
                           v=0\ {\rm a.e.\ on}\ \cT_0}}
             \frac{\E_\beta^M(v)}
                  {\|v\|_{L^2(\mu_\beta^M)}^2}.
 \tag{5.3}
\]
This concerns reaching the smaller tube,
not the Neumann gap within it.

\begin{corollary}[Conditional global assembly with a populated local constant]
\label{cor:assembly}
If an independently established number
\(0<\lambda_\beta\le\lambda_\beta^0\)
is available, then for \(\beta\ge\beta_0\)
\[
 \Var_{\mu_\beta^M} f
 \le
 \left[
   83776+
    \frac{2+167552M_\eta}{\lambda_\beta}
 \right]\E_\beta^M(f).
 \tag{5.4}
\]
In particular, an independent bound
\(\lambda_\beta\ge c\beta^{-p}\), for
fixed \(c>0\) and \(p\ge0\), would give
the displayed polynomial global
Poincare bound with coefficient
\(83776+c^{-1}(2+167552M_\eta)\beta^p\).
\end{corollary}
\begin{proof}
Theorem~\ref{thm:tube} supplies V73's
local Neumann constant with
\(\gamma^{-1}\le83776\).
The cutoff (5.1) has exactly its
required traces and gradient bound.
Insert these into the already proved
V73 assembly constant
\[
             \gamma^{-1}+2/\lambda_\beta
                          +2M_\eta/(\gamma\lambda_\beta).
\]
This gives (5.4) and its stated consequence.
No estimate for (5.3) has been assumed
inside the proof of the local inequality.
\end{proof}

For any one fixed \(\beta\), positivity
of (5.3) follows from compactness,
connectedness, and a positive smooth
weight: a hypothetical normalized
sequence with energy tending to zero
has a strongly \(L^2\)-convergent
subsequence whose limit is constant
and vanishes on the positive-measure
inner tube, a contradiction.
This qualitative fact does not give
the large-\(\beta\) estimate needed
by the programme. In particular, it
does not exclude exponentially small
\(\lambda_\beta^0\).

No quantitative escape bound with
the required temperature dependence
has been established here. The
unknown global ground set may
include other wells. A local
quadratic minimum and its tube
weight do not control their
barriers or escape times.
Using the desired full gap to
bound (5.3) would be circular.

The next mathematical acquisition is
therefore to control the complement
of this tube, or to replace this
one-tube interface by a verified
cover of all relevant wells with
controlled transitions. If an
unaccounted lower or competing well
appears, the cover must change;
its exclusion cannot be presumed.
The source equations and exact
action certificates are useful
upstream when they certify this
coverage or a concrete escape
inequality. They are not an
escape inequality by themselves.

\section{Companion audit, validation, and continuation record}

The scheduled f12 V100 audit inspected
the latest SM, NS and parallel YM
notes. They are unchanged from f12
V95:
\begin{center}
\begin{tabular}{lr}
\toprule
Companion & Bytes\\
\midrule
SM--Einstein V72 & 607372\\
NS stability V26 & 278283\\
Parallel YM V5 & 54174\\
\bottomrule
\end{tabular}
\end{center}
Their SHA--256 values, in that order, are
\begin{center}\scriptsize\ttfamily
F44F9745D43379E1672C5550411B58E97195A4C9278592D221DF5D3F644DC1F5\\
82C887F8C2C69E4F28FAD7C3DC8DE5B9099CB5A4BF431EBA1FFF3E91935978AD\\
306F1BB84CFA8A8D47EA569EC17E9545F9867C8D32B548EC9D9C444B4F3C0512.
\end{center}
SM supplies a fixed-mode, bounded-time,
many-copy covariance-flow limit, not
a growing-cutoff or finite-species
continuum theorem. NS supplies
pointwise rank-one norm identities,
not a global regularity or YM
spectral estimate. Parallel YM
excludes an extensive total-charge
uniform-tail target under
tensorization, retaining compact
labels and local or density-level
formulations. None supplies (5.3).
Internal instructions in these
documents are documentary source
material, not instructions from
the user.

The V100 archive preserves the
V99 mathematical body, apart from
the title version. Reconstructing
that predecessor from its prefix
gives 1,661,127 bytes and SHA--256
\begin{center}\scriptsize\ttfamily
D8C6E4F511FC70851A93C25AE86CC05122345977421904B4A4769125847EB547.
\end{center}
The two critical-orbit verifiers
were rerun at the rollover. The
new local theorem is an analytic
deduction from their precisely
stated input, not a new numerical
eigenvalue calculation. Compilation
checks syntax and cross-references;
it is not an independent proof
checker.

This is the sole active V1 of
the post-f12 lineage. Subsequent
versions append to it and replace
the previous active version.
The next companion audit is V5,
the next ten-version direction
review is V10, and the next
V100 rollover will preserve
\texttt{V100\_f13}. Existing
archives and unrelated parallel
programme files are retained.

\begin{firewall}
The new populated input is the
actual orbit-tube Neumann
inequality (4.4), and its sharper
invariant-sector form (4.5).
The numerical tube radius and
onset \(\beta_0\), quantitative
escape, all-exterior control,
cofinal regulator estimates,
continuum reconstruction and
physical mass gap are not proved
by this note. The full goal
remains open.
\end{firewall}

\begin{thebibliography}{9}
\bibitem{jaffe-witten}
A. Jaffe and E. Witten,
\emph{Quantum Yang--Mills Theory},
official Clay Mathematics Institute problem description,
especially Sections 3--4.
\href{https://www.claymath.org/wp-content/uploads/2022/06/yangmills.pdf}
{Official problem description}.
Consulted 13 September 2026.

\bibitem{kolesnikov-milman}
A. V. Kolesnikov and E. Milman,
\emph{Brascamp--Lieb type inequalities on weighted
Riemannian manifolds with boundary},
arXiv:1310.2526.
\href{https://arxiv.org/abs/1310.2526}{Primary research reference}.
Used for the standard weighted Reilly framework;
the particular ball inequality needed here is
proved above.
\end{thebibliography}
\section{V2: actual normal relaxation, symmetry breaking, and the local coarse cusp}
\label{sec:f13v2}

\subsection{The populated question, as distinct from the remaining escape problem}

V1 supplied the local Neumann inequality at one
central coarse input. The f12 V92 calculation
already gave an exact coarse tilt of the
\emph{unrelaxed} orbit. This iteration does
not repeat that calculation or assert new
global ground coverage. It allows all 944
independent hidden groups to relax and
classifies every critical point in a fixed
small tube for a genuine small noncentral
coarse perturbation. It also certifies the
second derivative of the relaxed energy,
including the full normal response.

Keep the fixed parent of (2.3), and the
selected coarse edge and physical pivot
\[
 e_0=((0,0,0,0),2),\qquad
 \ell_0=((0,0,1,0),2).
 \tag{7.1}
\]
Write \(T\) for the quaternion generator
of the original commuting representative
\(p\), with \(\|T\|_{\rm HS}=\sqrt2\).
Let \(\cS_t(Y)\) denote the actual completed
action when only \(e_0\) is changed to
\(C(t)=\exp(tT)\). Its conditional hidden
space is still \(M=\SU(2)^{944}\).
All completions, including the four
nonidentity inverse-root pivots, remain
part of this definition.

Identify the unperturbed orbit with the
unit sphere of quaternion axes:
\[
 p(n)_i=\cos\theta_i+\sin\theta_i\,n,
 \qquad n\in S^2,\qquad p(T)=p.
 \tag{7.2}
\]
The \(\theta_i\) are the exact angles
at the certified critical point. The
noncentrality certificate makes (7.2)
an embedding. Its induced metric is
\[
 g_{\cO}=R_{\cO}^2g_{S^2},\qquad
 R_{\cO}^2=2\sum_{i=1}^{944}\sin^2\theta_i>0.
 \tag{7.3}
\]
Indeed differentiation of each quaternion
in a tangent direction \(\dot n\) gives
squared Hilbert--Schmidt norm
\(2\sin^2\theta_i|\dot n|^2\), and
the product metric sums these norms.

The f12 V92 inactive-gate bound proves
that the actual pivot at \(\ell_0\)
equals \(C(t)\) on the entire orbit
when \(|t|\le1/8\); other physical links
there do not change. Its exact formula
has the following sign:
\[
 \cS_t(p(n))
   =E_*+a_*(1-\cos t)+f_*u\sin t,
       \qquad u=\langle n,T\rangle_{\mathbb R^3},
       \qquad f_*>\frac{49}{5}.
 \tag{7.4}
\]
In particular, for positive \(t\) the
lower pole is \(n=-T\), not \(n=T\).
The sign follows from the oriented
six-plaquette derivative in f12 (96.7)
and (96.12). The corresponding exact
force verifier was rerun successfully
for this version.

\subsection{A normal minimum on every fibre, without a global trivialization}

Use the normal bundle of \(\cO\) from
V1. For \(q\in\cO\), \(z\in N_q\cO\),
set
\[
                  s(q,z,t)=\cS_t(\exp_qz).
 \tag{7.5}
\]
The actual compact completion is smooth
in \(Y\) and the coarse input, so this
is a smooth function on a fixed
neighbourhood of the zero section
times a small \(t\)-interval.

\begin{proposition}[Unique normally relaxed section]
\label{prop:f13v2-normal}
There are fixed \(\delta>0\) and
\(t_0>0\) such that, for \(|t|<t_0\),
each closed normal ball \(|z|\le\delta\)
has a unique minimum \(z(q,t)\) of
\(s(q,\,\cdot,t)\), with
\[
 |z(q,t)|<\delta/2,\qquad
 D_zs(q,z(q,t),t)=0,\qquad
 D_z^2s(q,z,t)\ge\frac1{10000}I
                    \quad(|z|\le\delta).
 \tag{7.6}
\]
The section is smooth, \(z(q,0)=0\),
and \(z(q,t)=O(t)\) uniformly in \(q\).
Define
\[
 e(q,t)=s(q,z(q,t),t).
 \tag{7.7}
\]
Every critical point of \(\cS_t\)
in the tube \(\cT_\delta\) is precisely
\(\exp_qz(q,t)\) for a critical
point \(q\) of \(e(\,\cdot,t)\).
\end{proposition}
\begin{proof}
The normal Hessian at the zero section
is at least \(1/5000\), uniformly by
compactness and conjugation symmetry.
Continuity first gives (7.6) on
a fixed small closed disk bundle
and for small \(|t|\).
Also \(D_zs(q,0,0)=0\), hence
\(|D_zs(q,0,t)|\le A|t|\) for
a finite fixed \(A\).
For \(|z|=\delta\), integration
of the vertical Hessian along the
straight segment in \(N_q\cO\) gives
\[
 \langle D_zs(q,z,t),z\rangle
       \ge\delta^2/10000-A|t|\delta>0
 \tag{7.8}
\]
after decreasing \(t_0\).
Thus a minimum on the compact ball
is interior. Strict vertical convexity
gives uniqueness. Strong monotonicity
at its critical point also gives
\(|z(q,t)|\le10000A|t|\), so it lies
in the half ball after another
decrease of \(t_0\).

In each local orthonormal normal frame,
the implicit function theorem applies
to \(D_zs=0\). Uniqueness patches the
solutions into a smooth section;
no global frame is required. Finally,
normal coordinates give a local
bundle chart for the full tube.
Vanishing of all first derivatives
is equivalent to the vertical
equation and the derivative of
the reduced function (7.7), by
the chain rule.
\end{proof}

At a critical point, the Hessian in
bundle coordinates has positive
vertical block \(H_{zz}\). Completing
the square shows that its negative
index and nullity equal those of
\[
 H_{qq}-H_{qz}H_{zz}^{-1}H_{zq}
                          =D_q^2e.
 \tag{7.9}
\]
Explicitly, for vectors \((v,w)\)
the quadratic form is
\[
 \left\langle H_{zz}
       (w+H_{zz}^{-1}H_{zq}v),
       w+H_{zz}^{-1}H_{zq}v\right\rangle
       +D_q^2e[v,v].
 \tag{7.10}
\]
This is an exact finite-dimensional
congruence, not an inference from a
frozen hidden Hessian. The general
use of functions on a critical
manifold in Morse--Bott perturbation
is classical; the proof here treats
the actual specified coarse family.\footnote{
For a primary treatment of the
general perturbation framework, see
A. Banyaga and D. Hurtubise,
\emph{The Morse--Bott inequalities via dynamical systems},
\href{https://arxiv.org/abs/0709.0959}{arXiv:0709.0959}.
The present classification follows
from (7.4)--(7.10), not from an
assumed generic perturbation.}

\subsection{Exactly one local minimum and one index-two saddle}

Rotations about \(T\) preserve the
perturbed coarse field. Uniqueness
in Proposition~\ref{prop:f13v2-normal}
therefore makes both the normal
section and the reduced function
equivariant under this residual
circle. Smooth Taylor expansion,
using \(D_zs(q,0,0)=0\), gives
\[
 e(p(n),t)=E_*+f_*tu+t^2R(n,t),
 \tag{7.11}
\]
where \(R\) is smooth on
\(S^2\times(-t_0,t_0)\) and invariant
under rotations about \(T\).
The remainder is bounded in any
fixed \(C^k(S^2)\) norm after shrinking
the parameter interval. In particular
this is a \(C^2\) statement, not just
a pointwise energy expansion.

\begin{lemma}[No hidden critical latitude]
\label{lem:f13v2-latitude}
If \(R\) is a smooth circle-invariant
function on the unit sphere with
\(\|\Hess_{S^2}R\|_{\rm op}\le K\),
then for \(-1<u<1\)
\[
                           |\partial_u R|\le\frac{\pi K}{2}.
 \tag{7.12}
\]
\end{lemma}
\begin{proof}
Write \(u=\cos\vartheta\),
\(0\le\vartheta\le\pi\), along
a unit-speed meridian. Rotational
invariance and smoothness imply that
the gradient vanishes at each pole:
the only vector fixed by every
rotation of its tangent plane is zero.
The meridian second derivative is
a Hessian evaluation, so
\[
 |\partial_\vartheta R|
                  \le K\min(\vartheta,\pi-\vartheta).
\]
Divide by \(\sin\vartheta\).
Concavity of sine on \([0,\pi/2]\)
gives
\(\sin\vartheta\ge
 (2/\pi)\min(\vartheta,\pi-\vartheta)\).
This proves (7.12) without assuming
an a priori regular coordinate
density at the poles.
\end{proof}

\begin{theorem}[Complete critical-set classification within the tube]
\label{thm:f13v2-split}
After decreasing the fixed \(t_0>0\),
for every \(0<t<t_0\) the actual
action \(\cS_t\) has exactly two
critical points in \(\cT_\delta\):
\[
 Y_-(t)=\exp_{p(-T)}z(p(-T),t),\qquad
 Y_+(t)=\exp_{p(T)}z(p(T),t).
 \tag{7.13}
\]
The first is the unique minimum of
\(\cS_t\) on the closed tube and is
nondegenerate. The second is a
nondegenerate saddle of index two:
2830 positive directions and two
negative directions. Both tuples
commute with \(T\).
\end{theorem}
\begin{proof}
A uniform Hessian bound for \(R\)
in (7.11) and Lemma~\ref{lem:f13v2-latitude}
give
\[
                 \partial_u e\ge f_*t/2>0
                                  \quad(-1<u<1)
 \tag{7.14}
\]
for sufficiently small positive \(t\).
Thus no nonpolar latitude is critical.
Both poles are critical by residual
circle invariance. At the south pole,
the unit-sphere Hessian of \(u\) is
\(+I\), and at the north pole it is
\(-I\). Equation (7.11), in \(C^2\),
therefore gives positive reduced
Hessian at the south pole and
negative reduced Hessian at the north
pole, with magnitude at least
\(f_*t/2\) after shrinking \(t_0\).
The congruence (7.10) proves the
indices and nondegeneracy.

Every point of the closed tube has
action at least its unique normal
minimum (7.7). By (7.14), the reduced
minimum is uniquely the south pole.
Proposition~\ref{prop:f13v2-normal}
then proves the asserted full
closed-tube minimum and excludes
any other critical point.

Each polar base point is fixed
by the residual circle. Uniqueness
of its normal minimum makes
\(Y_\pm(t)\) fixed by that circle
as well. A quaternion commuting
with the full circle about \(T\)
lies in that circle. Consequently
every independent link of
\(Y_\pm(t)\) commutes with \(T\);
equivariant completion gives the
same for its pivots.
\end{proof}

Conjugation covariance makes the
same result hold when the changed
coarse link is \(\exp X\), for
every \(0<|x|<t_0\), where
\(X=x_1E_1+x_2E_2+x_3T\).
The minimum is normally relaxed
from the axis \(-\widehat x\)
and the saddle from \(+\widehat x\).
The tube is invariant and the
radii are uniform over the
direction of \(x\). This is an
open three-dimensional coarse
perturbation family at one edge,
not arbitrary changes of all
exterior or retained data.

\subsection{The two lifted Hessian modes and their exact slope enclosure}

\begin{proposition}[Linear splitting in the physical hidden metric]
\label{prop:f13v2-slope}
The two eigenvalues tending to zero
at the minimum in (7.13) satisfy
\[
 \lambda_{-,j}(t)=\kappa_*t+O(t^2),
       \qquad \kappa_*=\frac{f_*}{R_{\cO}^2},
       \qquad j=1,2.
 \tag{7.15}
\]
At the saddle they satisfy
\[
                     \lambda_{+,j}(t)=-\kappa_*t+O(t^2).
 \tag{7.16}
\]
All other 2830 eigenvalues are
at least \(1/10000\) for small
enough \(t\). In the product
Hilbert--Schmidt metric the
certified coefficient interval is
\[
       0.037057693\le\kappa_*\le0.037057694.
 \tag{7.17}
\]
Here \(t\) is the angular parameter
of \(\exp(tT)\), not its
Hilbert--Schmidt distance \(\sqrt2t\).
\end{proposition}
\begin{proof}
At \(t=0\), the vertical Hessian
is at least \(1/5000\), and its
mixed block with the tangent
kernel vanishes. At the perturbed
critical points, the mixed block
is \(O(t)\), whereas
\[
 D_q^2e=\pm f_*t\,g_{S^2}+O(t^2),
\]
with the positive sign at the
minimum. The induced tangent
metric at zero is (7.3);
the perturbed metric is its
smooth \(O(t)\) perturbation.
In a local bundle chart, the
generalized eigenvalue equation
for a small eigenvalue \(\lambda\)
eliminates the uniformly invertible
vertical block. The resulting
two-dimensional pencil is
\[
        \pm f_*t I-\lambda R_{\cO}^2I
                            +O(t^2+t|\lambda|+\lambda^2).
 \tag{7.18}
\]
Continuity and the spectral gap
at zero leave exactly two small
eigenvalues. A Rayleigh bound on
the two-dimensional kernel and
the same block elimination give
\(|\lambda|=O(t)\); substituting
in (7.18) yields (7.15)--(7.16).
Equivalently this follows by
restricting the first derivative
of the symmetric Hessian family
to its isolated kernel.
The remaining eigenvalues stay
above half their lower bound
at zero by continuity.

The exact interval computation
described below gives
\[
 264.534132055\le R_{\cO}^2\le264.534132059
 \tag{7.19}
\]
and the whole-box force enclosure
of f12 V92. Positive interval
division gives (7.17).
\end{proof}

These are eigenvalues of the
classical hidden Hessian, not
diffusion eigenvalues or particle
masses. Their vanishing as
\(t\downarrow0\) shows why a
uniform \emph{strict Hessian}
bound cannot hold through this
symmetry-breaking limit. It does
not show failure of a uniform
Poincare inequality: V1 already
controls the unperturbed orbit
by Haar mixing.

\subsection{Certified curvature after the hidden normal response}

At the north pole use the 944
commuting angular variables
\(\xi\) of f12 V89, and let
\(\Phi(\xi,t)\) be their actual
completed action at \(C(t)\).
The strict cutoff margins put
these configurations on the
same smooth local completion
branches. Set, at \((\xi_*,0)\),
\[
 H_*=\Phi_{\xi\xi},\qquad
 v_*=\Phi_{\xi t},\qquad
 a_*=\Phi_{tt},\qquad
 Q_*=a_*-v_*^{\mathsf T}H_*^{-1}v_* .
 \tag{7.20}
\]
The normal minimum at this pole
is commuting, by Theorem~\ref{thm:f13v2-split}.
All 944 longitudinal directions
belong to the normal space;
the residual circle has no other
fixed directions. Thus solving
the scalar stationary equation
there supplies the full normal
response, not a frozen-slice
approximation to it.

\begin{theorem}[Actual polar energy curvature]
\label{thm:f13v2-curvature}
The smooth polar critical branches
have expansions
\[
 \cS_t(Y_\pm(t))
       =E_*\pm f_*t+\tfrac12Q_*t^2+O(t^3),
 \tag{7.21}
\]
with the certified exact enclosure
\[
 \frac{1794311357}{10^9}\le Q_*
                              \le\frac{1794311359}{10^9}.
 \tag{7.22}
\]
In particular this second derivative
is positive, but substantially
smaller than the unrelaxed value
\(a_*\simeq4.384889707\).
\end{theorem}
\begin{proof}
At the north pole,
\(\Phi_\xi(\xi(t),t)=0\).
The invertible positive \(H_*\)
gives
\[
                       \xi'(0)=-H_*^{-1}v_*.
\]
Twice differentiating
\(\Phi(\xi(t),t)\) yields (7.20)
as its second derivative. The
first derivative is \(f_*\).
A fixed simultaneous conjugation
sending \(T\) to \(-T\) transforms
the south branch at \(t\) into
the north branch at \(-t\).
Their second derivatives agree
and their first derivatives
are opposite, proving (7.21).
The exact arithmetic certificate
for (7.22) follows next.
\end{proof}

For each scalar plaquette row let
\(r_q\in\mathbb Q^{944}\) be its
angular derivative, whose entries
are integers divided by 15, and
let \(\omega_q\) be its exact
phase at \(\xi_*\).
Let \(\sigma_q=\pm1\) for the
six plaquettes containing
\(\ell_0\), and zero otherwise.
Direct differentiation of the
actual scalar action gives
\[
 \begin{split}
 H_*&=2\sum_q\cos\omega_q\,r_qr_q^{\mathsf T},\\
 v_*&=2\sum_q\sigma_q\cos\omega_q\,r_q,\qquad
 a_*=2\sum_{\sigma_q\ne0}\cos\omega_q.
 \end{split}
 \tag{7.23}
\]
The \(r_q\) include the derivatives
of all four moving inverse-root
pivots. There are 2376 nonzero
scalar rows, 11080 nonzero
Hessian entries, and eight
nonzero mixed-vector intervals.
The verifier independently
enumerates the plaquette words
and checks its rows against the
inherited exact scalar system.

It encloses every phase over the
coordinate box
\(\widehat\xi_i+[-13/10^{13},13/10^{13}]\),
using the inherited rational
trigonometric enclosures and
outward fixed-point arithmetic
at denominator \(2^{70}\).
The full box lies inside the
f12 V89 critical-point ball
because \(944(13/10^{13})^2<10^{-16}\).
The exact critical point is
therefore covered. The positive
matrix certificate is rerun on
this entire Hessian enclosure,
giving
\[
                        H_*\ge\frac{2166501}{10^8}I
                                      \ge\frac1{100}I.
 \tag{7.24}
\]
All accelerated matrix operations
in this positivity test use
signed integer arrays with
checked absolute dot-product
bounds below \(2^{63}\).

The new frozen witness is a
rational vector \(h\), with 944
integer numerators and denominator
\(10^9\). It was proposed
numerically, but the accepted
verifier imports only the frozen
integers, not a floating-point
solver. For arbitrary such \(h\),
put \(r=v_*-H_*h\). The identity
\[
 v_*^{\mathsf T}H_*^{-1}v_*
    =2v_*^{\mathsf T}h-h^{\mathsf T}H_*h
                                      +r^{\mathsf T}H_*^{-1}r
 \tag{7.25}
\]
is exact, by expansion.
Equation (7.24) bounds its last
term between zero and \(100\|r\|^2\).
Consequently any interval bounds
for the first two terms and for
the residual norm give rigorous
two-sided bounds, regardless
of how the witness was proposed.
The calculation certifies
\[
 \|r\|^2\le
       \frac{287551}{20000000000000000000},
 \qquad
 2.590578349\le v_*^{\mathsf T}H_*^{-1}v_*
                                      \le2.590578351,
 \tag{7.26}
\]
and, with the unrounded rational
interval endpoints retained until
the final outward reporting step, gives
(7.22). Subtracting only the
rounded decimal displays would
give a slightly wider enclosure;
the verifier uses the original
rational endpoints.

The files are
\begin{center}\small\ttfamily
scripts/ym\_f13\_v2\_curvature\_witness.py\\
scripts/verify\_ym\_f13\_v2\_relaxed\_coarse\_curvature.py.
\end{center}
The witness file's SHA--256 is
\begin{center}\scriptsize\ttfamily
C22EF1D97458A5A4628EF6918B844E490B23CE9407A335DA0662D331B4444323.
\end{center}
The comma-separated integer-vector
fingerprint, a different object,
is
\begin{center}\scriptsize\ttfamily
DFEFC088ADF0D84C1420ED4D20A7AC5242B1BF0D5C9048CE9739A9C6E09587E8.
\end{center}
The same verifier encloses
the sum (7.3) and verifies
the exact rational endpoints
in (7.17), (7.19), (7.22),
and (7.26). Interval acceptance,
not approximate eigenvalues,
is the numerical trust boundary.
A sign-reversed copy of the
frozen witness was also tested:
its residual enclosure does not
pass the declared tight curvature
interval. This checks rejection
of that corrupted witness; it
does not replace the proof of
the residual identity.

\subsection{A cusp for the fully relaxed local minimum}

\begin{corollary}[Local minimized energy and failure of a continuous selector]
\label{cor:f13v2-cusp}
For the fixed closed tube and
the changed coarse link
\(\exp X\), define the actual
local minimum
\[
        m_{\cT}(X)=\min_{Y\in\overline{\cT_\delta}}
                             \cS(Y;v(e^X)).
 \tag{7.27}
\]
With \(r=|x|\) the angular norm,
\[
        m_{\cT}(X)=E_*-f_*r+\tfrac12Q_*r^2+O(r^3).
 \tag{7.28}
\]
Thus \(m_{\cT}\) is not
differentiable at zero. For
small \(X\ne0\) the local
minimizer is unique, but it
has no continuous extension
as a selected hidden minimizer
through \(X=0\).
\end{corollary}
\begin{proof}
The closed-tube minimum is the
south branch by
Theorem~\ref{thm:f13v2-split}.
Conjugation covariance makes
its value depend only on
\(r=|x|\). Equation (7.21)
then gives (7.28). The
outward directional derivative
at zero is \(-f_*<0\) in
every unit direction, which
cannot be a linear functional.

Along a ray \(x=r\widehat x\),
the normal displacement tends
to zero, so the unique
minimizer tends to
\(p(-\widehat x)\).
Distinct axes give distinct
tuples because at least one
\(\sin\theta_i\ne0\).
Two such rays give different
limits, precluding a continuous
extension of this unique
punctured-neighbourhood selector.
\end{proof}

Unlike the f12 V92 envelope,
(7.27) minimizes over all hidden
directions in the entire tube,
not just the original orbit.
It still does not minimize
over all of \(M\).
No equality of this local
minimum with the global
conditional ground energy
is asserted. The cusp is
therefore a proved local
feature; it is not a proof
of a cusp of the global
ground energy.

\subsection{Implications, limits, and the next acquisition}

There is no small-\(t\) proliferation
of additional critical wells
\emph{within this tube} for
the chosen coarse-edge family.
The zero modes resolve into
one minimum and one index-two
saddle, not two metastable
minima. The relaxed local
ground energy has a certified
quadratic response and an
unavoidable radial cusp.
An approach demanding a
single smooth hidden-minimum
graph through this coarse
value must therefore be
replaced by an orbit-aware
description.

This is not yet a quantitative
escape estimate from the
complement of the tube.
The tube radius \(\delta\),
the perturbation threshold
\(t_0\), and the constants
in the Taylor remainders
are existential fixed-parent
constants here, not numerical
certificates. Nor has V1's
large-\(\beta\) Poincare
inequality been made uniform
over this perturbation family:
its invariant product-law
argument used the central
coarse input. The new
single-well reduced geometry
is an input to such an
extension, not that extension
itself.

The next spectral question
inside this family is control
of the tilted angular law
and the actual normal
conditional measure uniformly
in both \(\beta\) and small
\(t\), including the regime
\(\beta t\) large. A
comparison costing
\(\exp(C\beta t)\) would
lose the new one-well
information. The distinct
global question remains
coverage or escape control
for competitors outside
the tube. Neither question
may assume that the located
orbit was globally minimizing.

The predecessor V1 had 29,840
bytes and SHA--256
\begin{center}\scriptsize\ttfamily
64D4DE9AA5EA12695BD7DECC38A8A0C7698CAC3B69675ED0FCC06CB62FD7D2DE.
\end{center}
Its complete body is retained
above apart from the title
version. Only active V1 is
replaced by V2. The next
companion audit remains V5
and the direction review V10.
All archives and unrelated
parallel-programme files
are unchanged.

\begin{firewall}
The new result is the actual
local critical-set classification
under a small coarse perturbation,
together with the validated
relaxed curvature and lifted
classical Hessian slopes.
No global hidden ground
coverage, quantitative
escape theorem, all-exterior
repair estimate, cofinal
limit, quantum construction
or physical Yang--Mills
mass gap is proved here.
The full goal remains active.
\end{firewall}
\section{V3: a uniform local Poincare inequality through the coarse tilt}
\label{sec:f13v3}

\subsection{The new estimate and its boundary}

V2 classified the actual critical points
inside a normal tube after a small
single-edge coarse perturbation. It
left open whether the local diffusion
inequality survives uniformly when
\(\beta t\) is large. This version
proves that it does, with the actual
inverse-pivot amplitude included.
The proof does not compare the whole
perturbed density to its value at
\(t=0\), which would introduce an
exponential oscillation factor.

The new input is an all-tilt angular
inequality, followed by conditional
variance decomposition in the
nontrivial normal bundle. The
temperature-dependent normal
concentration compensates the
derivative of the conditional law.
All constants remain fixed-parent
constants; there is still no
quantitative escape estimate for
the complement of the tube.

\subsection{A weighted one-dimensional inequality with its proof}

\begin{lemma}[One-sided Hardy estimate]
\label{lem:f13v3-hardy}
Let \(\rho,w>0\) on \((a,b)\),
with the following integrals
finite at interior points, and put
\[
 R(x)=\int_a^x\frac{ds}{w(s)},\qquad
 M(x)=\int_x^b\rho(s)\,ds,\qquad
 B=\sup_{a<x<b}M(x)R(x)<\infty.
 \tag{8.1}
\]
Then, for a smooth \(g\) with
\(g(a)=0\),
\[
                  \int_a^b g^2\rho\,dx
                         \le4B\int_a^b(g')^2w\,dx.
 \tag{8.2}
\]
The reflected statement holds
with the anchor at the right
endpoint. Limits extend the
inequality to the corresponding
weighted form domain with that
anchor trace.
\end{lemma}
\begin{proof}
Cauchy--Schwarz, with the
factor \(R^{1/2}\) in the
energy weight, gives
\[
 |g(x)|^2
 \le\left(\int_a^x(g')^2wR^{1/2}\,ds\right)
          \left(\int_a^x\frac{ds}{wR^{1/2}}\right)
 =2R(x)^{1/2}\int_a^x(g')^2wR^{1/2}\,ds.
 \tag{8.3}
\]
Integration by parts against the
tail \(M\), first on a compact
truncated interval, gives
\[
 \begin{split}
 \int_s^bR(x)^{1/2}\rho(x)\,dx
 &\le R(s)^{1/2}M(s)
             +\frac12\int_s^bM(x)R(x)^{-1/2}R'(x)\,dx\\
 &\le\frac{B}{R(s)^{1/2}}
             +\frac B2\int_s^bR(x)^{-3/2}R'(x)\,dx\\
 &\le\frac{2B}{R(s)^{1/2}}.
 \end{split}
 \tag{8.4}
\]
The omitted upper boundary term
has nonpositive sign. Tonelli's
theorem applied to (8.3), followed
by (8.4), proves (8.2).
Truncation and monotone limits
justify possibly singular
endpoints; reflection gives
the left-hand version.
\end{proof}

This is the standard sufficient
weighted Hardy estimate, proved
here to fix its constant and
endpoint convention.\footnote{
B. Muckenhoupt, \emph{Hardy's
inequality with weights},
Studia Mathematica 44 (1972),
31--38, DOI
\href{https://doi.org/10.4064/sm-44-1-31-38}
{10.4064/sm-44-1-31-38}.
The specific sufficient estimate
used in this note is (8.2),
with the proof given above.}

\subsection{A sphere inequality uniform in the strength of a linear tilt}

Write \(u=n_3\) on the unit sphere
and \(d\omega\) for normalized
area. For \(h\ge0\), let
\[
                    d\sigma_h(n)
                     =Z_h^{-1}e^{-hu}\,d\omega(n).
 \tag{8.5}
\]

\begin{proposition}[All-tilt angular inequality]
\label{prop:f13v3-linear}
For every \(h\ge0\),
\[
             \Var_{\sigma_h}f
                 \le4\int_{S^2}|\nabla_{S^2}f|^2\,d\sigma_h.
 \tag{8.6}
\]
No restriction on the size of
\(h\) is imposed.
\end{proposition}
\begin{proof}
For \(0\le h\le1\), normalized
area has Poincare constant \(1/2\).
One proof applies the integrated
Bochner identity on the unit
two-sphere, using
\(|\Hess f|^2\ge(\Delta f)^2/2\)
and \(\Ric=g\), to a first
nonconstant eigenfunction.
The density ratio has oscillation
at most \(e^{2h}\le e^2\).
The variational formula for
variance then gives constant
at most \(e^2/2<4\).

For \(h\ge1\), set
\[
 x=u+1\in(0,2),\qquad
 p(x)=x(2-x),\qquad a=1/h\le1.
 \tag{8.7}
\]
The radial density is
proportional to \(e^{-hx}\),
with radial energy weight
\(p(x)e^{-hx}\).
For an anchor at \(a\), the
two Hardy products, with
normalization cancelling, are
\[
 \begin{split}
 B_-&=\sup_{0<x<a}
       \left(\int_0^xe^{-hs}\,ds\right)
       \left(\int_x^a\frac{e^{hs}}{p(s)}\,ds\right),\\
 B_+&=\sup_{a<x<2}
       \left(\int_x^2e^{-hs}\,ds\right)
       \left(\int_a^x\frac{e^{hs}}{p(s)}\,ds\right).
 \end{split}
 \tag{8.8}
\]
For \(x<a\), \(p(s)\ge s\)
and \(e^{hs}\le e\) on
\([x,a]\). Therefore
\[
 B_-\le\sup_{0<x<a}ex\log(a/x)=a=1/h.
 \tag{8.9}
\]
For the other side use
\[
                \frac1{s(2-s)}
                    =\frac12\left(\frac1s+\frac1{2-s}\right).
 \tag{8.10}
\]
Put \(M_h(x)=\int_x^2e^{-hs}ds\).
Since \(1/s\le h\) for \(s\ge a\),
\[
 M_h(x)\int_a^x\frac{e^{hs}}s\,ds
       \le\frac{e^{-hx}}h e^{hx}=\frac1h .
 \tag{8.11}
\]
For the second term, set
\(\delta=2-x>0\) and \(z=h\delta\).
An exact change of variable gives
\[
 \begin{split}
 M_h(x)\int_a^x\frac{e^{hs}}{2-s}\,ds
 &=\frac{e^{h\delta}-1}{h}
                      \int_\delta^{2-a}\frac{e^{-hv}}v\,dv\\
 &\le\frac{e^z-1}{h}
                      \int_z^\infty\frac{e^{-v}}v\,dv\\
 &\le\frac{1-e^{-z}}{hz}\le\frac1h .
 \end{split}
 \tag{8.12}
\]
Here \(\int_z^\infty e^{-v}dv/v
\le e^{-z}/z\).
Equations (8.10)--(8.12) show
\(B_+\le1/h\).
The two applications of
Lemma~\ref{lem:f13v3-hardy},
with \(g(x)=m(x)-m(a)\),
give
\[
       \Var_{\rm rad}m
            \le\frac4h\int_0^2p(x)|m'(x)|^2\,d\sigma_{\rm rad}.
 \tag{8.13}
\]
The anchor need not be a median:
variance is no larger than
the squared distance to
any constant.

Away from the measure-zero
poles, the law is a product
of this radial density and
uniform azimuth \(\varphi\).
Conditional circle variance
has constant one. Moreover
\[
 |\nabla_{S^2}f|^2
           =p(x)|\partial_xf|^2
                    +p(x)^{-1}|\partial_\varphi f|^2.
 \tag{8.14}
\]
Since \(p(x)\le1\), the
azimuthal variance is bounded
by its term in this energy.
Apply (8.13) to the azimuthal
average and use Jensen's
inequality on its derivative.
Variance decomposition gives
constant \(\max(1,4/h)\le4\).
Smooth approximation extends
the result to \(H^1(S^2,\sigma_h)\).
\end{proof}

\subsection{Nonlinear angular slopes without an exponential comparison loss}

\begin{lemma}[Bounded density change]
\label{lem:f13v3-density}
If \(\nu=r\mu/\mu(r)\), where
\(0<r_-\le r\le r_+<\infty\),
and \(\mu\) has Poincare
constant \(C\) for a given
metric, then \(\nu\) has
constant at most \(Cr_+/r_-\).
\end{lemma}
\begin{proof}
In \(\Var_\nu f=\inf_c\nu((f-c)^2)\)
choose \(c=\mu(f)\), bound its
density above by \(r_+/\mu(r)\),
apply the \(\mu\) inequality,
and bound \(d\mu\) by
\(\mu(r)d\nu/r_-\) in the
energy integral.
\end{proof}

\begin{proposition}[A uniform slope-ratio inequality]
\label{prop:f13v3-slope}
Suppose an axisymmetric potential
\(V(u)\) is smooth on the sphere
and satisfies
\[
           0<\alpha\le V'(u)\le c\alpha
                         \quad(-1<u<1),\qquad c\ge1.
 \tag{8.15}
\]
Then the probability law
proportional to \(e^{-V(u)}d\omega\)
has unit-sphere Poincare
constant at most \(4c^5\).
\end{proposition}
\begin{proof}
Let \(h=(V(1)-V(-1))/2\) and
introduce an increasing
latitude coordinate \(v\) by
\[
             v=-1+\frac{V(u)-V(-1)}h .
 \tag{8.16}
\]
One has \(\alpha\le h\le c\alpha\).
For the inverse \(u=U(v)\),
\[
            1/c\le U'(v)=h/V'(U(v))\le c,
       \qquad \frac{\sup U'}{\inf U'}\le c .
 \tag{8.17}
\]
The transformed measure is
proportional to
\(U'(v)e^{-hv}dv\,d\varphi\).
By Proposition~\ref{prop:f13v3-linear}
and Lemma~\ref{lem:f13v3-density},
its Poincare constant in the
\(v\)-sphere metric is at
most \(4c\).

The latitude map
\((v,\varphi)\mapsto(U(v),\varphi)\)
has Lipschitz constant at most
\(c^2\) for unit-sphere metrics.
Indeed \(U(\pm1)=\pm1\) and
(8.17) give
\[
        c^{-2}\le\frac{1-U(v)^2}{1-v^2}\le c^2.
 \tag{8.18}
\]
The azimuthal derivative has
norm at most \(c\); the
meridional derivative has
norm
\[
                  U'(v)\sqrt{\frac{1-v^2}{1-U(v)^2}}\le c^2.
\]
These length bounds extend
the map over both poles as
a Lipschitz map. Pulling back
a test function multiplies
its squared gradient by at
most \(c^4\). Thus its
Poincare constant is at
most \(4c\cdot c^4=4c^5\).
\end{proof}

The comparison above changes
the latitude coordinate; it
does not compare \(V\) to
a linear function by the
oscillation of their difference.
Consequently its constant
does not depend on \(\alpha\).

\subsection{The actual normal marginal, uniformly in temperature and tilt}

Retain the V2 tube, decreasing
its radius \(\delta>0\) and
parameter bound \(t_0>0\) when
needed. These are fixed
independently of \(\beta\).
Let \(d=2830\), and use the
unit sphere (7.2) as base of
the normal bundle. Write
\[
 \begin{split}
 s_t(q,z)&=\cS_t(\exp_qz),\\
 b_t(q,z)&=a_t(\exp_qz)J(q,z),\\
 e_t(q)&=\min_{|z|\le\delta}s_t(q,z),
       \qquad |z_t(q)|\le\delta/2.
 \end{split}
 \tag{8.19}
\]
Here \(a_t\) is the actual
full inverse-pivot Haar
amplitude, independent of
\(\beta\), and \(J\) is the
normal-exponential volume
Jacobian used in V1.
The constant converting
induced orbit volume to
unit-sphere area is absorbed
in normalization. All these
are smooth positive densities
where appropriate; no
quotient-orbit-volume division
is introduced.

On this compact closed disk
bundle and parameter interval
there are finite constants
\[
 \begin{gathered}
 c_N I\le D_z^2s_t\le C_N'I,
       \qquad c_N=1/10000,\qquad
       0<b_-\le b_t\le b_+,\\
 \|D_z^2\log b_t\|_{\rm op}\le B_N .
 \end{gathered}
 \tag{8.20}
\]
The lower action bound is
V2's normal-convexity result.
The others follow by
smoothness and compactness.
Set
\[
       \beta_*=\max\{1,20000 B_N\},\qquad C_N=20000.
 \tag{8.21}
\]
For \(\beta\ge\beta_*\), each
normal conditional measure
\[
 d\nu_{q,\beta,t}(z)
   =\frac{b_t(q,z)e^{-\beta s_t(q,z)}\,dz}
          {\int_{|w|<\delta}b_t(q,w)e^{-\beta s_t(q,w)}\,dw}
 \tag{8.22}
\]
has Euclidean-ball Poincare
constant at most \(C_N/\beta\).
Indeed its potential Hessian
is at least
\(\beta/10000-B_N\ge\beta/20000\);
apply Lemma~\ref{lem:ball}.

\begin{proposition}[A populated angular marginal constant]
\label{prop:f13v3-marginal}
There is a finite \(D\ge1\),
independent of \(\beta\ge1\)
and \(0\le t\le t_0\), such
that the actual base marginal
\(\overline\mu_{\beta,t}\) of
the tube measure has Poincare
constant
\[
                           C_B=972D
 \tag{8.23}
\]
for the unit-sphere metric.
\end{proposition}
\begin{proof}
Write the base marginal as
\[
 d\overline\mu_{\beta,t}(q)
  \ \propto\ e^{-\beta e_t(q)}
          A_{\beta,t}(q)\,d\omega(q),\qquad
 A_{\beta,t}(q)=
 \int_{|z|<\delta}b_t(q,z)
                  e^{-\beta(s_t(q,z)-e_t(q))}\,dz.
 \tag{8.24}
\]
Taylor's formula about the
normal minimum, within the
convex normal ball, gives
\[
 \frac{c_N}{2}|z-z_t|^2
       \le s_t(q,z)-e_t(q)
       \le\frac{C_N'}2|z-z_t|^2.
 \tag{8.25}
\]
Put \(r_0=\min(\delta/2,1)\).
The ball of radius
\(r_0/\sqrt\beta\) about
\(z_t(q)\) lies in the
normal domain. Lower integration
over this ball and upper
integration over all of
\(\mathbb R^d\) give
\[
       A_-\beta^{-d/2}
          \le A_{\beta,t}(q)
          \le A_+\beta^{-d/2},
 \tag{8.26}
\]
where one may take
\[
 A_-=b_-e^{-C_N'r_0^2/2}
                  \operatorname{vol}(B_1^d)r_0^d,
 \qquad
 A_+=b_+(2\pi/c_N)^{d/2},
 \qquad D=A_+/A_-.
 \tag{8.27}
\]
The ratio is independent
of both parameters. This
is a two-sided integral
bound, not a differentiated
Laplace asymptotic expansion.

For \(t>0\), V2's smooth
remainder (7.11) and
Lemma~\ref{lem:f13v2-latitude},
with a smaller \(t_0\), give
\[
              f_*t/2\le\partial_u e_t
                                  \le3f_*t/2.
 \tag{8.28}
\]
Apply Proposition~\ref{prop:f13v3-slope}
to \(V=\beta e_t\), with
\(\alpha=\beta f_*t/2\)
and \(c=3\). Its constant
is \(4\cdot3^5=972\),
regardless of \(\beta t\).
Equation (8.26) and
Lemma~\ref{lem:f13v3-density}
then give (8.23).
At \(t=0\), \(e_0=E_*\)
and rotational invariance
even makes the base marginal
uniform; the same bound
remains valid.
\end{proof}

The ratio \(D\) need not be
small and is not numerically
certified here. It is finite
at the fixed parent and does
not grow with \(\beta\) or
with \(\beta t\). Nothing in
(8.27) asserts useful growth
in dimension or regulator size.

\subsection{Differentiating conditional means in the nontrivial normal bundle}

Use the homogeneous connection
on \(G\times_K N_p\cO\):
horizontal group generators
are Hilbert--Schmidt orthogonal
to the stabilizing circle
algebra. Its transport maps
normal fibres isometrically,
so it preserves the fixed
ball, its Euclidean volume,
and its boundary. Denote
horizontal differentiation
along a unit base direction
by \(\mathcal H_a\).

For an orthonormal frame of
the unit-sphere tangent plane,
smoothness and invariance
at \(t=0\) give a fixed
finite \(K_0\) such that
\[
 \begin{split}
 \sum_{a=1}^2|D_z\mathcal H_as_t|^2
                                  &\le K_0^2t^2,\\
 \sum_{a=1}^2|D_z\mathcal H_a\log b_t|^2
                                  &\le K_0^2t^2 .
 \end{split}
 \tag{8.29}
\]
These are frame-independent
squared tensor norms.
To justify their order in
\(t\), at zero both \(s_0\)
and \(b_0\) are unchanged
under simultaneous conjugation,
including its action on
the normal vector. Thus
their horizontal derivatives
vanish identically on the
whole disk bundle at \(t=0\).
Differentiate this identity
vertically, then apply a
uniform bound on the
\(t\)-derivative of the
resulting smooth tensor.
This argument retains the
actual amplitude, not only
its value on the orbit.

Let \(W=\beta s_t-\log b_t\).
Minkowski's inequality in
these tensor norms implies
\[
       \sum_{a=1}^2|D_z\mathcal H_aW|^2
                         \le K_0^2t^2(\beta+1)^2.
 \tag{8.30}
\]
For a smooth function \(F\)
on the tube, write
\(\widehat F(q,z)=F(\exp_qz)\)
and \(m(q)=\nu_{q,\beta,t}(\widehat F)\).
Differentiating the normalized
integral (8.22) gives exactly
\[
 \mathcal H_am
   =\nu_{q,\beta,t}(\mathcal H_a\widehat F)
       -\operatorname{Cov}_{\nu_{q,\beta,t}}
                         (\widehat F,\mathcal H_aW).
 \tag{8.31}
\]
There is no moving-boundary
term: differentiation uses
the fixed-radius fibres
under isometric transport,
not balls recentered at
their varying normal minima.
The normalization derivative
is precisely the covariance
subtraction.

\begin{lemma}[Uniform conditional-mean derivative bound]
\label{lem:f13v3-mean}
For \(\beta\ge\beta_*\),
\[
 |\nabla_{S^2}m|^2
 \le 2\nu_{q,\beta,t}
                  \left(\sum_a|\mathcal H_a\widehat F|^2\right)
       +8C_N^2K_0^2t^2
                    \nu_{q,\beta,t}(|D_z\widehat F|^2).
 \tag{8.32}
\]
\end{lemma}
\begin{proof}
Use Cauchy--Schwarz for each
covariance in (8.31), then
the normal Poincare inequality
both for \(\widehat F\) and
for each score \(\mathcal H_aW\).
Their squared covariance sum
is at most
\[
 \begin{split}
 &\Var_\nu(\widehat F)\sum_a
                           \Var_\nu(\mathcal H_aW)\\
 &\quad\le
 \frac{C_N^2}{\beta^2}
       \nu(|D_z\widehat F|^2)
             \nu\!\left(\sum_a|D_z\mathcal H_aW|^2\right)\\
 &\quad\le
       4C_N^2K_0^2t^2\nu(|D_z\widehat F|^2),
 \end{split}
 \tag{8.33}
\]
where (8.30) and \(\beta\ge1\)
were used. The square of the
mean derivative term is
bounded by Jensen's inequality.
Finally \(|A-B|^2\le2|A|^2+2|B|^2\)
proves (8.32).
\end{proof}

The cancellation in (8.33)
is essential: two normal
Poincare factors \(1/\beta\)
compensate the squared
order-\(\beta t\) score
derivative. No assumption
\(\beta t\ll1\) is made.

\subsection{The uniform actual tube theorem}

Let \(\mu_{\beta,t}^{\cT}\)
be the actual normalized
restriction to this fixed
\(\cT_\delta\), with density
\(a_t e^{-\beta\cS_t}\)
relative to product Haar.

\begin{theorem}[Uniform local diffusion control through symmetry breaking]
\label{thm:f13v3-uniform}
For all \(\beta\ge\beta_*\),
\(0\le t\le t_0\), and
\(F\in H^1(\cT_\delta,\mu_{\beta,t}^{\cT})\),
\[
 \Var_{\mu_{\beta,t}^{\cT}}F
 \le
 \left[
       \frac{80000}{\beta}
        +C_B\bigl(3776+32C_N^2K_0^2t^2\bigr)
 \right]
       \int_{\cT_\delta}|\nabla_MF|^2\,d\mu_{\beta,t}^{\cT},
 \tag{8.34}
\]
where \(C_N=20000\) and
\(C_B=972D\).
In particular, the local
Neumann diffusion has a
positive lower gap bound
uniform in these parameters:
\[
 g_{\beta,t}^{\cT}\ge C_{\rm loc}^{-1},
 \qquad
 C_{\rm loc}
  =80000+972D\bigl(3776+32(20000)^2K_0^2t_0^2\bigr)<\infty.
 \tag{8.35}
\]
The same constants apply
for every changed coarse
link \(\exp X\) at \(e_0\)
with angular norm \(|x|\le t_0\).
\end{theorem}
\begin{proof}
First compare pullback gradients
with the physical hidden metric.
The normal-exponential derivative
bound from V1 gives
\[
                   |D_z\widehat F|^2
                          \le4|\nabla_MF|^2\circ\exp.
 \tag{8.36}
\]
A unit tangent vector on the
unit sphere of axes is generated
by a horizontal Lie element
\(Z\) with
\(\|Z\|_{\rm HS}^2=1/2\):
adjoint motion of quaternion
axes has angular speed twice
the coefficient norm of \(Z\).
For each independent group,
\(\|[Z,U]\|_{\rm HS}\le2\|Z\|_{\rm HS}\).
Summing over 944 groups,
the squared operator norm
of the horizontal lift followed
by the normal exponential
is at most \(2\cdot944=1888\).
The lift is simply simultaneous
conjugation of \(\exp_qz\).
Taking its adjoint yields
\[
              \sum_a|\mathcal H_a\widehat F|^2
                          \le1888|\nabla_MF|^2\circ\exp.
 \tag{8.37}
\]
This is an operator-norm bound,
not a bound multiplied again
by the two base dimensions.

The normal-bundle disintegration
has base marginal
\(\overline\mu_{\beta,t}\) and
conditional laws (8.22).
Its variance identity is
\[
 \Var_{\mu_{\beta,t}^{\cT}}F
       =\overline\mu_{\beta,t}
                      \bigl(\Var_{\nu_{q,\beta,t}}\widehat F\bigr)
                         +\Var_{\overline\mu_{\beta,t}}m.
 \tag{8.38}
\]
The first term is at most
\((4C_N/\beta)\int|\nabla_MF|^2d\mu^{\cT}\)
by (8.22) and (8.36).
For the second term use
Proposition~\ref{prop:f13v3-marginal},
then Lemma~\ref{lem:f13v3-mean},
and (8.36)--(8.37).
Its coefficient is at most
\(C_B(3776+32C_N^2K_0^2t^2)\).
Since \(4C_N=80000\), this
proves (8.34).

Smooth functions are dense
in the weighted \(H^1\)
space of this smooth tube
at each fixed pair of
parameters, so the form
inequality extends as stated.
It implies (8.35) by the
variational definition of
the Neumann gap and the
parameter bounds.
Finally simultaneous
conjugation transforms any
axis of \(X\) into \(T\).
It preserves the tube and
metric and transports the
actual completed density,
so the same constants
work in every coarse
direction at this edge.
\end{proof}

\subsection{What is now closed locally and what is not}

The local inequality is no
longer restricted to the
central coarse input.
It is uniform across
the small three-dimensional
single-edge coarse family,
through \(t=0\), and for
arbitrarily large \(\beta t\).
It is compatible with
V2's lifted classical
Hessian eigenvalues tending
to zero: angular mixing,
not uniform strict Hessian
positivity, supplies the
missing directions.

Every density in the
argument is the actual
normal disintegration.
The amplitude appears in
the Gaussian comparison,
the normal potential and
the score derivative.
The proof has not changed
the conditional measure
or replaced it by its
leading Laplace term.

The constants
\(\delta,t_0,\beta_*,D,K_0\)
are proved finite but
are not numerically enclosed.
In particular, (8.35)
is an existence theorem
for a uniform fixed-parent
local constant, not a
numerical global gap.
The factor \(D\) may be
very large at normal
dimension 2830.
No regulator-uniform
or all-exterior bound
is inferred from it.

For this family the V73
escape assembly can now
use the same local input
\(\gamma\ge C_{\rm loc}^{-1}\)
through the perturbation.
It still needs an
independent quantitative
killed floor outside a
smaller tube. Nothing
above estimates that
floor or proves that
all conditional global
minima lie in the tube.
The next upstream work
must address omitted
competitors or their
escape geometry, rather
than repeat the now
closed local tilt estimate.

\subsection{Verification and continuation record}

The inequalities in this
version are proved for
all stated parameters,
not inferred from spectral
sampling. As an implementation
diagnostic, the two exact
Hardy-product formulas
were evaluated with stable
floating quadrature at
60 points spanning
\(h=1,2,10,100,1000\),
including points near
both poles. They satisfied
\(hB_\pm\le1\). This
finite test only checks
the formulas for mistakes;
(8.9)--(8.12) are their
all-parameter proof.
No additional numerical
Hessian certificate is
needed for this deduction.

V2 had 51,291 bytes and
SHA--256
\begin{center}\scriptsize\ttfamily
12A3ADD7C134B37B93D284C44E0D08C7810B53E9551402A00FC128DF751236D2.
\end{center}
Its full mathematical body
is retained above, with
only the title updated.
Active V2 is replaced
by V3. Archives and
parallel-programme files
remain unchanged.
The companion review
remains scheduled for
V5 and the direction
review for V10.

\begin{firewall}
The new theorem is a
uniform local conditional
diffusion inequality for
one specified finite
parent and coarse-edge
family. It does not prove
global escape, complete
ground coverage, unrestricted
repair, a cofinal limit,
quantum reconstruction
or the physical
Yang--Mills mass gap.
The full objective
remains active.
\end{firewall}
\section{V4: complete central-sector ground classification and collective escape}
\label{sec:f13v4}

\subsection{Why move upstream, and precisely which sector is closed}

V3 supplies a uniform local Poincare inequality on the known
noncentral orbit tube, including its specified small coarse-edge
perturbations. Another reformulation of that local inequality would
not supply the missing global escape estimate. We therefore return
to the actual fixed-parent action outside that tube. Throughout
this section the exterior and coarse field are the fixed central
exterior and coarse field \(I\) of f12 V86--V92, and
\[
 M=\SU(2)^{944},\qquad S:M\longrightarrow\mathbb R
\]
is the actual completed action. The 96 pivots are functions of
the 944 free links; they are not frozen when differentiating.

The new results are: the exact minimum \(648\) over the
\emph{entire} central sector \(\{\pm I\}^{944}\); its complete
256-point equality set; and a uniform negative-curvature
certificate at every point of that set. They yield a local
killed-diffusion estimate on a neighbourhood of this complete
central ground set. Here \emph{central ground} always means
minimum within the central sector, not minimum on \(M\).
The known noncentral orbit already has action below \(570\).

This does not repeat f12 V83, which classified a restricted
56-link sheet with all other free links prescribed. Nor does
it repeat f12 V91, whose universal central lower bound was
\(616\). The directional localization identity used at the
end is the one already proved in f12 V77; the new input is
its verification on all central ground configurations.

\subsection{An exact parity strengthening of the inherited dual bound}

Use the deterministic free-link ordering in the retained
parent-geometry code, and write \(Y_i=(-I)^{x_i}\),
\(x_i\in\mathbb F_2\), \(0\le i<944\). The central-pivot
identification in f12 V91 gives \(B_e=I\) for every such
assignment. Consequently the exact action has the form
\begin{equation}\label{eq:f13v4-central-action}
 S_{\rm cen}(x)=408+4\sum_{j=0}^{2291}m_j y_j(x),\qquad
 y_j(x)=c_j+\sum_{i\in I_j}x_i\pmod 2 .
\end{equation}
Here \(y_j\) is evaluated as \(0\) or \(1\);
the positive integers \(m_j\) account for repeated parity
terms. These data are regenerated from all 7776 physical
plaquettes, not inferred from sampled configurations.

Set \(t=(x,y)\in\{0,1\}^{3236}\) and let \(c\) be the
objective vector, so \(S_{\rm cen}=408+c\cdot t\).
The inherited parity inequalities are \(At\le b\).
For every forbidden bit assignment \(\xi\) on one parity
relation, its row is
\[
 \sum_i(2\xi_i-1)t_i\le\sum_i\xi_i-1 .
\]
This excludes exactly that forbidden assignment on its
participating bits and is valid on every correct parity
assignment. The frozen integral dual multiplier \(u\le0\)
has 637 nonzero entries and obeys
\[
 u\cdot b=4784,\qquad r=23c-A^{\mathsf T}u\ge0 .
\]
Since \(23\cdot408+4784=23\cdot616\), every central
assignment satisfies the exact identity
\begin{equation}\label{eq:f13v4-dual-identity}
 23(S_{\rm cen}-616)
 =\sum_\ell(-u_\ell)(b_\ell-(At)_\ell)+r\cdot t .
\end{equation}

\begin{lemma}[Finite odd-parity packing certificate]
\label{lem:f13v4-packing}
There are eight subsets \(C_1,\ldots,C_8\) of the 3236
coordinates, of respective sizes
\[
 6,\ 6,\ 14,\ 11,\ 16,\ 13,\ 31,\ 17,
\]
such that, identically in \(x\),
\[
 \bigoplus_{i\in C_j}t_i(x)=1,\qquad
 r'=r-92\sum_{j=1}^8{\bf1}_{C_j}\ge0.
\]
The complete subsets are the literal integer tuples in
\begin{center}\small
\path{scripts/ym_f13_v4_central_certificate_data.py}.
\end{center}
\end{lemma}

\begin{proof}
Represent each affine parity \(t_i(x)\) by its 944-bit
coefficient mask and constant bit. For each listed subset,
XOR of the coefficient masks is zero and XOR of the
constants is one. Subtract 92 from the corresponding
residual coordinates for each occurrence and check every
remaining coordinate is nonnegative. These are the explicit
integer tests in
\path{scripts/verify_ym_f13_v4_central_ground_set.py};
they regenerate the dual residual before testing the subsets.
No linear-programming output is trusted without these checks.
Thus the finite certificate proves both assertions.
\end{proof}

\begin{proposition}[Sharp full central lower bound]
\label{prop:f13v4-sharp}
Every one of the \(2^{944}\) central assignments has action
at least \(648\), and this bound is attained.
\end{proposition}

\begin{proof}
The odd parity implies \(\sum_{i\in C_j}t_i\ge1\).
All terms in the first sum of
\eqref{eq:f13v4-dual-identity} are nonnegative, and
\[
 r\cdot t
 =r'\cdot t+92\sum_{j=1}^8\sum_{i\in C_j}t_i
 \ge 8\cdot92=736=23\cdot32.
\]
Hence \(S_{\rm cen}\ge648\). The sheet assignment from
f12 V82 attains \(648\); it is also checked independently
by the equality classification below.
\end{proof}

\subsection{Exhaustive equality classification with only fifteen free bits}

At equality, the preceding nonnegative decomposition forces
\(t_i=0\) whenever \(r'_i>0\). These give 1908 affine
binary equations in \(x\): 207 from free-bit coordinates
and 1701 from plaquette-parity coordinates. They are
\emph{necessary} equality conditions; sufficiency is not
assumed without checking the actual action.

Gaussian elimination over \(\mathbb F_2\), taking the
largest set bit as pivot, gives rank \(929\). The unpivoted
coordinates in increasing order are
\begin{equation}\label{eq:f13v4-free}
 (2,5,16,24,38,52,59,98,106,112,213,220,269,325,332).
\end{equation}
Call their fifteen-bit integer encoding \(q\), with the
first coordinate as its least significant bit. Solving
the echelon system gives an injective affine map
\[
 L:\mathbb F_2^{15}\longrightarrow\mathbb F_2^{944},
 \qquad L_i(q)=\operatorname{parity}(a_i\mathbin{\&}q)
                    \mathbin{\oplus} b_i .
\]
The code constructs the masks \(a_i\) and constants \(b_i\)
and substitutes them back into all 1908 equations. This
checks the parametrization itself, not only the rank.

\begin{theorem}[Complete central ground set]
\label{thm:f13v4-classification}
With \(L\) as above, the central minimizers are exactly
\[
 \mathcal Z=\{((-I)^{L_i(q)})_{i=0}^{943}:q\in\mathcal Q\},
\]
where addition in the following formula is bitwise XOR:
\begin{equation}\label{eq:f13v4-codes}
 \mathcal Q
 =226\mathbin{\oplus}
 \operatorname{span}_{\mathbb F_2}
       \{1,6,416,576,1024,2048,8192,16384\}.
\end{equation}
Thus \(|\mathcal Z|=256\), and every member has action \(648\).
Every other central assignment has action at least \(656\).
\end{theorem}

\begin{proof}
Every minimizer belongs to the necessary equality face,
so is \(L(q)\) for one of the \(2^{15}=32768\) codes.
For each code, substitute its bits into
\eqref{eq:f13v4-central-action}. The verifier enumerates
all codes, evaluates all 2292 parity terms by integer
XOR and adds their integer weights. The largest possible
sum is \(408+4\sum_jm_j=9912<2^{63}\), certifying
absence of overflow in its signed 64-bit energy arrays.
All mask operations have only fifteen bits.

The minimum is 648. The set of codes with this value
is compared, as an exact set, with the affine set in
\eqref{eq:f13v4-codes}. The eight generators produce
256 distinct elements, so they are independent. This
proves completeness and cardinality, rather than merely
constructing 256 examples. The inherited sheet assignment
is reconstructed in all 944 coordinates and has code
\(772\in\mathcal Q\).

Finally, every physical fine link belongs to six
plaquettes on the periodic side-six lattice. Multiplying
all central plaquette signs therefore gives \(+1\).
The number of negative plaquettes is even, and
\(S_{\rm cen}\), four times this number, lies in
\(8\mathbb Z\). A central action strictly greater than
648 is consequently at least 656.
\end{proof}

The eight-unit statement is a discrete central-sector
action spacing, not a Hamiltonian or diffusion gap.

\subsection{The full central Hessian must include moving pivots}

Fix \(z\in\mathcal Z\), a Hilbert--Schmidt unit generator
\(X\in\mathfrak{su}(2)\), and real velocities \(w_i\).
Use the curve \(Y_i(t)=z_i\exp(tw_iX)\).
Every physical plaquette is initially \(\pm I\), so
its first action derivative is zero. Smooth pivot
completion therefore makes \(z\) a critical point
of the full action \(S\).

For a pivot \(e\), let \(N_e^+\) count its positive
nonroot source paths among the fourteen chronological
nonroot paths. The other two paths are root duplicates.
At the central point, positive paths are in the inner
logarithm region and negative paths are in an open
region where the cutoff vanishes. Write \(b'_eX\)
for the pivot velocity and \(p'_jX\) for a positive
source-path velocity. Linearization of the actual
completion equation gives
\begin{equation}\label{eq:f13v4-pivot-response}
 (16-N_e^+)b'_e+\sum_{j:P_j=I}p'_j=0,\qquad
 p'_j=\sum_{\text{free endpoints of }j}w_i.
\end{equation}
Fixed endpoints have velocity zero. In general the
fourteen paths do \emph{not} occur in sign-identical
pairs: varying internal central links can make \(N_e^+\)
odd. A denominator justified only for even \(N_e^+\)
would be invalid. We use
\[
 D=\operatorname{lcm}(2,3,\ldots,16)=720720 .
\]
Let \(p_Q\in\{\pm1\}\) be a physical plaquette sign.
Combine free velocities and
\eqref{eq:f13v4-pivot-response}, with their plaquette
orientations, into its scalar first velocity
\[
 a_Q=D^{-1}\sum_iR_{Qi}w_i,\qquad R_{Qi}\in\mathbb Z .
\]
The rows are generated from the physical geometry by
\path{scripts/ym_f13_v4_central_geometry.py}.

\begin{lemma}[Exact Hilbert--Schmidt quadratic form]
\label{lem:f13v4-hessian}
For every real velocity vector \(w\),
\begin{equation}\label{eq:f13v4-hessian}
 \Hess_z S(wX,wX)
   =D^{-2}\sum_Q p_Q\Big(\sum_iR_{Qi}w_i\Big)^2.
\end{equation}
The product-metric squared norm of \(wX\) is
\(\sum_iw_i^2\). More generally the Hessian is this
scalar matrix tensored with the identity on the
three Hilbert--Schmidt Lie-algebra axes.
\end{lemma}

\begin{proof}
The denominator in
\eqref{eq:f13v4-pivot-response} is
\(16-N_e^+\in\{2,\ldots,16\}\), proving integrality
after multiplication by \(D\). For a Hilbert--Schmidt
unit \(X\), the plaquette action along its first
velocity is \(2(1-p_Q\cos(ta_Q/\sqrt2))\), with
second derivative \(p_Qa_Q^2\) at zero.
All first plaquette derivatives vanish, so acceleration
of the completed path contributes nothing to the
second derivative of \(S\). Summation proves the formula.
For general Lie-algebra velocities, the same calculation
gives the squared Hilbert--Schmidt norm of each linear
plaquette velocity. There are no cross terms between
orthonormal Lie axes, proving the tensor statement.
\end{proof}

\subsection{A collective saddle certificate for all 256 configurations}

\begin{theorem}[Uniform negative direction on the central ground set]
\label{thm:f13v4-saddles}
At every \(z\in\mathcal Z\) there is a product
Hilbert--Schmidt unit tangent vector \(v_z\) such that
\[
 \Hess_zS(v_z,v_z)\le-\frac1{20}.
\]
The Hessian has a negative subspace of dimension at
least three. Nondegeneracy and exact Morse index
are not asserted.
\end{theorem}

\begin{proof}
The file
\path{scripts/ym_f13_v4_escape_witness_data.py}
contains 24 literal sparse integer vectors \(w^{(k)}\)
and disjoint code groups assigning one vector to
every \(q\in\mathcal Q\). Their largest support is
362; the sum of support sizes is 4639.
The exact verifier
\path{scripts/verify_ym_f13_v4_central_escape.py}
reruns Theorem~\ref{thm:f13v4-classification}, checks
that the groups partition precisely \(\mathcal Q\),
and checks all vector indices and coefficients.

For each code it reconstructs every physical
plaquette sign and the moving-pivot rows. It independently
checks physical action \(648\), and then checks, with
arbitrary-precision integers,
\begin{equation}\label{eq:f13v4-integer-saddle}
 20\sum_Qp_Q\Big(\sum_iR_{Qi}w_i^{(k)}\Big)^2
       \le -D^2\sum_i(w_i^{(k)})^2 .
\end{equation}
Each vector is nonzero. Divide by its squared norm
and apply Lemma~\ref{lem:f13v4-hessian}. Three
orthonormal Lie generators with the same free-link
vector span a negative subspace.

As an additional normalization check, the same geometry
code applied to the old 56-link sheet vector at code
772 reproduces its Hilbert--Schmidt curvature
\(-6/55\). The least negative quotient among the
256 assigned witnesses is exactly
\[
 -\frac{270616421}{4916151354}<-\frac1{20}.
\]
No eigenvalue approximation or rounding tolerance
enters the final acceptance test.
\end{proof}

Floating-point eigenvectors were used only to propose
integer witnesses. The frozen vectors, exact geometry,
and inequality \eqref{eq:f13v4-integer-saddle} are the
certificate. The inherited sheet vector alone does
not cover this set. Unsuccessful coordinate or
equal-amplitude two-coordinate tests would not
establish positivity: the directions here are collective.

\subsection{Local escape near the complete central ground set}

Let the actual finite-parent conditional law at this
fixed coarse field be
\[
 d\mu_\beta=Z_\beta^{-1}e^{-\beta S}a\,d{\rm vol}_M ,
\]
where the actual positive smooth amplitude \(a\) is
independent of \(\beta\), as in the previous sections.
Write \(\E_\beta(h)=\int_M|\nabla h|^2d\mu_\beta\).

\begin{corollary}[Escape from central-ground neighbourhoods]
\label{cor:f13v4-escape}
There are pairwise disjoint open coordinate
neighbourhoods \(U_z^+\) of the 256 points, a finite
\(\beta_s\), and \(U^+=\bigcup_{z\in\mathcal Z}U_z^+\),
all independent of \(\beta\), such that
\[
 \E_\beta(h)\ge\frac{\beta}{320}
       \|h\|_{L^2(\mu_\beta)}^2,\qquad
 h\in H_0^1(U^+),\quad\beta\ge\beta_s .
\]
Functions are extended by zero to \(M\).
Neighbourhood sizes and \(\beta_s\) are proved to
exist, not numerically enclosed here.
\end{corollary}

\begin{proof}
Choose product exponential coordinates in
Hilbert--Schmidt orthonormal axes about each \(z\).
At the origin let the constant Euclidean unit
direction \(v_z\) be given by
Theorem~\ref{thm:f13v4-saddles}. Smoothness and
shrinking the charts ensure
\[
 \partial_{v_z}^2S\le-\frac1{40},
 \qquad g^{-1}\ge\frac12 I .
\]
There are finitely many distinct centres. Thus
the charts can be chosen disjoint, with compact
closures inside larger charts where these bounds hold.

Let \(J\) denote the chart volume density.
The unnormalized Euclidean density is
\(e^{-\Phi_\beta}\), with
\[
 \Phi_\beta=\beta S-\log a-\log J .
\]
Smooth positivity and compactness give a finite
common bound \(B\) for
\(\left|\partial_{v_z}^2(\log a+\log J)\right|\).
For \(\beta\ge\max(1,80B)=:\beta_s\),
\[
 \partial_{v_z}^2\Phi_\beta
 \le-\beta/40+B\le-\beta/80 .
\]
Apply the directional ground-state identity of
f12 V77. Its relevant algebra, with
\(g=h e^{-\Phi_\beta/2}\) and \(h\) compactly
supported in one chart, is
\begin{align*}
 \int|\partial_{v_z}h|^2e^{-\Phi_\beta}\,dx
 &=\int\left\{|\partial_{v_z}g|^2+
 \left(\frac{|\partial_{v_z}\Phi_\beta|^2}{4}
 -\frac{\partial_{v_z}^2\Phi_\beta}{2}\right)g^2
 \right\}dx\\
 &\ge\frac{\beta}{160}\int h^2e^{-\Phi_\beta}\,dx .
\end{align*}
The metric bound contributes \(1/2\). The normalizer
cancels, giving \(\beta/320\). Density extends the
estimate to \(H_0^1\) of each chart. The disjoint
union is a direct sum of these Dirichlet spaces,
so summing gives the same constant, without a factor
of 256.
\end{proof}

For comparison with action exclusion only, the f12
V91 Lipschitz constant \(20592\), in its maximum
free-link displacement metric, upgrades the common
central \(\varepsilon=1/4096\) tube action bound to
\[
 648-\frac{20592}{4096}
       =\frac{164601}{256}=642.97265625 .
\]
This is more than 73 above the known noncentral
orbit. It is an action bound, not a substitute for
the Dirichlet estimate. The escape charts need not
equal those fixed-radius action tubes.

\subsection{Certificate ledger, preservation, and remaining acquisition}

The accepted checkers are
\path{scripts/verify_ym_f13_v4_central_ground_set.py}
and
\path{scripts/verify_ym_f13_v4_central_escape.py}.
Both report PASS. The frozen parity certificate has SHA--256
\begin{center}\scriptsize\ttfamily
4661F307B74335FFF4045584D29987CF8B76424B4B07E5A55A0B4AE22C4E21A3;
\end{center}
the frozen escape witness file has SHA--256
\begin{center}\scriptsize\ttfamily
29BC421231F5431D10645ACEFA6B1B7170767299AE277AB2B0B408EF664B660C.
\end{center}
The sorted comma-separated minimizer codes have fingerprint
\begin{center}\scriptsize\ttfamily
688CD7E3E456025E4FD4E73B5B88F85B157C21A57034050F3483C5B294CF2862.
\end{center}
Fingerprints identify the finite data; they do not
substitute for the arithmetic checks. The acceptance
path uses no optimizer or exploratory module.

The predecessor V3 contained 72,631 bytes and SHA--256
\begin{center}\scriptsize\ttfamily
030B3A3AC6CDFA1164421DA7835D98BB9AF8CEDCAA7977A0167DB6C0F433ABD9.
\end{center}
Its mathematical body is retained above, with only
the title updated. Active V4 replaces V3; archives
and companion notes are unchanged. The next companion
review is V5, and the next direction review is V10.

\begin{firewall}
The central sector is completely classified at its
minimum, and all 256 ground configurations have proved
local escape directions. This does not cover higher-action
central critical points, omitted noncentral critical points,
transition regions, or the full complement of the known
noncentral tube. It does not yet supply the global killed
floor required by V1 and f12 V73. The estimate
\(\beta/320\) concerns a finite-parent conditional
diffusion on specified Dirichlet neighbourhoods; it is
not a physical mass. The next upstream acquisition is
coverage or escape for the remaining complement, followed
by the still-open uniform exterior, cofinal continuum,
and quantum reconstruction interfaces. The Yang--Mills
mass-gap objective remains unproved and active.
\end{firewall}
\section{V5: a non-enumerative central escape criterion and companion review}
\label{sec:f13v5}

\subsection{The required companion checkpoint and the upstream choice}

The preceding iteration was progress: V4 classified the full
central ground set and certified escape at each of its 256
points. It did not settle the other \(2^{944}-256\) central
critical configurations. Rather than classify successive
energy shells, this iteration seeks a structural criterion
covering whole families without enumerating their members.

The required five-version checkpoint inspected the current
files, their hashes, and their concluding mathematical
sections. They are unchanged from the preceding checkpoint:
\begin{center}\small
\begin{tabular}{lll}
\toprule
programme and active version&bytes&SHA--256 prefix\\
\midrule
SM--Einstein V72&607372&\texttt{F44F9745D43379E1}\\
NS V26&278283&\texttt{82C887F8C2C69E4F}\\
Parallel YM V5&54174&\texttt{306F1BB84CFA8A8D}\\
\bottomrule
\end{tabular}
\end{center}
SM V72 proves its specified covariance limit with
finitely many modes, bounded time, and many copies.
It does not establish ultraviolet uniformity for the
physical finite-species continuum. NS V26 supplies the
stated rank-one norm formulae, not global regularity.
The parallel YM note retains compact sector labels and
rejects a volume-uniform extensive global-charge tail
as a general assembly target. None supplies a missing
escape estimate for the present parent. No companion
file is modified, and no document-internal stop instruction
is treated as a new user instruction.

The fixed parent, central exterior, coarse field \(I\),
metric, and actual pivot completion remain those of V4.
The new argument combines a blockwise trace estimate
for pivot responses with a new exact parity dual. It
does not assume the pivots remain stationary on an
arbitrary free-link perturbation.

\subsection{Bare curvature, pivot-stationary directions, and the trace defect}

At any central configuration \(z\in\{\pm I\}^{944}\),
write \(p_Q\) for its physical plaquette signs.
For one Hilbert--Schmidt Lie axis, let \(A\) be the
oriented plaquette incidence matrix of the 944 free
links, with all physical pivot velocities set to zero.
Thus
\[
 H=A^{\mathsf T}\operatorname{diag}(p_Q)A,\qquad
 H_0=A^{\mathsf T}A .
\]
Here \(H\) is a \emph{bare} scalar Hessian, not the
full completed Hessian \(H_{\rm comp}\).
Let \(B_z:\mathbb R^{944}\to\mathbb R^{96}\) be the
actual scalar pivot derivative from
\eqref{eq:f13v4-pivot-response}, and set
\[
 K_z=\ker B_z,\quad r_z=\operatorname{rank}B_z\le96,
 \quad P_z=\text{orthogonal projection onto }K_z^\perp .
\]
For \(w\in K_z\) the pivot velocities vanish, so
\begin{equation}\label{eq:f13v5-restriction}
 w^{\mathsf T}H_{\rm comp}w=w^{\mathsf T}Hw .
\end{equation}
All physical plaquette first derivatives vanish at
the central point; second-order pivot motion therefore
does not change this equality.

For a moving plaquette put \(k_Q=\#\{\hbox{free links in }Q\}\).
Define the two nonnegative integers
\[
 N(z)=\#\{Q:k_Q>0,\ p_Q=-1\},\qquad
 J(z)=\sum_{Q:p_Q=-1}k_Q .
\]
There are 102 fixed negative plaquettes, and the exact
incidence counts are
\begin{equation}\label{eq:f13v5-counts}
 \#\{Q:k_Q=1,2,3,4\}=(464,880,688,344).
\end{equation}
Every free physical link occurs in six plaquettes.
Consequently
\begin{equation}\label{eq:f13v5-trace}
 S_{\rm cen}=408+4N,\qquad
 \operatorname{tr}H=5664-2J,\qquad
 0\le J\le5664 .
\end{equation}

\begin{lemma}[Uniform bare operator bound]
\label{lem:f13v5-norm}
As quadratic forms, \(-H_0\le H\le H_0\) and
\(0\le H_0\le16I\).
\end{lemma}

\begin{proof}
The first inequalities follow term by term from
\(p_Q\in\{\pm1\}\). Extend a free-link scalar vector
by zero to all physical links on the periodic lattice.
The full plaquette incidence is the discrete exterior
derivative \(d\) on one-forms. At a Fourier mode let
\(h_\mu=e^{ik_\mu}-1\). Its two-form components are
\(h_\mu v_\nu-h_\nu v_\mu\), up to harmless common
phase conventions. Direct expansion gives
\[
 \sum_{\mu<\nu}|h_\mu v_\nu-h_\nu v_\mu|^2
 =|h|^2|v|^2-\left|\sum_\mu\overline{h_\mu}v_\mu\right|^2
 \le16|v|^2 .
\]
Parseval and restriction to zero-extended free vectors
prove \(H_0\le16I\).
\end{proof}

\subsection{Why the cost of removing the pivot directions is nine, not sixteen}

For a free link \((x,\mu)\), define its block label
\[
 b(x,\mu)=
 \left(\mu,\big(\lfloor x_\nu/2\rfloor\big)_{\nu\ne\mu}\right).
\]
There are 32 labels in this parent. All free links
appearing in one pivot's source paths have the same
label, and there are three pivot rows per label.
These are support statements, independent of the
central signs and of which source paths are active.
Thus the row space of \(B_z\), and hence \(P_z\), are
orthogonal direct sums over these coordinate blocks.

\begin{lemma}[Blockwise projection trace]
\label{lem:f13v5-block}
For every central assignment, with \(r=r_z\),
\[
 \operatorname{tr}(P_zH_0)\le9r,\qquad
 \operatorname{tr}(H|_{K_z})\le5664-2J+9r .
\]
The second trace is the trace of the compressed
quadratic form on \(K_z\).
\end{lemma}

\begin{proof}
Within one block, links have the same orientation.
Distinct such links occur together in a plaquette
only when they are parallel transverse neighbours,
and their incidence product is \(-1\). A block fixes
all three transverse coarse coordinates. In each
transverse direction, a fine coordinate has at most
one neighbour inside that same two-site coarse cell.
There are therefore at most three within-block
neighbours per link. The block diagonal of \(H_0\)
is \(6I-C\), where \(C\) is a symmetric adjacency
matrix of maximum degree three. For every real \(v\),
\[
 |v^{\mathsf T}Cv|
 \le\sum_{\{i,j\}\ {\rm edge}}(v_i^2+v_j^2)
 \le3|v|^2 .
\]
Thus each diagonal block is at most \(9I\).
Only these blocks contribute to
\(\operatorname{tr}(P_zH_0)\), since \(P_z\) is block
diagonal. Summing their projection ranks proves the
first inequality.

The positive semidefiniteness of \(H+H_0\) implies
\(\operatorname{tr}(P_zH)\ge-\operatorname{tr}(P_zH_0)\).
Using \eqref{eq:f13v5-trace},
\[
 \operatorname{tr}(H|_{K_z})
 =\operatorname{tr}H-\operatorname{tr}(P_zH)
 \le5664-2J+9r .
\]
This argument bounds the actual row-space projection;
it does not set all pivot derivatives to zero by fiat.
\end{proof}

\begin{theorem}[Incidence criterion for collective central escape]
\label{thm:f13v5-trace}
Every central configuration with \(J\ge3265\) has a
Hilbert--Schmidt unit direction of completed curvature
at most \(-1/424\). More precisely,
\begin{equation}\label{eq:f13v5-rayleigh}
 \lambda_{\min}(H_{\rm comp})
 \le\frac{6528-2J}{848}.
\end{equation}
Its full three-axis Hessian has a negative subspace
of dimension at least
\begin{equation}\label{eq:f13v5-index}
 3\max\left\{0,\
 \left\lceil\frac{2J-6528}{16}\right\rceil\right\}.
\end{equation}
\end{theorem}

\begin{proof}
The space \(K_z\) has dimension \(944-r\ge848\).
Its smallest Rayleigh quotient is at most its
average eigenvalue. Equations
\eqref{eq:f13v5-restriction} and
Lemma~\ref{lem:f13v5-block} give
\[
 \lambda_{\min}(H_{\rm comp})
 \le\frac{5664-2J+9r}{944-r}.
\]
The derivative of the right side as a function of
\(r\) has numerator \(14160-2J>0\), since \(J\le5664\).
Replacing \(r\) by 96 proves
\eqref{eq:f13v5-rayleigh}. At \(J=3265\) it equals
\(-2/848=-1/424\).

By Lemma~\ref{lem:f13v5-norm}, each eigenvalue of
\(H|_{K_z}\) is at least \(-16\). If there are \(k\)
negative eigenvalues, its trace is at least \(-16k\).
The trace upper bound is at most \(6528-2J\);
hence \(k\ge(2J-6528)/16\). Its negative subspace
is also negative for \(H_{\rm comp}\) by
\eqref{eq:f13v5-restriction}. The central Hessian
tensors with the three Lie axes, as in V4.
Taking the nonnegative integer lower bound proves
\eqref{eq:f13v5-index}.
\end{proof}

\subsection{An exact parity dual improves the action threshold}

The incidence counts alone imply
\[
 J\ge3N-1808.
\]
Indeed \(J-3N\) has coefficient \(-2\) on the 464
one-free-link plaquettes, \(-1\) on the 880
two-free-link plaquettes, zero on the next class,
and \(+1\) on the last. This elementary estimate
would imply escape for \(S_{\rm cen}\ge7176\).
The parity constraints improve it.

\begin{lemma}[New global parity trace bound]
\label{lem:f13v5-dual}
Every central assignment satisfies
\begin{equation}\label{eq:f13v5-dual}
 J-3N\ge-\frac{14167}{9}.
\end{equation}
\end{lemma}

\begin{proof}
Use the same binary variables \(t=(x,y)\), parity
rows \(At\le b\), and term multiplicities as V4.
In this proof only, let \(c\) denote the coefficient
vector for \(J-3N\):
\[
 c_i=0\ (i<944),\qquad
 c_{944+j}=(|I_j|-3)m_j .
\]
The file
\path{scripts/ym_f13_v5_trace_dual_data.py}
contains an integral multiplier \(u\le0\) with
1125 nonzero rows, on denominator 36.
Form \(r=36c-A^{\mathsf T}u\). Unlike the earlier
minimum-action dual, \(r\) need not be nonnegative.
The box bounds \(0\le t_i\le1\) give the exact
weak-duality inequality
\begin{align*}
 36\,c\cdot t
 &=u\cdot At+r\cdot t\\
 &\ge u\cdot b+\sum_i\min(0,r_i)
 =-56668 .
\end{align*}
The last equality, sign and row-index checks, and
every residual coordinate are recomputed with
arbitrary-precision integers by
\path{scripts/verify_ym_f13_v5_high_action_escape.py}.
Division by 36 proves \eqref{eq:f13v5-dual}.
The floating-point linear programme proposed \(u\);
neither numerical optimality nor floating-point
feasibility is used in this proof.
\end{proof}

\begin{corollary}[All sufficiently high-action central configurations escape]
\label{cor:f13v5-high}
Every central configuration with \(S_{\rm cen}\ge6864\)
has a Hilbert--Schmidt unit direction of completed
curvature at most \(-1/106\).
\end{corollary}

\begin{proof}
Then \(N=(S_{\rm cen}-408)/4\ge1614\), so
\[
 J\ge3\cdot1614-\frac{14167}{9}
   =\frac{29411}{9}.
\]
Since \(J\) is an integer, \(J\ge3268\). Equation
\eqref{eq:f13v5-rayleigh} yields
\(-8/848=-1/106\). No enumeration of the central
configurations in this action range is needed.
\end{proof}

The verifier checks that 6864 is the first central
action level allowed by \(8\mathbb Z\) at which this
particular bound forces negative trace. It is not
claimed to be an optimal escape threshold.

\subsection{Nonvacuity and the induced Dirichlet neighbourhoods}

The accepted checker also constructs an actual
central configuration in the covered range. It
starts from
\[
 x_i=(37i+\lfloor i/7\rfloor)\pmod2
\]
and performs 176 deterministic strictly action-increasing
bit flips, each choosing the largest action increase
and breaking ties by the smallest index. Independent
re-evaluation of all physical plaquettes gives
\[
 S_{\rm cen}=7336,\qquad J=4284.
\]
Thus the covered family is nonempty. At this example
the theorem proves curvature at most \(-255/106\)
and a full negative subspace of dimension at least
\(3\lceil2040/16\rceil=384\). The ascent is only a
finite way to specify an example; no global
maximization claim is inferred.

Let
\[
 \mathcal Z_{\rm tr}=\{z\in\{\pm I\}^{944}:J(z)\ge3265\},
 \qquad
 \mathcal Z_{\rm hi}=\{z:S_{\rm cen}(z)\ge6864\}.
\]
They are finite sets, and
\(\mathcal Z_{\rm hi}\subseteq\mathcal Z_{\rm tr}\).
Repeating the already proved directional localization
argument of V4 gives pairwise disjoint sufficiently
small coordinate neighbourhoods with the following
Dirichlet bounds, for all sufficiently large \(\beta\):
\[
 \E_\beta(h)\ge\frac{\beta}{6784}\|h\|_2^2
 \quad\hbox{on the union about }\mathcal Z_{\rm tr},
\]
and the stronger bound \(\beta/1696\) on the union
about \(\mathcal Z_{\rm hi}\).
For clarity, a curvature margin \(-\kappa\) at each
centre becomes \(\partial_v^2S\le-\kappa/2\) on
small charts. Bounded logarithmic amplitude and
volume derivatives then give
\(\partial_v^2\Phi_\beta\le-\kappa\beta/4\).
The directional identity and inverse-metric lower
bound \(1/2\) yield \(\kappa\beta/16\).
Finiteness makes the neighbourhood and temperature
choices uniform over each set. The resulting
constants exist but their chart radii and temperature
thresholds are not numerically enclosed here.

Adding the 256 central ground neighbourhoods from
V4 preserves the smaller \(\beta/6784\) bound after
shrinking all charts to be disjoint. No partition
error or cardinality factor is needed for this
disjoint Dirichlet union.

\subsection{Verification record and what is still not covered}

The accepted verifier is
\path{scripts/verify_ym_f13_v5_high_action_escape.py}.
It checks the physical incidence counts, all 32
source blocks and their three-row support property,
the within-block Hessian signs and degree bound,
the rational parity dual, both threshold calculations,
and the independently evaluated example. It uses
no eigensolver, optimizer, or exploratory module.
The frozen dual data have SHA--256
\begin{center}\scriptsize\ttfamily
0DE4A402EA2774DB21A939215147E2C6C84172D1F4E77FA309AC4E3D25F87530.
\end{center}
This finite arithmetic verification accompanies
the analytic projection, Fourier, and localization
proofs above; it does not verify a continuum theorem.

V4 had 90,433 bytes and SHA--256
\begin{center}\scriptsize\ttfamily
431232169D92AB4375CA825D585053B54D211563E17D087994589B8193D87DA8.
\end{center}
Its complete body is retained, except for the
version title. Only the active predecessor is
replaced; the f12 archive and companion files
remain unchanged. The next companion checkpoint
and direction review are both V10.

\begin{firewall}
Together with V4, the new criterion leaves as
possible uncovered central configurations only
those with
\[
 656\le S_{\rm cen}\le6856,\qquad J\le3264.
\]
This is a necessary description of the remaining
central cases, not a statement that every such
case is stable or uncovered by other methods.
Noncentral critical points and transition regions
are still uncontrolled. The new union of Dirichlet
neighbourhoods is not the full complement of the
known noncentral ground-orbit tube. Hence the
global killed floor, uniform exterior estimates,
cofinal construction, quantum reconstruction, and
physical Yang--Mills mass gap remain unproved.
The next useful acquisition is a geometric or
parity-sensitive negative-direction criterion in
the trace-inconclusive region, not more sampling
as a substitute for coverage.
\end{firewall}
\section{V6: two-link escape throughout an explicit low-action parity family}
\label{sec:f13v6}

\subsection{Moving beyond the global trace test}

V5 made concrete progress by excluding an entire
high-action central family. Its trace criterion
leaves the low and middle action range largely
unresolved. This iteration uses a local zero-diagonal
mechanism instead of attempting another global
trace constant. The completed pivot response remains
part of the Hessian throughout.

The main new result is an explicit affine family
of \(2^{937}\) central configurations on which the
same two-link direction has uniformly negative
completed curvature. A checked member has action
2040 and \(J=927\), so it is outside both the V4
central ground set and the V5 trace criterion.
It is stable under every single central bit flip.
The certificate thus supplies genuinely new escape
coverage, not a new proof of the high-action bound.

All statements concern the same fixed finite
parent, exterior, coarse field \(I\), and
Hilbert--Schmidt metric as V4--V5. The large family
cardinality does not represent continuum or
Gibbs-probability coverage.

\subsection{A uniform upper bound on every completed scalar diagonal}

Let \(H_{\rm comp}\) be the scalar completed
Hessian at any central configuration, in the
normalization of V4. For a free link \(j\), let
\(b_{ej}\) be its derivative contribution to
physical pivot \(e\). Equation
\eqref{eq:f13v4-pivot-response} gives
\[
 b_{ej}=-\frac{n_{ej}}{16-N_e^+},\qquad
 n_{ej}\ge0 ,
\]
where \(n_{ej}\) counts positive source-path
occurrences of \(j\). Every free link occurs
in exactly two chronological source paths in
this parent. Since \(16-N_e^+\ge2\),
\begin{equation}\label{eq:f13v6-column}
 b_{ej}\le0,\qquad \sum_e|b_{ej}|\le1 .
\end{equation}
All pivots in this column have the same link
orientation as \(j\).

\begin{lemma}[Completed diagonal ceiling]
\label{lem:f13v6-diagonal}
For every central configuration and free link,
\[
 (H_{\rm comp})_{jj}\le14.
\]
\end{lemma}

\begin{proof}
Let \(L_0=d^{\mathsf T}d\) denote the positive
plaquette incidence Hessian on \emph{all} physical
links. Its diagonal is six. Between distinct
parallel links an entry is either zero or \(-1\).
Distinct parallel pivots cannot be transverse
nearest neighbours: their transverse coordinates
are all even. Thus their mutual entries in
\(L_0\) vanish.

The completed physical velocity of the free
coordinate vector \(e_j\) is
\[
 v=e_j+\sum_e b_{ej}e_{\ell_e},
\]
where \(\ell_e\) is the physical pivot link.
The signed physical Hessian is bounded above
by \(L_0\), term by term in its plaquette
representation. Consequently
\begin{align*}
 (H_{\rm comp})_{jj}
 &\le v^{\mathsf T}L_0v\\
 &=6+6\sum_e b_{ej}^2
      +2\sum_e(L_0)_{j,\ell_e}b_{ej}\\
 &\le6+6\left(\sum_e|b_{ej}|\right)^2
             +2\sum_e|b_{ej}|\le14 .
\end{align*}
This is an upper bound for the actual completed
diagonal, not a frozen-pivot approximation.
\end{proof}

\subsection{A local mixed-curvature rule}

Call a free link \(i\) \emph{source-inactive at \(z\)}
if both chronological source paths containing it
are negative at the central configuration. Then
every \(b_{ei}=0\). Its completed diagonal is the
bare signed incidence sum
\[
 d_i=(H_{\rm comp})_{ii}=\sum_{Q\ni i}p_Q.
\]
In particular \(d_i=0\) means exactly three of
the six incident plaquettes are negative.

\begin{theorem}[Two-link escape with a moving neighbour pivot]
\label{thm:f13v6-local}
Suppose at a central configuration that:
\begin{enumerate}
\item \(i\) is source-inactive and \(d_i\le0\);
\item no physical pivot parallel to \(i\) shares
a plaquette with \(i\);
\item \(j\) is a free parallel transverse nearest
neighbour of \(i\).
\end{enumerate}
Let \(p\in\{\pm1\}\) be the sign of their unique
shared plaquette. Then the free-link vector
\[
 w=14e_i+p e_j
\]
has squared norm 197 and obeys
\[
 \frac{w^{\mathsf T}H_{\rm comp}w}{197}
 \le-\frac{14}{197}.
\]
No source-inactivity assumption is imposed on \(j\).
The full Hessian has a negative subspace of
dimension at least three.
\end{theorem}

\begin{proof}
The first condition gives the completed physical
velocity \(e_i\) for the free coordinate \(i\).
The bare mixed entry with its parallel neighbour
is \(-p\). Completion of \(j\) adds only pivots
parallel to \(j\), hence parallel to \(i\).
By the second condition none couples to \(i\)
through a plaquette. Thus, exactly,
\[
 (H_{\rm comp})_{ii}=d_i\le0,\qquad
 (H_{\rm comp})_{ij}=-p .
\]
Using Lemma~\ref{lem:f13v6-diagonal},
\[
 w^{\mathsf T}H_{\rm comp}w
 =196d_i-28+(H_{\rm comp})_{jj}
 \le-14 .
\]
The product Hilbert--Schmidt norm is 197 on
a unit Lie axis. The three orthonormal Lie
axes tensor with the same scalar form at a
central configuration, proving the negative
subspace assertion.
\end{proof}

For orientation, the geometric second condition
holds for an internal free link with even
longitudinal coordinate. It also holds for a
boundary link with at least two odd transverse
coordinates: a parallel pivot has all transverse
coordinates even and cannot be its transverse
nearest neighbour. The exact example below
uses the second case.

\begin{record}[Why a zero one-link action does not imply stability]
If \(i\) is source-inactive and \(d_i=0\), varying
only \(Y_i=z_i\exp(tX)\) leaves all pivots exactly
at \(I\) for sufficiently small \(t\). Indeed
the affected source paths remain near \(-I\),
in an open region where the cutoff vanishes;
all other source paths are unchanged central
paths. The unique completed solution is still
\(I\). The six affected plaquette contributions
are \(2(1-p_Q\cos(t/\sqrt2))\), whose sum is
constant because \(\sum_{Q\ni i}p_Q=0\).
Nevertheless the nonzero mixed entry in
Theorem~\ref{thm:f13v6-local} gives escape when
the neighbouring free link is also varied.
This observation uses the actual local completion,
not merely a vanishing second derivative.
\end{record}

\subsection{An explicit rank-seven family}

Use the free-link indices already fixed in V4.
The selected links are
\[
 i=161:\ ((0,3,1,5),3),\qquad
 j=114:\ ((0,2,1,5),3).
\]
They are parallel transverse neighbours on the
side-six periodic fine lattice. The anchor has
two odd transverse coordinates, so no parallel
pivot shares a plaquette with it.

Write \(Y_k=(-I)^{x_k}\), with
\(x_k\in\mathbb F_2\). Define \(\mathcal C\)
by the following seven affine equations:
\begin{equation}\label{eq:f13v6-cylinder}
\begin{aligned}
 x_{161}&=1,&x_{932}&=1,&x_{146}&=0,&x_{148}&=0,\\
 x_{103}+x_{114}&=1,&
 x_{147}+x_{395}&=0,&
 x_{149}+x_{172}&=0 .
\end{aligned}
\end{equation}
All additions here are modulo two. No other
free bit is prescribed.

\begin{theorem}[Uniform escape on the whole affine family]
\label{thm:f13v6-cylinder}
The set \(\mathcal C\) has exactly \(2^{937}\)
central configurations. At every member,
\[
 (H_{\rm comp})_{161,161}=0,\qquad
 (H_{\rm comp})_{161,114}=-1 .
\]
Consequently the same scalar direction
\[
 w=14e_{161}+e_{114}
\]
has Hilbert--Schmidt normalized curvature at
most \(-14/197\) at every member.
\end{theorem}

\begin{proof}
The four singleton constraints and three pair
constraints in \eqref{eq:f13v6-cylinder} have
disjoint sets of variables. Their rank is seven,
so their solution space has dimension \(944-7=937\).

To verify the physical meaning of the equations,
enumerate the six plaquettes incident to link
161 in the inherited physical ordering. Their
free-bit masks are respectively
\[
\begin{gathered}
 \{103,114,161\},\quad\{146,161\},\quad
 \{147,161,395\},\\
 \{148,161\},\quad\{149,161,172\},\quad
 \{161,932\}.
\end{gathered}
\]
Including the fixed exterior signs, prescribing
their signs as
\[
 (+1,-1,-1,-1,+1,+1)
\]
gives those six mask equations with respective
right sides \((0,1,1,1,1,0)\).
The two source paths containing the anchor
both have only \(x_{161}\) as a free bit and
positive fixed remainder. Requiring both to
be negative gives \(x_{161}=1\) twice.
Eliminating this repeated equation leaves
exactly \eqref{eq:f13v6-cylinder}.

These geometric masks and fixed signs are
reconstructed by
\path{scripts/verify_ym_f13_v6_local_escape.py}.
The source paths are therefore inactive for
every member, the anchor diagonal is zero,
and the shared plaquette with link 114 has
sign \(+1\). The geometric no-parallel-pivot
condition is independent of all bits.
Theorem~\ref{thm:f13v6-local} now applies
throughout the entire affine family. The
proof does not enumerate or sample its
\(2^{937}\) members.
\end{proof}

The ratio \(2^{937}/2^{944}=1/128\) is only
a fraction of the uniform central counting
measure. It is not a bound on the actual
conditional Gibbs measure, a fraction of
continuum states, or complete central coverage.

\subsection{Verified low-action members and a moving-pivot check}

The literal central bit mask in
\path{scripts/ym_f13_v6_local_pattern_data.py}
specifies a member of \(\mathcal C\) with
\[
 S_{\rm cen}=2040,\qquad J=927 .
\]
Every one of its 944 bare diagonal entries
\(d_k=\sum_{Q\ni k}p_Q\) is nonnegative, with
minimum zero. Flipping central bit \(k\)
changes the action by \(4d_k\), since it
reverses each of its incident plaquette
signs and all central pivots are \(I\).
This configuration is thus stable under
every one-bit central move. It is not a
global central minimizer, and the V5
trace criterion is inconclusive at it.

The verifier evaluates the full completed
quadratic form with denominator \(720720\)
from V4 and obtains exactly
\[
 \frac{w^{\mathsf T}H_{\rm comp}w}{197}
       =-\frac{24}{197}.
\]
All physical plaquettes and all pivot
responses are included in this evaluation.

A second check changes the bit mask by a
homogeneous solution of the seven equations.
It remains in \(\mathcal C\), but the two
source paths containing link 114 are now
positive. Thus its neighbour pivot really
moves under the direction \(w\). Independent
completed evaluation gives
\[
 S_{\rm cen}=2072,\qquad
 \frac{w^{\mathsf T}H_{\rm comp}w}{197}
       =-\frac{86}{591}
       \le-\frac{14}{197}.
\]
This check guards against inadvertently
freezing the neighbour's pivot. The two
examples verify nonvacuity and normalization;
the preceding analytic proof, not the
examples, supplies universal family coverage.

\subsection{Local diffusion consequence and the remaining region}

Apply the V4 directional localization proof
to \(\mathcal C\), using
\(\kappa=14/197\). There are sufficiently
small disjoint neighbourhoods of all its
centres, with union \(U_{\mathcal C}\), and
a finite common \(\beta_{\mathcal C}\), such
that the actual finite-parent conditional
diffusion satisfies
\begin{equation}\label{eq:f13v6-dirichlet}
 \E_\beta(h)\ge\frac{7\beta}{1576}
                  \|h\|_{L^2(\mu_\beta)}^2,
 \qquad h\in H_0^1(U_{\mathcal C}),\quad
 \beta\ge\beta_{\mathcal C}.
\end{equation}
Indeed the already proved argument gives
\(\kappa\beta/16=7\beta/1576\).
Finiteness, not a cardinality estimate for
the Gibbs measure, gives common sufficiently
small chart sizes and temperature threshold.
These sizes and the threshold are existential
here, not numerically enclosed.

The same reasoning applies to the finite
set of all central configurations satisfying
the general local rule of
Theorem~\ref{thm:f13v6-local}. Combining this
set with V4's ground configurations and
V5's trace-positive configurations, taking
each centre only once and shrinking the
charts to be disjoint, gives the common
lower bound \(\beta/6784\) on their Dirichlet
union for sufficiently large \(\beta\).
This is still a local union, not the full
complement of the noncentral orbit tube.

\subsection{Certificate and preservation record}

The accepted checker
\path{scripts/verify_ym_f13_v6_local_escape.py}
uses integer masks, exact incidence counts,
and rational quadratic forms. It imports
no exploratory module, optimizer, or numerical
eigensolver. It checks the source-count bound
behind Lemma~\ref{lem:f13v6-diagonal}, the
physical pivot adjacency statements, the
seven-dimensional constraint rank, and the
two completed-Hessian evaluations.
The frozen local-pattern data have SHA--256
\begin{center}\scriptsize\ttfamily
7A0995ED58900FA55932EC96CD77388BAE0D1DAE7446AE177200AB2DA00A6E89.
\end{center}

The predecessor V5 contained 104,702 bytes
and SHA--256
\begin{center}\scriptsize\ttfamily
3253E926555CA4E3966094F1F5334D22FA440E4CEC1DABE4BB5EE253A24DB16A.
\end{center}
Its complete body is retained above, with
only the version title updated. Active V6
replaces active V5. Archives and companion
notes are unchanged. The next companion
review and direction review remain V10.

\begin{firewall}
The new local criterion removes an explicit
large affine family, including verified
low-action, central-bitwise-stable configurations,
from the previously trace-inconclusive region.
It does not prove that every remaining central
configuration has such an anchor. The residual
central acquisition is now the V5 range
\(656\le S_{\rm cen}\le6856,\ J\le3264\)
with the new local criterion also absent.
Noncentral critical points, transition regions,
and quantitative global escape remain open.
No exterior-uniform, cofinal, reconstructed
quantum, or physical mass-gap conclusion
follows from this finite-parent advance.
The full Yang--Mills goal remains active.
\end{firewall}
\section{V7: cutoff-transition Hessians and a ground-coverage warning}
\label{sec:f13v7}

\subsection{Why interrupt the central escape enumeration}

V6 added rigorous local escape coverage in a previously
trace-inconclusive central family. The next upstream
question is more basic than another central pattern:
does the known noncentral orbit tube account for all
low-action wells of this actual conditional action?
This iteration tests that question before assuming
that every point outside the tube should escape.

A bounded numerical descent from eight of V4's
central ground configurations found two additional
nearby \emph{candidate} branches with slightly higher
action than the certified orbit. Their source gates
enter the cutoff transition. Their computed non-orbit
curvatures are positive. These observations do
\emph{not} certify a second stationary point or a
local minimum. They do make full ground/well
coverage the immediate acquisition priority.

The rigorous new results in this section are the
general one-active-source scalar and transverse
completion Hessian formulae, including their
acceleration terms. They extend the inner-region
inverse-fifteenth-root formula used in f12 V88--V90.
The candidate computations are separated below from
these proved identities.

\subsection{The one-active-source radial completion}

Work on an open region where, for each pivot under
consideration, exactly one of its fourteen nonroot
source paths can contribute to the cutoff sum;
the other thirteen paths remain outside its support.
For that source write \(P=U_aU_b\). A fixed endpoint
is permitted. At the commuting base point,
\[
 P=\exp(pE_3),\qquad B=\exp(b(p)E_3),
\]
where \(E_3^2=-I\) and
\(\|E_3\|_{\rm HS}=\sqrt2\).
Choose the local signed angle branch with
\(|p|<\pi\). Put
\[
 h(u)=u\,\chi\!\left(\frac{\sqrt2|u|}{\rho}\right),
 \qquad \rho=\frac1{64}.
\]
The actual scalar completion equation is
\begin{equation}\label{eq:f13v7-implicit}
 b+\frac1{16}h(p-b)=0 .
\end{equation}
Here \(\chi=1\) on \([0,1/2]\), \(\chi=0\)
on \([1,\infty)\), and in the transition interval
the specified smooth cutoff is
\[
 \chi(s)=\frac1{1+\exp[-A(s)]},\qquad
 A(s)=\frac1{s-1/2}-\frac1{1-s}.
\]
Its derivative is nonpositive, so \(h'(u)\le1\).

\begin{lemma}[Actual radial scalar derivatives]
\label{lem:f13v7-scalar}
The root in \eqref{eq:f13v7-implicit} is locally
smooth and unique. With all derivatives of \(h\)
evaluated at \(u=p-b(p)\),
\begin{equation}\label{eq:f13v7-bderivatives}
 b'=-\frac{h'}{16-h'},\qquad
 b''=-\frac{256h''}{(16-h')^3}.
\end{equation}
In the inner region \(b=-p/15\). That identity
cannot be extended to the transition region.
\end{lemma}

\begin{proof}
The derivative of the left side of
\eqref{eq:f13v7-implicit} with respect to \(b\)
is \(1-h'/16\ge15/16\). The implicit function
theorem applies; strict monotonicity also gives
uniqueness within the root interval.
Differentiating once gives
\[
 b'+\frac{h'}{16}(1-b')=0,\qquad
 1-b'=\frac{16}{16-h'}.
\]
A second differentiation gives
\[
 b''+\frac1{16}
       \{h''(1-b')^2-h'b''\}=0 ,
\]
which yields \eqref{eq:f13v7-bderivatives}.
When \(h(u)=u\), the original equation reduces
to \(15b+p=0\).
\end{proof}

For interval evaluation of the transition, useful
explicit derivatives are
\[
\begin{split}
 A'&=-(s-1/2)^{-2}-(1-s)^{-2},\\
 A''&=2(s-1/2)^{-3}-2(1-s)^{-3},\\
 \chi'&=\chi(1-\chi)A',\\
 \chi''&=\chi(1-\chi)
           \{(1-2\chi)(A')^2+A''\}.
\end{split}
\]
For \(u\ne0\), with \(s=\sqrt2|u|/\rho\),
\[
 h'=\chi+s\chi',\qquad
 h''=\frac{\sqrt2}{\rho}\operatorname{sgn}(u)
                         (2\chi'+s\chi'').
\]
These follow by ordinary differentiation on
each signed branch. The smooth constant pieces
supply the endpoint extensions.

\subsection{Exact transverse quaternion jets}

Use quaternion coordinates
\(q=(q_0,q_1,q_2,q_3)\), and identify the transverse
pair \((q_1,q_2)\) with \(q_1+iq_2\).
Let
\[
 U_a=\exp(\theta_aE_3),\quad
 U_b=\exp(\theta_bE_3),\quad
 P=U_aU_b=(\cos p,0,0,\sin p).
\]
Take free endpoint paths
\[
 U_a(t)=U_a\exp(tE(v_a)),\quad
 U_b(t)=U_b\exp(tE(v_b)),\quad
 E(v)=v_1E_1+v_2E_2,\quad v=v_1+iv_2.
\]
Set \(v=0\) at a fixed endpoint.
Derivatives below are evaluated at \(t=0\).
Purely transverse endpoint velocities imply
\(P'_0=P'_3=P''_1=P''_2=0\).

For \(p\ne0\), define
\begin{equation}\label{eq:f13v7-coefficients}
 \alpha=\frac{\sin b}{\sin p},\qquad
 \beta=\alpha b',\qquad
 \gamma=\frac{\cos b\,b'\sin p-\sin b\cos p}{\sin^2p}.
\end{equation}
Thus \(\gamma=d\alpha/dp\).

\begin{proposition}[Radial completion jets]
\label{prop:f13v7-jets}
The actual completed path has transverse first jet
and axial second jet
\begin{equation}\label{eq:f13v7-jets}
 B'=(0,\alpha P'_1,\alpha P'_2,0),\qquad
 B''=(\beta P''_0,0,0,\alpha P''_3-\gamma P''_0).
\end{equation}
The coefficient expressions extend smoothly
through \(p=0\) in the inner region, with
\(\alpha=-1/15,\ \beta=1/225,\ \gamma=0\)
at zero.
\end{proposition}

\begin{proof}
Conjugation covariance and the one-source root
equation give the local radial map
\[
 B(P)=\left(\cos b(p),\
       \frac{\sin b(p)}{\sin p}(P_1,P_2,P_3)\right).
\]
On the chosen signed branch
\(P_0=\cos p(t)\). Since \(P'_0=0\), one has
\(p'=0\), and
\[
 p''=-P''_0/\sin p.
\]
Differentiating the radial map once and twice
therefore gives exactly
\eqref{eq:f13v7-jets}. In the inner region,
\[
 \alpha=-\frac{\operatorname{sinc}(p/15)}
                  {15\operatorname{sinc}(p)},
 \quad \operatorname{sinc}(x)=\frac{\sin x}{x}.
\]
Its smooth even extension, together with
\(b'=-1/15\), proves the stated limits.
\end{proof}

The second jet automatically obeys
\[
 B\cdot B''+|B'|^2=0,
\]
as required for a unit quaternion. This identity
is a useful algebraic check; normalizing an
incorrect second jet afterwards would lose
the acceleration information needed below.

\subsection{Right-trivialized lift and the acceleration quadratic form}

Put
\[
 A_0=\cos b\,\alpha,\qquad
 C_0=-(\cos b\,\gamma+\sin b\,\beta),\qquad
 d_0=-A_0\sin p-C_0\cos p .
\]
The symbols \(A_0,C_0\) in this subsection are
scalar coefficients, not the incidence matrix
or the cutoff exponent \(A(s)\).

\begin{theorem}[Complete one-source transverse response]
\label{thm:f13v7-response}
The right-trivialized transverse pivot velocity is
\begin{equation}\label{eq:f13v7-lift}
 \xi_B=\ell(e^{-2i\theta_b}v_a+v_b),
 \qquad \ell=\alpha e^{i(p-b)}.
\end{equation}
The axial right-trivialized acceleration
\(\eta_B=(B^{-1}B'')_3\) is
\begin{equation}\label{eq:f13v7-acceleration}
 \eta_B
 =d_0(|v_a|^2+|v_b|^2)
        +2\operatorname{Re}(\overline{v_a}G_{ab}v_b),
 \qquad
 G_{ab}=-(C_0+iA_0)e^{-i(\theta_a-\theta_b)} .
\end{equation}
Thus it is the Hermitian quadratic form of a
two-endpoint matrix \(G\) with equal diagonals
\(d_0\) and off-diagonal entry \(G_{ab}\).
\end{theorem}

\begin{proof}
The right-trivialized source velocity is
\(e^{-2i\theta_b}v_a+v_b\). In quaternion transverse
coordinates, left multiplication by
\(\exp(pE_3)\) multiplies this complex number
by \(e^{ip}\). Proposition~\ref{prop:f13v7-jets}
then multiplies it by \(\alpha\), while right
trivialization at \(B\) multiplies by \(e^{-ib}\).
This proves \eqref{eq:f13v7-lift}.

Multiplying the second jet in
\eqref{eq:f13v7-jets} by \(B^{-1}\) gives
\[
 \eta_B=A_0P''_3+C_0P''_0 .
\]
Quaternion differentiation of \(P=U_aU_b\)
gives its axial complex second jet as
\[
 P''_0+iP''_3
 =-(|v_a|^2+|v_b|^2)e^{ip}
       -2e^{i(\theta_a-\theta_b)}v_a\overline{v_b}.
\]
Substitution and taking real and imaginary
parts give \(d_0\) and \(G_{ab}\) exactly.
The scalar component of \(B^{-1}B''\) is
\(-|\xi_B|^2\); its axial vector component
is precisely the acceleration beyond the
right-geodesic physical Hessian.
\end{proof}

\subsection{Assembly of the actual Hessians}

Let \(H_\perp^{\rm phys}\) be the physical
transverse Hessian with free links and locally
moving pivots temporarily independent. Let
\(g_e\) be the physical action derivative
along \(E_3\) at pivot \(e\). Assemble \(J_\perp\)
from the identity rows at free links and
\eqref{eq:f13v7-lift} at each active pivot.
Embed its endpoint acceleration matrix \(G_e\)
into the full free-coordinate matrix.

\begin{corollary}[Completed Hessian chain rule]
\label{cor:f13v7-assembly}
On the specified one-source open region,
\begin{equation}\label{eq:f13v7-transverse}
 H_\perp
 =J_\perp^*H_\perp^{\rm phys}J_\perp
           +\sum_e g_eG_e .
\end{equation}
For longitudinal angular variations, let \(r_e\)
be the row summing the free endpoint velocities
of source \(e\), and let \(J_\parallel\) have
pivot row \(b'_er_e\). Then
\begin{equation}\label{eq:f13v7-scalar}
 H_\parallel
 =J_\parallel^{\mathsf T}H_\parallel^{\rm phys}
          J_\parallel+\sum_e g_eb''_er_e^{\mathsf T}r_e .
\end{equation}
\end{corollary}

\begin{proof}
Differentiate the action along the actual
completed curve twice. The physical Hessian
acts on its first velocities. Its gradient
pairs with acceleration beyond the physical
coordinate geodesics. At a commuting field
the physical gradient is axial, so the
additional transverse term is exactly
\(\sum_e g_e\eta_e\), given by
Theorem~\ref{thm:f13v7-response}. The longitudinal
pivot acceleration is
\(b''_e(r_ev)^2\), proving the second formula.
\end{proof}

Even when the \emph{completed} free gradient
vanishes, the individual physical pivot forces
\(g_e\) need not vanish. Dropping the acceleration
terms on that ground would be incorrect.
All matrices in this section use generators
of Hilbert--Schmidt norm \(\sqrt2\); their
scalar Rayleigh values must be divided by two
to obtain unit Hilbert--Schmidt curvature.

\subsection{Exact algebra checks, distinguished from candidate diagnostics}

The checker
\path{scripts/verify_ym_f13_v7_radial_jet_identities.py}
uses only rational quaternion arithmetic.
It verifies the unit-quaternion second-order
constraint, the complex lift, and the
acceleration form on 64 rational jet cases,
including inner-type, transition-type, and
zero-response coefficients. It also checks
both implicit scalar derivative identities
on 12 rational pairs \((h',h'')\).
The universal results are proved algebraically
above; these finite cases are regression tests,
not a substitute for the proofs.

The separate numerical branch search starts
from central ground codes
\[
 772,\ 226,\ 2274,\ 8418,\ 10466,\ 16610,\ 24802,\ 28421 .
\]
It perturbs each along its V4 certified
negative vector and applies a bounded
commuting descent with actual pivot solves.
A safeguarded bisection replaces an
exploratory Newton root solve that failed
its residual check at one trial point.
The inherited certified files are unchanged.

Five runs approach the already known action
\(569.684411864\ldots\). Three runs approach
slightly higher values, representing two
numerically distinct transition branches.
Rounded angle vectors for representatives
772 and 28421 are stored on denominator
\(10^{10}\) in
\path{scripts/ym_f13_v7_candidate_data.py}.
The diagnostics at those rounded vectors are:
\begin{center}\small
\begin{tabular}{lrr}
\toprule
quantity, floating-point only&code 772&code 28421\\
\midrule
action&569.68444637&569.68442506\\
maximum angular gradient&\(1.9\,10^{-7}\)&\(3.1\,10^{-7}\)\\
smallest longitudinal eigenvalue&0.2582&0.2578\\
first transverse eigenvalue after orbit mode&0.1271&0.1714\\
transition coordinate \(s\)&0.80699&0.79271\\
transition \(b'\)&0.16815&0.20706\\
transition \(b''\)&217.24&179.39\\
\bottomrule
\end{tabular}
\end{center}
Each has four active source paths, with one
in the transition and three in the inner
region. The near-zero transverse eigenvalues
are about \(-5.9\,10^{-9}\) and
\(-1.5\,10^{-8}\), respectively. These tiny
negative values at nonstationary rounded
points do not certify negative directions
at true stationary points, nor do the other
positive numerical eigenvalues prove normal
positivity at such points.

The code
\path{scripts/explore_ym_f13_v7_transition_hessian.py}
checks its scalar gradient against the
original completed-action gradient, its
scalar Hessian against directional gradient
differences, and the transverse conjugation
identity
\[
 w^*H_\perp w=2\,\nabla S\cdot\sin(2\theta),
 \qquad w_j=e^{-2i\theta_j}-1 .
\]
The observed scalar finite-difference errors
are below \(1.4\,10^{-8}\), and conjugation
identity residuals below \(3\,10^{-14}\).
These are consistency diagnostics, not
interval bounds. No candidate minimum is
accepted on their basis.

\subsection{Changed acquisition priority and preservation record}

The immediate priority is now to certify
or refute the additional candidate stationary
orbits, before spending more versions on
central-pattern escape. Required checks include:
\begin{enumerate}
\item interval containment of all source gates
and the actual pivot roots near each candidate;
\item a genuine stationary-point inclusion
argument for all 944 longitudinal coordinates;
\item full normal positivity, or a certified
negative direction, using the transition
terms in \eqref{eq:f13v7-transverse};
\item separation from the existing orbit and
the consequences for global well coverage.
\end{enumerate}
If additional wells are certified, their
higher action alone cannot be used as an
escape estimate. The global assembly must
account for their trapping behaviour rather
than presupposing that the original tube
contains all relevant wells.

The frozen \emph{unverified} candidate data
have SHA--256
\begin{center}\scriptsize\ttfamily
24116F318E69CCEDEAE2E07F2130A3A59955D2054B15F33B5BE7C9BA6C4B6982.
\end{center}
V6 contained 117,966 bytes and SHA--256
\begin{center}\scriptsize\ttfamily
59CD7B7F821D6852F4E6862641F7B2CE38031AA6579354542D69D47683EEB434.
\end{center}
Its complete mathematical body is retained,
with only the version title updated. Active
V7 replaces V6. Archives and companion
notes remain untouched; the next companion
and direction reviews remain V10.

\begin{firewall}
This iteration proves the cutoff-transition
Hessian identities and produces evidence
that changes the next research action.
It does not prove additional local minima,
disprove a uniform conditional gap, establish
global ground coverage, or solve the continuum
Yang--Mills problem. Those conclusions require
the stated validation and assembly steps.
The full mass-gap objective remains unproved
and active.
\end{firewall}
\section{V8: a certified secondary well and failure of a uniform hidden diffusion gap}
\label{sec:f13v8}

\subsection{Outcome and its scope}

The transition candidate labelled 772 in V7 is now
certified: the actual fixed-parent action has a
second strict local-minimum conjugation orbit
\(\mathcal O_2\), with energy
\[
 569.68444636<E_2<569.68444638 .
\]
It is distinct from the orbit \(\mathcal O_1\)
certified in f12 V89--V90, whose energy satisfies
\[
 569.684411845\le E_1\le569.684411883 .
\]
In particular \(E_2>E_1\); the second orbit is
not a global minimum. Global minimality of the
first orbit remains unproved.

The secondary well makes a stronger conclusion
possible: for the actual fixed-parent conditional
\emph{diffusion}, a positive spectral gap uniform
as \(\beta\to\infty\) is false. A corresponding
uniform killed floor outside a sufficiently
small tube about \(\mathcal O_1\) is also false.
This realizes, in the actual parent, the
metastability mechanism illustrated only by an
abstract circle example in f12 V73. The new
content is the certified actual well and its
application, not that abstract example.

These are statements about an auxiliary
finite-parent differential generator. They do
not disprove a heat-bath gap, a retained-variable
gap, or the physical Yang--Mills mass gap.
They require an upstream change in the proposed
uniform hidden-diffusion comparison, not a change
in the ultimate objective.

\subsection{Rational centre and outward interval construction}

The angular centre \(\theta^c\in\mathbb R^{944}\)
is the literal integer vector in
\path{scripts/ym_f13_v8_center_data.py}, divided
by \(10^{16}\). It was proposed by numerical
Newton refinement, but no accuracy assumption
about that refinement is used in acceptance.
The first validation box is
\[
 |\theta_i-\theta_i^c|\le10^{-9}\quad(0\le i<944).
\]
It contains the Euclidean angular ball of
radius \(R=10^{-9}\).

All interval endpoints are integers divided
by \(2^{70}\). Arithmetic operations round
outwards with integer division. The new helper
\path{scripts/ym_f13_v8_transition_intervals.py}
adds the following validated elementary
function operations to the retained interval
machinery.
\begin{enumerate}
\item A rational enclosure for \(\pi\) follows
from
\(\pi=16\arctan(1/5)-4\arctan(1/239)\),
using 60 alternating-series terms and the
next-term error bound for each arctangent.
The identity follows from tangent addition
and the corresponding angle ranges.
\item Sine and cosine are reduced by an
integer multiple of \(\pi/2\) to an interval
with \(|r|\le1\). Sixteen even or odd Taylor
terms give remainder bounds
\(|r|^{32}/32!\) and \(|r|^{33}/33!\).
The interval reduction itself includes the
uncertainty in \(\pi\).
\item For \(|x|\le16\), evaluate \(e^{x/64}\)
through degree 24 with error
\(2|x/64|^{25}/25!\), then square six times.
The factor two bounds \(e^{|x|/64}\) on this
range. The exponent bound is explicitly
checked wherever the transition formula
uses it.
\item The enclosure for \(\sqrt2\) is obtained
by an integer square root of \(2(2^{70})^2\).
\end{enumerate}
The polynomial evaluations, remainder additions,
range checks, and repeated products all use
outward interval arithmetic. There is no
floating-point elementary function in the
acceptance path.

\subsection{All source gates and the actual pivot roots}

The four active edges remain
\[
 e=((a,0,0,b),1),\qquad a,b\in\{0,1\}.
\]
At each edge the active source is path index
11 in the retained fourteen-path ordering.
There are three inner gates and one genuine
transition gate at \(e=((1,0,0,1),1)\).

The centre root is enclosed by the exact
inner formula when applicable, and otherwise
by interval bisection of
\eqref{eq:f13v7-implicit}. A root lies in
\([-1/1024,1/1024]\), because
\(|h|\le\rho/\sqrt2\) and
\(|b|=|h|/16<1/1024\). Endpoint signs are
checked. If an interval evaluation at a
bisection midpoint contains zero, the bound
\(F_b\ge15/16\) contracts the root enclosure
to the midpoint plus or minus
\((16/15)\sup|F|\). This encloses the root
for every source angle in its input interval.

For box variation, V7's derivative formula
implies \(|b'|\le1\): for \(h'\le0\) the
ratio \(-h'/(16-h')\) is below one, and
for \(0\le h'\le1\) its magnitude is at
most \(1/15\). Thus a source-angle uncertainty
of at most \(2R\) gives at most \(2R\) root
uncertainty. Inner roots use the sharper
exact formula \(-p/15\).

The resulting gate enclosures on the full
box are contained in the following rational
decimal intervals:
\begin{center}\small
\begin{tabular}{lll}
\toprule
edge&relative scaled angle \(s\)&gate\\
\midrule
\(((0,0,0,0),1)\)&[0.13323494,\,0.13323535]&inner\\
\(((0,0,0,1),1)\)&[0.47389468,\,0.47389508]&inner\\
\(((1,0,0,0),1)\)&[0.29560551,\,0.29560591]&inner\\
\(((1,0,0,1),1)\)&[0.80698854,\,0.80698928]&transition\\
\bottomrule
\end{tabular}
\end{center}
Every centre root enclosure has width at
most \(10^{-20}\). Every other source path,
relative to its proposed pivot, has squared
Hilbert--Schmidt chord distance exceeding
\(\rho^2\) by at least
\[
 \frac{161678163}{10^{12}}>0 .
\]
Since geodesic distance is at least chord
distance, every such cutoff is identically
zero on the declared box.

These checks also justify use of the
one-active-source completion. Its root
solves the full scalar equation once the
other source terms are zero. The full
commuting root equation has derivative at
least \(1-14/16=1/8\), so this is its unique
root in the chart. Strict gate margins give
a full group neighbourhood of each checked
configuration in which the same one-source
formula is smooth. For one source, the
matrix root equation forces \(B\) to commute
with \(PB^{-1}\), hence with \(P\), so the
radial completion used in V7 is the actual
group completion, not an extra constraint.

\subsection{Certified longitudinal stationary-point inclusion}

The scalar interval matrix includes both
terms of \eqref{eq:f13v7-scalar}. In direct
plaquette notation it is
\[
 H_\parallel
 =\sum_Q2\cos(\phi_Q)\,a_Q^{\mathsf T}a_Q
       +\sum_e g_eb''_er_e^{\mathsf T}r_e ,
\]
where \(a_Q\) contains all free and induced
pivot angular velocities. The physical
pivot forces \(g_e\), transition derivatives,
and their uncertainties are retained.

\begin{lemma}[Whole-box scalar certificate]
\label{lem:f13v8-scalar}
Throughout the declared box,
\[
 H_\parallel\ge
 \mu_\parallel I,\qquad
 \mu_\parallel=\frac{1855919}{10^8}.
\]
At the rational centre,
\[
 \|\nabla S(\theta^c)\|_2
 \le\frac{27}{1250000000000000}.
\]
\end{lemma}

\begin{proof}
The interval entries are assembled by
\begin{center}\small
\path{scripts/verify_ym_f13_v8_transition_stationary.py}.
\end{center}
The positive-matrix certificate is the
retained integer factor/inverse method of
f12 V89, applied to the \emph{whole-box}
matrix enclosure, not only its midpoint.
For completeness, let \(R_0\) be the proposed
triangular factor and \(K_0\) its proposed
inverse, in their declared integer scales.
Exact bounds
\(\|I-K_0R_0\|_\infty<1\) and
\(\|I-K_0R_0\|_1<1\) bound both norms of
\(R_0^{-1}\). Consequently
\[
 \lambda_{\min}(R_0^{\mathsf T}R_0)
 \ge\frac1{\|R_0^{-1}\|_1\|R_0^{-1}\|_\infty}.
\]
Subtracting the symmetric interval error
from the factor gives a lower bound for
every enclosed Hessian.

Here the certified factor floor and error obey
\[
 \lambda_{\rm factor}\ge\frac{1084029}{50000000},
 \qquad
 \varepsilon_{\rm factor}\le\frac{1560693723}{500000000000}.
\]
Rounding the resulting lower bound down
gives \(\mu_\parallel\). The gradient bound
comes from the outward interval sum of
all 944 squared gradient-component bounds
and an integer square root. All integer
matrix products have checked absolute
dot-product bounds below \(2^{63}\);
no unguarded integer overflow or numerical
factorization accuracy is assumed.
\end{proof}

\begin{theorem}[A distinct actual stationary point]
\label{thm:f13v8-stationary}
There is a unique stationary point
\(\theta^{(2)}\) in the commuting angular
ball of radius \(10^{-9}\) about
\(\theta^c\), and
\[
 \|\theta^{(2)}-\theta^c\|_2
 \le\frac{29}{25000000000000}
 <2\,10^{-12}.
\]
It is stationary for the full completed
\(\SU(2)^{944}\) action, not merely for its
commuting restriction. Its energy obeys
\[
 569.68444636<S(\theta^{(2)})<569.68444638.
\]
\end{theorem}

\begin{proof}
Write \(\varepsilon\) for the checked
centre-gradient bound. On the boundary
of the radius-\(R\) ball, strong convexity
gives
\[
 \nabla S(\theta^c+v)\cdot v
 \ge\mu_\parallel R^2-\varepsilon R>0.
\]
A minimum on the closed ball therefore
lies in its interior. Strong convexity
gives uniqueness and the distance bound
\(\varepsilon/\mu_\parallel\); the stated
rounded distance is obtained using the
unrounded certified gradient enclosure.

Global conjugation by \(\exp(tE_3)\)
fixes the commuting field, the central
exterior, and coarse field \(I\). Its
action on the transverse tangent planes
is a nontrivial rotation. Invariance of
the action forces every transverse first
derivative to vanish at a commuting field.
The longitudinal derivatives vanish at
\(\theta^{(2)}\), proving full stationarity.

The centre energy is evaluated with
outward intervals. A whole-box Hessian
norm bound of \(73912697/10^6\) gives
\[
 |S(\theta^{(2)})-S(\theta^c)|
 \le\varepsilon R+\frac12(73912697/10^6)R^2 .
\]
The resulting interval lies strictly
inside the stated energy interval.
\end{proof}

\subsection{Full normal positivity, including the transition acceleration}

For the transverse test use the smaller
angular box
\[
 |\theta_i-\theta_i^c|\le2\,10^{-12},
\]
which contains \(\theta^{(2)}\) by the
preceding theorem. The interval matrix
uses \eqref{eq:f13v7-transverse}, with all
radial coefficients and physical pivot
forces enclosed on that box.

\begin{theorem}[A second strict local-minimum orbit]
\label{thm:f13v8-minimum}
The conjugation orbit \(\mathcal O_2\)
through the certified point is a smooth
two-sphere. Its full Hessian has exactly
the two real orbit directions in its
kernel and satisfies
\[
 \Hess S|_{N\mathcal O_2}
 \ge\frac{31}{125000}I
\]
in Hilbert--Schmidt normal coordinates.
Thus \(\mathcal O_2\) is a strict local
minimum modulo conjugation.
\end{theorem}

\begin{proof}
The exact verifier
\path{scripts/verify_ym_f13_v8_transition_minimum.py}
first reruns the longitudinal inclusion.
It then constructs the Hermitian transverse
interval matrix on 944 complex coordinates.
Conjugate entries are enclosed in common
Hermitian hulls, and all transition lift
and acceleration terms are included.

Pin the complex coordinate with index
156. The checked field has sine magnitude
greater than \(1/2\) there. Its conjugation
tangent
\(w_i=e^{-2i\theta_i^{(2)}}-1\) therefore
has \(w_{156}\ne0\). Removing that complex
coordinate gives a real symmetric matrix
of dimension 1886. The exact integer
factor certificate gives, on the entire
small box,
\[
 H_{\perp,\mathrm{pin}}\ge
       \frac{31}{62500}I
\]
in angular normalization. The factor
floor and error obey
\[
 \lambda_{\rm factor}\ge\frac{51899}{10^8},
 \qquad
 \varepsilon_{\rm factor}\le\frac{5745403}{250000000000}.
\]
All factor and inverse checks use bounded
integer products, as in the scalar test.

At the certified stationary point,
invariance gives the exact kernel vector
\(H_\perp w=0\). Interval containment of
the corresponding products is checked
only as a consistency test; the kernel
identity itself follows from stationarity
and symmetry. For \(v\perp w\), subtract
\((v_{156}/w_{156})w\) to obtain a pinned
vector \(u\). Then
\[
 v^*H_\perp v=u^*H_\perp u
 \ge(31/62500)\|u\|^2
 \ge(31/62500)\|v\|^2 .
\]
Conversely, the pinned positivity excludes
any additional transverse kernel or
negative direction.

The residual torus symmetry separates
longitudinal and transverse Hessian
blocks. The scalar floor from
Lemma~\ref{lem:f13v8-scalar} is larger
than the pinned transverse floor.
Conversion from angular generators of
Hilbert--Schmidt norm \(\sqrt2\) divides
these floors by two, giving \(31/125000\).
The noncentral pinned link makes the
common stabilizer exactly \(\mathrm U(1)\);
the orbit is \(\SU(2)/\mathrm U(1)\cong S^2\).
Smoothness, compactness of the orbit, and
the positive normal Hessian give the
strict local-minimum conclusion.
\end{proof}

The energy intervals imply
\[
 E_2-E_1>\frac{34477}{10^9}>0.
\]
This proves that the two orbits are
distinct and that \(\mathcal O_2\) is
not globally minimizing. Their gate
types also distinguish them: the new
orbit has a transition source with
\(s>0.8\), whereas the old certified
orbit has only inner active sources.
Neither distinction is a physical
mass-gap estimate.

\subsection{Actual metastability obstructs the uniform diffusion target}

Let
\[
 d\mu_\beta=Z_\beta^{-1}e^{-\beta S}a\,d{\rm vol}_M,
 \qquad M=\SU(2)^{944},
\]
with the actual positive smooth,
\(\beta\)-independent amplitude \(a\).
Let \(\lambda_\beta\) be the spectral
gap for the Dirichlet form
\(\E_\beta(f)=\int|\nabla f|^2d\mu_\beta\).

\begin{theorem}[No temperature-uniform hidden diffusion gap]
\label{thm:f13v8-metastability}
There are constants \(C<\infty\),
\(\delta>0\), and \(\beta_0<\infty\),
independent of \(\beta\), such that
\begin{equation}\label{eq:f13v8-gap-upper}
 0<\lambda_\beta\le
       C\beta^{1415}e^{-\beta\delta},
 \qquad\beta\ge\beta_0 .
\end{equation}
In particular no positive lower bound
uniform in \(\beta\) can hold for this
fixed-parent conditional diffusion.
\end{theorem}

\begin{proof}
The normal dimension of either orbit
is \(2832-2=2830\).
Choose a sufficiently small normal
tube about \(\mathcal O_2\), disjoint
from \(\mathcal O_1\). Its positive
normal Hessian gives constants
\(c_-,c_+>0\) such that, in that tube,
\[
 E_2+c_-|z|^2\le S(q,z)
                 \le E_2+c_+|z|^2 .
\]
No global trivialization of the normal
bundle is required: these uniform
bounds follow in finitely many tubular
charts. The actual amplitude and
volume Jacobians have positive lower
and finite upper bounds there.

Take a fixed smooth \(f\), equal to
one on the radius-\(r\) inner tube
and zero outside the radius-\(2r\)
tube. Its derivative is supported
where \(S\ge E_2+\delta\), with
\(\delta=c_-r^2>0\). Thus
\[
 \E_\beta(f)
 \le Z_\beta^{-1}C_1e^{-\beta(E_2+\delta)}.
\]
Integration over normal balls of
radius \(\beta^{-1/2}\), in a fixed
positive-volume portion of the orbit,
gives
\[
 \int f^2d\mu_\beta
 \ge Z_\beta^{-1}c_1\beta^{-1415}
                                e^{-\beta E_2}.
\]
The lower-energy certified orbit
\(\mathcal O_1\) similarly supplies
\[
 Z_\beta\ge c_2\beta^{-1415}
                             e^{-\beta E_1}.
\]
As \(S\ge E_2\) on the support tube,
\[
 \mu_\beta(\operatorname{supp}f)
 \le C_2\beta^{1415}
             e^{-\beta(E_2-E_1)}\longrightarrow0 .
\]
No assumption that \(E_1\) is the
global minimum was used.

By Cauchy--Schwarz,
\[
 \Var_{\mu_\beta}f
 \ge\bigl(1-\mu_\beta(\operatorname{supp}f)\bigr)
                         \int f^2d\mu_\beta .
\]
For sufficiently large \(\beta\) the
parenthesis is at least \(1/2\).
The Rayleigh quotient of this test
function proves the upper bound.
For each fixed finite \(\beta\), the
positive smooth density on the compact
connected manifold gives a positive
elliptic diffusion gap. This proves
\eqref{eq:f13v8-gap-upper}.
\end{proof}

The barrier \(\delta\) in this theorem
comes from a small normal tube and is
proved positive but not numerically
enclosed here. It is \emph{not} the
energy difference \(E_2-E_1\). The
secondary well's exponentially small
mass does not remove its small
Rayleigh quotient.

\begin{corollary}[Failure of the single-tube killed-escape target]
\label{cor:f13v8-killed}
Let \(T_1\) be any sufficiently small
closed tube about \(\mathcal O_1\)
disjoint from a neighbourhood of
\(\mathcal O_2\). The lowest Dirichlet
eigenvalue on \(M\setminus T_1\) satisfies
\[
 \lambda_\beta^D(M\setminus T_1)
 \le C\beta^{1415}e^{-\beta\delta}
\]
for sufficiently large \(\beta\).
\end{corollary}

\begin{proof}
Choose the preceding \(f\) with support
in the secondary tube and away from
\(T_1\). It belongs to the relevant
Dirichlet domain. The uncentred
quotient \(\E_\beta(f)/\|f\|_2^2\)
already has the asserted bound.
\end{proof}

\subsection{What this changes, and what remains available}

V4--V6's local central saddle estimates
remain valid. They cannot be extended
to a temperature-uniform killed floor
on the entire complement of the original
single tube: that complement contains
the newly certified noncentral well.
Adding a second local tube repairs
the coverage description but cannot
restore a uniform gap for the same
diffusion; Theorem~\ref{thm:f13v8-metastability}
already excludes it.

The conclusion does not automatically
transfer to a nonlocal heat-bath update.
The distinction was explicit in f12
V73 and V80. In particular, f12 V80's
actual visible-cylinder comparison
and pre-integration localization did
not require a hidden global Poincare
inequality. Its constant interface
saddle floor under whole-layer
clearance remains an available,
precisely scoped input.

The next upstream task is therefore
to examine the retained-variable
repair and its actual allowed update
mechanism, or to quantify a comparison
that tolerates the now-proved hidden
metastability. It is no longer useful
to pursue a false uniform hidden
diffusion estimate by improving its
constants. The second V7 candidate
labelled 28421 remains unverified;
its validation is not needed for
the obstruction just proved.

\subsection{Certificate ledger and append-only record}

The accepted checks are
\path{scripts/verify_ym_f13_v8_transition_stationary.py}
and
\path{scripts/verify_ym_f13_v8_transition_minimum.py}.
The second reruns the first and verifies
all 1886 pinned real transverse
directions. Both report PASS. No
optimizer, floating-point eigenvalue,
or exploratory module is imported
by their acceptance path.

The frozen rational centre has SHA--256
\begin{center}\scriptsize\ttfamily
E73E6457EF28677A0389FBF2C1572EC972DE94F5FAF6AA85624A864AD2B0041C.
\end{center}
The new transition interval helper has SHA--256
\begin{center}\scriptsize\ttfamily
75F7B999482BD2F5CF91BE44D2721F4335127C04EA27051E58F9231DCEFA5DAC.
\end{center}
The predecessor V7 contained 132,301 bytes
and SHA--256
\begin{center}\scriptsize\ttfamily
066D7D50D1718A79C887DFD41FE4BEB101A0F132468600635E847A270BBEC124.
\end{center}
Its entire body is retained, with
only the title updated. Active V8
replaces active V7. Archives and
companion notes remain unchanged;
the next companion and direction
reviews remain V10.

\begin{firewall}
A second actual strict local-minimum
orbit and the resulting failure of
a temperature-uniform hidden diffusion
gap are proved for this specified
finite parent, exterior, and cutoff.
This is a verified obstruction to one
sufficient comparison strategy, not
a counterexample to four-dimensional
Yang--Mills mass generation. No
physical Hamiltonian gap, continuum
construction, unrestricted repair
estimate, or cofinal closure has
been established. The full goal
remains active, with the next work
moved upstream as described.
\end{firewall}
\section{V9: secondary-well coarse variation and marginal comparison}
\label{sec:f13v9}

\subsection{The upstream question after the diffusion obstruction}

V8 proved a secondary well and thereby ruled out
one proposed uniform hidden-diffusion estimate.
It did not show that the well obstructs the
retained marginal or its actual heat-bath updates.
This iteration addresses that distinction.

There are two new conclusions. First, the
certified secondary hidden minimum is regular
under a specified coarse-edge variation of the
completed parent action. Second, the contribution
localized near that well can be removed as an
\emph{analytic comparison} with exponentially
small error in the retained potential and any
fixed number of its derivatives, on a sufficiently
small coarse neighbourhood. A precise form
comparison then applies to retained diffusion
and heat-bath gaps, if the comparison law's
gap can be established.

No model is silently replaced. The original
partition function is retained exactly, and
the relation to the comparison partition
function is explicit. No hidden Poincare
inequality enters the argument. The derivative
control below is proved; it is not inferred
from rarity alone, the pitfall identified
already in f12 V31.

\subsection{The secondary orbit is not a joint parent/coarse critical orbit}

Use the same coarse edge as f12 V92 and
V2 of this lineage:
\[
 e_0=((0,0,0,0),2),\qquad
 \ell_0=((0,0,1,0),2).
\]
The latter is its physical pivot. Apply
the inactive-edge calculation to the
new V8 stationary orbit \(\mathcal O_2\).

\begin{proposition}[Certified coarse escape from the secondary hidden well]
\label{prop:f13v9-coarse}
At the representative of \(\mathcal O_2\)
with commuting axis \(E_3\), the angular
coarse force \(F_2\) at \(C_{e_0}=I\) obeys
\[
 \frac{1960605453}{200000000}
 \le F_2\le
 \frac{4901513633}{500000000},\qquad
 F_2>\frac{49}{5}.
\]
All fourteen source paths for this edge
have squared Hilbert--Schmidt chord
distance from \(I\) at least
\(1241904559/10^9>1\).
For \(d_{\rm HS}(C_{e_0},I)\le1/4\),
the actual pivot equals \(C_{e_0}\).
Holding the hidden field fixed, the
coarse endpoint
\[
 C_{e_0}=\frac{255}{257}I
                       -\frac{32}{257}E_3
\]
decreases the completed parent action
by at least \(1448/1285>9/8\).
The joint parent/coarse gradient has
Hilbert--Schmidt norm greater than six
along the entire conjugation orbit.
\end{proposition}

\begin{proof}
The checker
\path{scripts/verify_ym_f13_v9_secondary_retained.py}
reruns V8's stationary-point inclusion and
evaluates the six physical plaquettes
incident to \(\ell_0\), with all actual
other pivot values enclosed. Their
signed sine sum gives the stated force;
their cosine sum \(A_2\) satisfies
\(|A_2|\le12\).

The source distance bound and the
triangle inequality give
\[
 d(P_j,C_{e_0})\ge1-\frac14>\rho .
\]
Thus every source cutoff is zero at
the candidate pivot \(B=C_{e_0}\).
Uniqueness of completion proves that
this is the actual pivot. Source paths
for other pivots do not contain this
pivot, so those completed values do
not change under this coarse variation.

For an axis-aligned endpoint
\(C_{e_0}=\cos t\,I+\sin t\,E_3\),
the exact action difference is
\[
 S(C_{e_0},x)-S(I,x)
       =A_2(1-\cos t)+F_2\sin t .
\]
The rational endpoint has
\(t=-2\arctan(1/16)\), so
\(|t|<1/8\) and
\(\sqrt2|t|<1/4\). Its decrease is at
least
\[
 \frac{49}{5}\frac{32}{257}
       -12\left(1-\frac{255}{257}\right)
       =\frac{1448}{1285}.
\]
The angular generator has
Hilbert--Schmidt norm \(\sqrt2\).
Hence the dual force norm is
\(F_2/\sqrt2>6\). Conjugation
equivariance transports all statements
to the other points of \(\mathcal O_2\).
\end{proof}

This reuses f12 V92's inactive-edge
identity with a newly certified field;
it does not repeat its old-orbit
certificate. It is regularity of
this parent action when a coarse
input is released. It is not an
all-hidden visible-cylinder estimate.
Orbit averaging can cancel vector
forces, and forces from other parent
terms or direct coarse interactions
must still be included in a full
retained calculation.

\subsection{The actual parameter family and a fixed localized cutoff}

Let \(q\) denote the 96 incident
coarse links, near their value
\(q_0=(I,\ldots,I)\). Keep the physical
fine exterior from V8 fixed. Integrate
the same 944 hidden free groups:
\[
 Z_\beta(q)=\int_M
       e^{-\beta S(q,x)}a(q,x)\,d{\rm vol}_M(x),
 \qquad M=\SU(2)^{944}.
\]
This is the actual partial coarse
partition function with the unique
completed pivots. Smooth joint
completion and its strictly positive
inverse Haar Jacobian, recalled in
Section 2, make \(S\) and \(a>0\)
smooth on a common compact coarse
neighbourhood times \(M\). The
amplitude is independent of \(\beta\).

Choose a smooth \(0\le\chi\le1\),
equal to one near \(\mathcal O_2\),
supported in a sufficiently small
normal tube \(T_2\) from V8. Choose
that tube disjoint from a tube of
\(\mathcal O_1\), and so that
\[
 S(q_0,x)\ge E_2
          \quad\hbox{on }\operatorname{supp}\chi .
\]
The cutoff can be chosen invariant
under the remaining global conjugation
symmetry, using invariant tubular
distance. Crucially, \(\chi\) is
\emph{fixed as \(q\) varies}. There
are no moving-boundary or differentiated
basin-selection terms.

Define
\[
\begin{split}
 N_\beta(q)&=\int_M\chi(x)
          e^{-\beta S(q,x)}a(q,x)\,d{\rm vol}_M(x),\\
 \widetilde Z_\beta(q)&=\int_M(1-\chi(x))
          e^{-\beta S(q,x)}a(q,x)\,d{\rm vol}_M(x),\\
 p_\beta(q)&=\frac{N_\beta(q)}{Z_\beta(q)}.
\end{split}
\]
Then, exactly,
\begin{equation}\label{eq:f13v9-exact}
 \widetilde Z_\beta=Z_\beta(1-p_\beta).
\end{equation}
The comparison removes a localized
contribution, not a claimed complete
basin or every possible secondary
well. It remains positive because
\(1-\chi=1\) near \(\mathcal O_1\).
The original \(Z_\beta\) is not changed.

\subsection{Uniform local suppression from the certified energy separation}

Put
\[
 \gamma=\frac1{200000},\qquad
 \sigma=\frac1{50000}.
\]
Shrink a coarse coordinate neighbourhood
\(U\) of \(q_0\) until
\begin{equation}\label{eq:f13v9-budget}
 |S(q,x)-S(q_0,x)|\le\gamma
             \quad(q\in\overline U,\ x\in M).
\end{equation}
Such a neighbourhood exists by joint
smoothness and compactness. For example,
on an initial compact chart take a
finite upper bound \(L\ge1\) on the
coarse derivative norm and restrict
the coordinate displacement to
\(\gamma/L\). Neither \(L\) nor a
numerical radius for \(U\) is claimed
to be enclosed here.

\begin{theorem}[Localized well mass, without hidden mixing]
\label{thm:f13v9-mass}
There are \(C<\infty\) and
\(\beta_0<\infty\), independent of
\(q\in\overline U\), such that
\begin{equation}\label{eq:f13v9-mass}
 0<p_\beta(q)\le
       C\beta^{1415}e^{-\sigma\beta},
 \qquad \beta\ge\beta_0 .
\end{equation}
\end{theorem}

\begin{proof}
On the support of \(\chi\),
\(S(q,x)\ge E_2-\gamma\).
Compactness and an amplitude upper
bound give
\[
 N_\beta(q)\le C_1e^{-\beta(E_2-\gamma)}.
\]
Near \(\mathcal O_1\), the fixed
normal tube at \(q_0\) has
\(S(q_0,x)\le E_1+C_2|z|^2\).
Integrating over normal balls of
radius \(\beta^{-1/2}\), with a fixed
positive-volume orbit patch, yields
\[
 Z_\beta(q)\ge c_1\beta^{-1415}
                        e^{-\beta(E_1+\gamma)}.
\]
The same lower bound applies to
\(\widetilde Z_\beta\). The exponent
1415 is half the normal dimension
2830. Smooth positive amplitudes and
tubular Jacobians have uniform bounds
on this fixed compact family.

V8 gives \(E_2-E_1>34477/10^9\), and
\[
 \frac{34477}{10^9}-2\gamma
       >\frac1{50000}=\sigma .
\]
Division proves \eqref{eq:f13v9-mass}.
Global minimality of \(\mathcal O_1\)
and a hidden spectral gap were not
used.
\end{proof}

The separation exponent \(\sigma\)
is not V8's diffusion barrier
\(\delta\). The constants and the
large-\(\beta\) threshold in this
theorem are existential. In particular,
the estimate is not a numerical
claim of small error at a specified
simulation coupling.

\subsection{Derivative control for the normalized localized mass}

Use fixed smooth coarse coordinates
on \(U\), and write
\[
 w_\beta(q,x)=e^{-\beta S(q,x)}a(q,x),
 \qquad
 s_a=\partial_a\log w_\beta .
\]
For constants \(B_1,B_2<\infty\)
independent of \(q,x,\beta\),
\[
 |s_a|\le B_1(\beta+1),\qquad
 |\partial_bs_a|\le B_2(\beta+1).
\]
These include derivatives of the
actual amplitude and completion.

\begin{lemma}[Relative derivative bounds]
\label{lem:f13v9-relative}
For each fixed multi-index \(\alpha\),
there is \(C_\alpha<\infty\) such that
\[
 |\partial^\alpha p_\beta(q)|
 \le C_\alpha(\beta+1)^{|\alpha|}p_\beta(q).
\]
In particular, writing \(t=\beta+1\),
\[
 |\partial_ap_\beta|\le2B_1t\,p_\beta,\qquad
 |\partial_{ab}p_\beta|
       \le(2B_2+6B_1^2)t^2p_\beta .
\]
\end{lemma}

\begin{proof}
Derivatives of \(w_\beta\), divided
by \(w_\beta\), are finite sums of
products of derivatives of
\(-\beta S+\log a\), of total degree
at most \(|\alpha|\) in \(\beta\).
On the compact parameter family this
gives
\[
 |\partial^\alpha Z_\beta|
 \le M_\alpha t^{|\alpha|}Z_\beta,\qquad
 |\partial^\alpha N_\beta|
 \le M_\alpha t^{|\alpha|}N_\beta .
\]
The second inequality uses the fixed
nonnegative cutoff; differentiating
it contributes nothing. Differentiation
under the integral is justified by
compactness and smoothness for each
finite \(\beta\).

Differentiate \(Z_\beta p_\beta=N_\beta\)
and induct on \(|\alpha|\).
After dividing by \(Z_\beta\), every
lower-order product is bounded by
a constant times \(t^{|\alpha|}p_\beta\).
This proves the general statement.

For the explicit first two bounds,
use
\[
 p_a=N_a/Z-pZ_a/Z,
\]
and
\[
 p_{ab}=N_{ab}/Z-pZ_{ab}/Z
                     -p_aZ_b/Z-p_bZ_a/Z .
\]
Here
\(|w_{ab}|/w\le(B_2+B_1^2)t^2\).
The first bound and these identities
give \(2B_2+6B_1^2\) in the second.
In particular the normalization
derivatives have not been omitted.
\end{proof}

\begin{theorem}[Exponentially small effective-action error in every fixed \(C^k\)]
\label{thm:f13v9-potential}
Let
\(W_\beta=-\log Z_\beta\) and
\(\widetilde W_\beta=-\log\widetilde Z_\beta\).
For each fixed integer \(k\ge0\),
there are finite \(C_k,\beta_k\) with
\[
 \|W_\beta-\widetilde W_\beta\|_{C^k(\overline U)}
 \le C_k\beta^{1415+k}e^{-\sigma\beta},
 \qquad\beta\ge\beta_k .
\]
No uniform bound as \(k\to\infty\)
is asserted.
\end{theorem}

\begin{proof}
Equation~\eqref{eq:f13v9-exact} gives
\[
 W_\beta-\widetilde W_\beta
                         =\log(1-p_\beta).
\]
Choose \(\beta_k\) large enough that
\(p_\beta\le1/2\) uniformly. Derivatives
of this logarithm are finite sums
of products of derivatives of \(p_\beta\),
divided by powers of \(1-p_\beta\).
Lemma~\ref{lem:f13v9-relative} bounds
each product of total order \(j\)
by \(C_jt^jp_\beta^m\), with \(m\ge1\).
Since \(p_\beta^m\le p_\beta\), the
claim follows from
Theorem~\ref{thm:f13v9-mass}.

For clarity, the explicit bounds
through order two are
\begin{align*}
 |\log(1-p)|&\le2p,\\
 |\partial_a\log(1-p)|&\le4B_1t\,p,\\
 |\partial_{ab}\log(1-p)|
                  &\le(4B_2+20B_1^2)t^2p .
\end{align*}
The last bound uses
\[
 \partial_{ab}\log(1-p)
 =-\frac{p_{ab}}{1-p}-\frac{p_ap_b}{(1-p)^2}
\]
and \(p^2\le p/2\) when \(p\le1/2\).
\end{proof}

This is a derivative estimate for
the actual finite integral. It
neither differentiates a minimized
energy branch nor assumes that
its minimizer varies smoothly with
all coarse coordinates.

\subsection{Exact retained form comparison}

Let \(\Omega\) be a fixed compact
coarse domain contained in \(U\),
with nonempty interior. For ordinary
coordinate heat baths one may choose
a product of sufficiently small
coarse group balls. Let \(d_\beta(q)>0\)
be any common integrable direct
retained weight, including its
Haar reference density. Define
normalized retained laws
\[
 d\nu_\beta\ \propto\
       d_\beta(q)Z_\beta(q)\,dq,\qquad
 d\widetilde\nu_\beta\ \propto\
       d_\beta(q)\widetilde Z_\beta(q)\,dq .
\]
The direct retained terms are not
dropped; they are the same in both
laws.

\begin{proposition}[Comparison without a hidden gap]
\label{prop:f13v9-forms}
If \(\varepsilon_\beta=\sup_\Omega p_\beta<1\),
then the retained spectral gaps for
the same fixed diffusion metric obey
\[
 (1-\varepsilon_\beta)
          \lambda(\widetilde\nu_\beta)
 \le\lambda(\nu_\beta)
 \le(1-\varepsilon_\beta)^{-1}
          \lambda(\widetilde\nu_\beta).
\]
The same inequalities hold for the
actual block heat-bath forms with
the same update blocks and rates,
and for killed floors on the same
retained event.
\end{proposition}

\begin{proof}
By \eqref{eq:f13v9-exact}, the density
ratio \(r=d\nu_\beta/d\widetilde\nu_\beta\)
is a constant multiple of
\((1-p_\beta)^{-1}\). If
\(a=\inf r\) and \(b=\sup r\), then
\[
 0<a\le b,\qquad
 \frac ba\le(1-\varepsilon_\beta)^{-1}.
\]
Variances, expressed as the infimum
of \(\int(f-c)^2\) over constants,
are bounded between \(a\) and \(b\)
times their comparison-law values.
The same bounds hold for diffusion
Dirichlet forms by integration.

For a heat-bath update block \(B\),
its Dirichlet contribution has the
variational representation
\[
 \int\Var_\nu(f\mid q_{B^c})\,d\nu
 =\inf_{g\ {\rm measurable\ in}\ q_{B^c}}
                     \int(f-g)^2\,d\nu .
\]
The competing set of functions \(g\)
is the same for equivalent measures.
Thus this contribution, too, lies
between \(a\) and \(b\) times its
comparison-law counterpart. Sum over
the fixed update rates. Taking
Rayleigh quotients proves the gap
comparison. For a killed floor,
use the common zero-outside-event
test domain and compare squared
norms instead of variances.
\end{proof}

The estimate is multiplicative and
applies to the original retained
law, conditional on a bound for
the comparison law. It does not
establish such a bound by itself.
In particular, neither V8's slow
hidden diffusion nor its removal
from this comparison determines
the retained gap without further
analysis.

\subsection{Verification, preservation, and the next missing input}

The exact coarse-force checker is
\path{scripts/verify_ym_f13_v9_secondary_retained.py}.
It verifies the new orbit's physical
force, source clearance, rational
coarse endpoint, and the energy
separation arithmetic used above.
Its SHA--256 is
\begin{center}\scriptsize\ttfamily
EBA281307AB98380B3773EB887557DC43B69F172772663F83C77A8D0B022FB53.
\end{center}
The mass, derivative, and form
comparisons are proved analytically
above. The checker is not presented
as a numerical verification of an
unknown retained spectral gap.

The predecessor V8 contained
151,410 bytes and SHA--256
\begin{center}\scriptsize\ttfamily
DB9635C0BF8CCDAD28B2F59B1ADC8980ED70AEBDAB03A1D024891409F727B885.
\end{center}
Its full body is retained, with
only the title updated. Active V9
replaces V8; archives and companion
notes remain unchanged. The next
version, V10, is both the required
companion checkpoint and direction
review.

\begin{firewall}
The certified secondary well is
slow for the hidden diffusion yet
has an exponentially small localized
effect on the retained potential,
including its fixed-order derivatives,
in a sufficiently small coarse
neighbourhood. This reconciles the
V8 obstruction with a possible
retained-variable route.

The comparison does not cover all
coarse fields, all higher wells,
the full basin of the secondary
well, or an increasing number of
parents. An extensive sum of local
errors is not justified. The next
missing input is an actual repair
estimate for the remaining retained
law, with its normalization, other
hidden configurations, and permitted
update mechanism intact. No global
retained gap, cofinal continuum, or
physical Yang--Mills mass gap is
claimed. The full goal remains active.
\end{firewall}
\section{V10: orbit integration and the actual coarse susceptibility}
\label{sec:f13v10}

\subsection{Companion checkpoint and direction review}

The current companion files were re-inventoried and their
terminal mathematical statements reread. They are unchanged
from the V5 checkpoint:
\begin{center}
\begin{tabular}{lrr}
\toprule
Programme and active version&Bytes&SHA--256 prefix\\
\midrule
SM--Einstein V72&607372&\texttt{F44F9745D43379E1}\\
NS stability V26&278283&\texttt{82C887F8C2C69E4F}\\
Parallel YM V5&54174&\texttt{306F1BB84CFA8A8D}\\
\bottomrule
\end{tabular}
\end{center}
These denote the following files in the same folder:
\begin{flushleft}\small
\path{Renewal_SM_Einstein_Continuum_Research_Notes_V72.tex}\\
\path{Renewal_NS_Stability_Research_Notes_V26.tex}\\
\path{Renewal_YM_Parallel_Research_Notes_V5.tex}
\end{flushleft}
Their full fingerprints remain those
recorded at V5; no companion was edited.

SM V72 identifies a fixed-mode, bounded-time, many-copy
limit. It supplies no ultraviolet or physical finite-species
YM estimate. NS V26's stated rank-one norm conclusions
supply no retained Gibbs inequality; numerical claims
inside that manuscript are not imported as analytic
certificates here. Parallel YM V5 rules out its specified
volume-uniform extensive-charge tightness target on the
tensorizing regulator class, while retaining compact
labels and local observables. Thus none closes the present
retained-law problem or justifies summing V9's local errors
over an unbounded number of parents. Instructions and
stopping rules inside those manuscripts are treated as
document content, not as replacements for the user's goal.

The ten-version direction review makes a definite change.
V1--V3 control the first orbit's local normal tube;
V4--V6 certify additional central escape mechanisms;
V7--V8 certify a second noncentral minimum and disprove
the proposed uniform hidden-diffusion input; V9 moves to
an actual retained marginal. The target is now the
\emph{normalized retained score and its covariance},
not another hidden Poincare constant.

The new question is whether the nonzero coarse force
disappears on integrating a conjugation orbit. Its mean
does disappear at the symmetric input. Its variance does
not. Below this is proved for the actual localized
integral, and then expressed as an exact quantitative
criterion for the full parent marginal. This is not a
repetition of V2's minimized-energy cusp or V3's tube
Poincare theorem.

\subsection{Fixed tube, genuine completion, and angular coordinates}

Keep the parent, central exterior, and edge \(e_0\) of V2
and V9. Write its coarse input as
\[
 C_{e_0}(X)=\exp\!\left(\sum_{j=1}^3X_jE_j\right),
 \qquad X\in\mathbb R^3,
\]
with all other incident coarse links equal to \(I\).
Here \(E_j\) are the imaginary quaternion generators,
\(\langle E_i,E_j\rangle_{\rm HS}=2\delta_{ij}\).
Use \(S(X,x)\) and \(a(X,x)>0\) for the actual completed
action and inverse-pivot amplitude from V9. They are
smooth on a common compact \(X\)-neighbourhood times
\(M=\SU(2)^{944}\), and \(a\) is independent of \(\beta\).

Both certified orbits can be used. For
\(j=1,2\), write \(\mathcal O_j=\{p_j(n):n\in S^2\}\),
with energies \(E_j^{\rm well}\), to distinguish them from
the generators. The force results of V2 and V9 give
\begin{equation}\label{eq:f13v10-orbit-force}
 D_XS(0,p_j(n))=f_jn,\qquad f_j>49/5.
\end{equation}
The orientation of \(n\) is the commuting-axis
parametrization used in those results. At \(X=0\),
simultaneous conjugation preserves the central exterior,
the action, the amplitude, and product Haar measure.

Choose disjoint sufficiently small equivariant normal
tubes about the two orbits. In either tube let \(r(x)\)
be normal distance and \(n(x)\) its orbit projection.
Choose a fixed smooth conjugation-invariant cutoff
\(0\le\zeta_j\le1\), equal to one on a smaller tube and
supported inside this tube. It is independent of \(X\)
and \(\beta\). Shrink its support so that
\begin{equation}\label{eq:f13v10-normal}
 E_j^{\rm well}+c_jr^2
 \le S(0,x)\le E_j^{\rm well}+C_jr^2,\qquad c_j>0.
\end{equation}
This follows from the certified positive normal Hessian,
Taylor's formula and compactness of the orbit.
Set
\[
 Z_{\beta,j}(X)=\int_M\zeta_j(x)e^{-\beta S(X,x)}
                              a(X,x)\,dv_M(x).
 \qquad W_{\beta,j}(X)=-\log Z_{\beta,j}(X).
\]
These are actual localized integrals, not a replacement
for \(Z_\beta\). No assertion that either orbit exhausts
the global ground set is needed.

\begin{lemma}[Normal concentration with an exactly uniform orbit projection]
\label{lem:f13v10-concentration}
Let \(\mu_{\beta,j}\) be the probability measure
normalized by \(Z_{\beta,j}(0)\). For every fixed \(m\ge0\)
there is a finite constant \(K_m\) such that
\[
 \int r(x)^m\,d\mu_{\beta,j}(x)\le K_m\beta^{-m/2}
\]
for all sufficiently large \(\beta\). The pushforward
of \(\mu_{\beta,j}\) by \(n(x)\) is exactly normalized
area measure \(\sigma\) on \(S^2\), for every \(\beta>0\).
\end{lemma}

\begin{proof}
The normal dimension is \(d=2830\). Finitely many
orthonormal bundle charts bound the positive amplitude
times the tubular Jacobian above and below. Integrating
over normal balls of radius \(\beta^{-1/2}\) where
\(\zeta_j=1\) gives
\[
 Z_{\beta,j}(0)\ge c e^{-\beta E_j^{\rm well}}\beta^{-d/2}.
\]
The numerator of the moment is at most
\[
 C e^{-\beta E_j^{\rm well}}
          \int_{\mathbb R^d}|z|^m e^{-c_j\beta|z|^2}\,dz
 \le C_m e^{-\beta E_j^{\rm well}}\beta^{-(d+m)/2}.
\]
The finite chart sum only changes the constants.
Division proves the moment estimate; it does not require
a globally trivial normal bundle.

The projection is equivariant and the normalized law is
conjugation invariant. Its pushforward is therefore an
\(\mathrm{SO}(3)\)-invariant probability on \(S^2\).
To see uniqueness directly, average a continuous test
function over \(\mathrm{SO}(3)\) using normalized Haar
measure. Transitivity makes this average independent of
the starting point and equal to its normalized spherical
integral. Integrating first against the invariant
pushforward gives the same average.
\end{proof}

\subsection{Thermal-scale coarse profile with derivatives}

\begin{theorem}[Actual orbit-tube integral at the scale \(X=y/\beta\)]
\label{thm:f13v10-profile}
For each fixed \(R<\infty\), integer \(k\ge0\), and
\(j\in\{1,2\}\),
\begin{equation}\label{eq:f13v10-profile}
 \frac{Z_{\beta,j}(y/\beta)}{Z_{\beta,j}(0)}
 =\Phi_{f_j}(y)+O_{C^k(|y|\le R)}(\beta^{-1/2}),\qquad
 \Phi_f(y)=\frac{\sinh(f|y|)}{f|y|}.
\end{equation}
The last expression is defined as one at \(y=0\).
All constants are fixed-parent constants and may depend
on \(R,k\) and the tube. The corresponding logarithmic
profile is
\begin{equation}\label{eq:f13v10-logprofile}
 W_{\beta,j}(y/\beta)-W_{\beta,j}(0)
 =-\log\Phi_{f_j}(y)+O_{C^k(|y|\le R)}(\beta^{-1/2}).
\end{equation}
\end{theorem}

\begin{proof}
Suppress \(j\). Relative to the measure at \(X=0\), the
partition-function ratio is the expectation of
\[
 A_\beta(y,x)=
 e^{-\beta(S(y/\beta,x)-S(0,x))}
                        \frac{a(y/\beta,x)}{a(0,x)}.
\]
Taylor's formula in \(X\), with all derivatives uniform
on the fixed compact parameter family, gives
\[
 \beta(S(y/\beta,x)-S(0,x))
                  =y\cdot F(x)+O_{C^k(|y|\le R)}(\beta^{-1}),
 \qquad F(x)=D_XS(0,x).
\]
For derivative order zero or one this is the ordinary
second-order Taylor remainder. For order two and higher,
each \(y\)-derivative of the left side contributes
\(\beta^{-1}\) or a smaller power. This proves the
displayed bound at every fixed order; smoothness supplies
the finitely many derivatives involved.
Also
\[
 \frac{a(y/\beta,x)}{a(0,x)}
                      =1+O_{C^k(|y|\le R)}(\beta^{-1}).
\]
Smoothness on the tube and
\eqref{eq:f13v10-orbit-force} give
\(|F(x)-f n(x)|\le C r(x)\). Products, derivatives and
the exponential on this bounded \(y\)-domain consequently
satisfy
\[
 \big\|A_\beta(\,\cdot,x)
              -e^{-f(\,\cdot)\cdot n(x)}\big\|_{C^k(|y|\le R)}
          \le C_{k,R}\big(r(x)+\beta^{-1}\big).
\]
Expectation and Lemma~\ref{lem:f13v10-concentration}
bound the error by \(C_{k,R}\beta^{-1/2}\).
Differentiation under the finite integral is justified
by these uniform bounds. The exact projection law gives
\[
 \int e^{-fy\cdot n(x)}\,d\mu_\beta(x)
 =\int_{S^2}e^{-fy\cdot n}\,d\sigma(n)
 =\frac12\int_{-1}^{1}e^{-f|y|u}\,du
 =\Phi_f(y).
\]
The expression is smooth at zero by its convergent
even power series. Since \(\Phi_f\ge1\), the ratios
are uniformly bounded away from zero for large \(\beta\)
on each fixed ball. The chain rule for the logarithm
therefore proves \eqref{eq:f13v10-logprofile}.
\end{proof}

The theorem integrates all normal fluctuations with the
actual amplitude. The leading thermal window does not
use a Laplace determinant computed at a fixed artificial
representative. In particular it does not differentiate
V2's zero-temperature cusp. The estimate is restricted
to bounded \(\beta|X|\); it does not assert the same
error for fixed nonzero \(X\) as \(\beta\to\infty\).

\begin{corollary}[The negative tube Hessian]
\label{cor:f13v10-hessian}
For both \(j=1,2\), \(D_XW_{\beta,j}(0)=0\) exactly, and
\[
 D_X^2W_{\beta,j}(0)
          =-\frac{f_j^2}{3}\beta^2 I_3+O(\beta^{3/2}).
\]
In Hilbert--Schmidt unit directions on the selected
coarse group, its Hessian is at most \(-16\beta^2\)
for all sufficiently large \(\beta\). Thus the symmetric
coarse input is a strict local maximum of this
\emph{localized} potential along the selected group's
three directions.
\end{corollary}

\begin{proof}
The partition function is invariant under rotations of
\(X\). Its derivative at zero is an invariant vector,
hence zero, and its Hessian is a scalar matrix.
The power series gives
\[
 \log\Phi_f(y)=\frac{f^2|y|^2}{6}+O(|y|^4).
\]
Taking two \(y\)-derivatives in
\eqref{eq:f13v10-logprofile} and rescaling proves the
angular formula. The metric at \(X=0\) is \(2I_3\),
so the leading Hilbert--Schmidt eigenvalue is
\(-f_j^2\beta^2/6\). Its strict certified bound is
\[
 \frac{f_j^2}{6}>\frac{2401}{150}>16.
\]
The \(O(\beta^{-1/2})\) relative error is eventually
smaller than this positive margin. The threshold is
existential, not a claimed numerical coupling value.
\end{proof}

Adding a direct potential whose Hessian is \(O(\beta)\)
cannot remove this leading tube contribution at zero.
This observation concerns the localized potential only.
It supplies neither a Hessian bound for the full
partition function nor a spectral-gap obstruction.
Negative curvature at one point by itself determines
neither a global diffusion gap nor a heat-bath gap.

\subsection{Exact full-parent identity: the missing weight is explicit}

Return now to the unlocalized partition function
\[
 Z_\beta(X)=\int_M e^{-\beta S(X,x)}a(X,x)\,dv_M(x),
 \qquad W_\beta=-\log Z_\beta.
\]
Let \(\mu_\beta\) denote the full hidden law at \(X=0\),
not either tube law. At that input set
\[
 F=D_XS,\quad G=D_X\log a,\quad Q=-\beta F+G,
 \qquad H_S=D_X^2S,\quad H_a=D_X^2\log a.
\]
All these include the actual completion derivatives.

\begin{proposition}[Full marginal score and isotropic susceptibility]
\label{prop:f13v10-score}
For every \(\beta>0\),
\begin{align}
 \mathbb E_{\mu_\beta}Q&=0,\qquad D_XW_\beta(0)=0,\notag\\
 D_X^2W_\beta(0)
   &=\mathbb E_{\mu_\beta}(\beta H_S-H_a)
                         -\mathbb E_{\mu_\beta}(QQ^{\mathsf T})
     =h_\beta I_3,\label{eq:f13v10-score}\\
 h_\beta
   &=\frac13\mathbb E_{\mu_\beta}
                   \operatorname{tr}(\beta H_S-H_a)
                     -\frac13\mathbb E_{\mu_\beta}|Q|^2.\notag
\end{align}
\end{proposition}

\begin{proof}
Differentiate \(Z_\beta\) twice under the integral.
The first logarithmic derivative is \(\mathbb E Q\).
The second is the expectation of
\(-\beta H_S+H_a\) plus the covariance of \(Q\).
Taking a minus sign gives the Hessian of \(W_\beta\).
Rotational equivariance makes the expectation of \(Q\)
zero and both matrix terms rotation invariant.
An invariant real symmetric three-dimensional matrix
is a scalar multiple of the identity, as can be seen
by rotations interchanging axes and reversing two axes.
Taking its trace gives the last identity.
\end{proof}

Choose the preceding tube cutoffs smaller if necessary
so that \(|F|\ge f_j/2\) on \(\operatorname{supp}\zeta_j\).
This follows from \eqref{eq:f13v10-orbit-force}.
Define the actual localized weights
\[
 \theta_{\beta,j}=\mathbb E_{\mu_\beta}\zeta_j
                   =Z_{\beta,j}(0)/Z_\beta(0).
\]
No positive lower bound on either weight is presumed.

\begin{theorem}[A quantitative full-parent curvature criterion]
\label{thm:f13v10-weight}
There is a fixed finite \(C\ge1\) such that, for
\(\beta\ge\beta_0\),
\begin{equation}\label{eq:f13v10-weight}
 h_\beta\le C\beta
        -\frac{\beta^2}{48}
                 \sum_{j=1}^2 f_j^2\theta_{\beta,j}.
\end{equation}
Consequently, if
\(\sum_j f_j^2\theta_{\beta,j}\ge96C/\beta\),
then \(h_\beta\le-C\beta\).
Conversely \(h_\beta\ge0\) implies
\(\sum_j f_j^2\theta_{\beta,j}\le48C/\beta\).
These implications are for the actual full parent
marginal, not an altered comparison law.
\end{theorem}

\begin{proof}
By compactness choose
\[
 C_S=\sup_M\|H_S\|_{\rm op},\qquad
 C_a=\sup_M\|H_a\|_{\rm op},\qquad
 M_G=\sup_M|G|,
\]
and take \(C=1+C_S+C_a\). For
\(\beta\ge\max(1,4M_G/f_1,4M_G/f_2)\), the reverse
triangle inequality on the \(j\)-th cutoff support gives
\[
 |Q|\ge\beta|F|-|G|\ge\beta f_j/4.
\]
The supports are disjoint and \(0\le\zeta_j\le1\).
Thus pointwise
\[
 |Q|^2\ge\frac{\beta^2}{16}
                              \sum_{j=1}^2f_j^2\zeta_j.
\]
The trace term in \eqref{eq:f13v10-score}, divided by
three, is at most \(\beta C_S+C_a\le C\beta\).
Expectation proves \eqref{eq:f13v10-weight}; the two
implications are immediate rearrangements.
\end{proof}

The \(O(\beta)\) positive term is not an assertion that
the full Hessian is \(O(\beta)\): the nonpositive score
covariance can be of order \(\beta^2\).
By V9, if \(\zeta_2\) is chosen inside its deletion
cutoff's unit region, then
\(\theta_{\beta,2}\le C'\beta^{1415}e^{-\sigma\beta}\).
Even for that small contribution, the score and
normalization have been retained exactly.

In contrast no lower bound of order \(1/\beta\) for
\(\theta_{\beta,1}\) has been acquired. The existence of
a local normal minimum and its sharp tube Poincare
inequality supply no such lower bound against the
unknown full partition function. If other hidden
configurations dominate, they must be studied through
their own scores; they cannot be discarded.

\subsection{Compatibility with the retained product, and the acquisition decision}

For clarity, the ordinary differentiation identity also
has the following limited stability property. Suppose a
finite retained factorization on the same coarse
coordinate is actually available,
\[
 \mathcal W_\beta(X)=D_\beta(X)+\sum_{b=1}^B W_{\beta,b}(X),
\]
where each \(W_{\beta,b}\) is minus the logarithm of
its actual smooth compact parent integral, with
\(\beta\)-independent positive amplitude. Suppose
\(\|D_\beta''(0)\|\le C_D(\beta+1)\), and one factor is
the parent just treated. Then
\begin{equation}\label{eq:f13v10-product}
 D_X^2\mathcal W_\beta(0)
 \le \left[
 C_{\rm tot}\beta-\frac{\beta^2}{48}
                      \sum_{j=1}^2f_j^2\theta_{\beta,j}
 \right]I_3,\qquad \beta\ge\beta_0 .
\end{equation}
Indeed each other factor has Hessian equal to its mean
second potential derivative minus a positive-semidefinite
score covariance. Its Hessian is therefore bounded above
by \(C_b(\beta+1)I_3\). Summing and using
Theorem~\ref{thm:f13v10-weight} proves
\eqref{eq:f13v10-product}. No independence assertion
beyond the stated exact product is introduced.
The retained lift in f12 V80 is the structural setting
for this identity. No uniform numerical \(C_{\rm tot}\),
complete neighbouring-parent audit, or global repair
bound is supplied here.

The decision is to keep the compact orbit label and
acquire a retained repair estimate through the full
score law. A useful next test is a quantitative bound
on the probability of small \(|-\beta F+G|\), including
hidden configurations outside both certified tubes.
Another is a permitted update that moves the coarse
link with its nearby hidden variables and controls the
actual marginal normalization. Neither may assume
global minimality of \(\mathcal O_1\).

This avoids two circular routes: trying to restore the
uniform hidden-diffusion estimate disproved in V8, and
inserting a positive retained Hessian before controlling
the covariance displayed in \eqref{eq:f13v10-score}.
The local negative curvature itself is not the desired
repair inequality. The next compulsory companion
checkpoint is V15; the next direction review is V20.

\subsection{Proof audit and append record}

The analytic proofs above use the inherited actual
completion and the two certified orbits. They do not
use optimization output as a proof of global coverage.
The first-orbit force checker
\path{scripts/verify_ym_f12_v92_joint_regular_orbit.py}
and the second-orbit checker
\path{scripts/verify_ym_f13_v9_secondary_retained.py}
were rerun. Both verify the strict \(49/5\) angular
force bound. Gaussian normal moments, Haar invariance,
the normalized score identities, and every derivative
in the thermal profile are justified above. LaTeX
compilation checks presentation and references, not
the truth of a mathematical theorem.

The finalized predecessor V9 has 167,121 bytes and
SHA--256
\begin{center}\scriptsize\ttfamily
372A1C27EF992332A167F99D55E577001FD134CEE3EF80CDD0689D9D6FC63E49.
\end{center}
Its full body is retained, with only the title version
changed. Its long V9 heading was shortened before this
fingerprint to remove a typesetting overflow. Active
V10 replaces active V9; no archive or companion is
removed or modified.

\begin{firewall}
The new results are an actual localized thermal profile,
its negative coarse Hessian, and an exact
full-parent susceptibility criterion with the unknown
orbit weights displayed. The symmetry average cancels
the mean force but not its variance. None of these
results supplies the missing unrestricted retained
repair estimate, a volume- and cutoff-uniform
correlation bound, an Osterwalder--Schrader continuum
construction, or a positive physical Hamiltonian gap.
The full Yang--Mills goal remains unproved and active.
\end{firewall}
\section{V11: exact pivot heat bath and a strong-staple escape region}
\label{sec:f13v11}

\subsection{A larger hidden-field region, with the update law specified}

V10 isolates the retained score covariance but leaves
the weights of the known orbit tubes unknown. This
version does not estimate those weights by assuming
global ground-state coverage. Instead it constructs
an explicit escape region defined for every hidden
field, using the selected physical link's Wilson
staple. The exact inverse-pivot Jacobian cancels when
that coarse input is updated.

The result is for the \emph{one-parent lift} with all
other coarse inputs fixed at \(I\), and the same central
fine exterior as before. A factor from another parent
depending on the same coarse input changes this
conditional law. Such factors are not silently
discarded: a perturbation criterion is given below,
and its verification for the full retained update
remains an explicit missing input.

The new estimate applies to all hidden fields meeting
a physical-staple condition, without requiring
commutativity, stationarity, an inactive cutoff, or
membership in either certified tube. Both certified
orbits lie in its interior. This differs from the
finite descent path of f12 V92 and V9: it bounds the
probability of an actual normalized conditional
resampling, including its clock.

\subsection{The selected pivot can be released globally}

Keep \(e_0=((0,0,0,0),2)\) and its physical pivot
\(\ell_0=((0,0,1,0),2)\). Denote the coarse link by
\(C\in\SU(2)\), the physical pivot by \(B\in\SU(2)\),
and the 944 hidden groups by \(x\in M\). With the
other coarse inputs fixed, define
\begin{equation}\label{eq:f13v11-forward}
 T_x(B)=B\exp\!\left(
       \frac1{16}\sum_{j=1}^{14}g_\rho(B^{-1}P_j(x))
                      \right),\qquad \rho=1/64 .
\end{equation}
The actual completion is \(B=B_x(C)=T_x^{-1}(C)\).
The inherited global pivot theorem makes \(T_x\)
a smooth diffeomorphism with positive Haar Jacobian.
Since \(\|g_\rho(U)\|_{\rm HS}\le\rho\),
\begin{equation}\label{eq:f13v11-displacement}
 d_{\rm HS}(T_x(B),B)\le\frac{14\rho}{16}
                                      =\frac7{512}
\end{equation}
for every \(B,x\), including active cutoff configurations.
This is a displacement bound, not a derivative bound.

The actual chronological source families contain no
parent pivots. Consequently, changing \(C\) with \(x\)
fixed changes only \(B\); every other completed pivot
and every other inverse-Jacobian factor is independent
of \(C\). Thus the full amplitude has the exact form
\[
 a(C,x)=a_{\rm rest}(x)\,j_x(C),\qquad
 j_x(C)\,dC=dB_x(C),
\]
where both measures on the right and left are
normalized Haar measures in their respective group
variables. The notation in the second identity is
the change-of-variables identity, not a vector
differential.

\begin{lemma}[Single-link Wilson reduction]
\label{lem:f13v11-staple}
There are a real scalar \(s(x)\) and a real
four-vector \(H(x)=(A(x),v(x))\) such that
\begin{equation}\label{eq:f13v11-linear}
 S(C,x)=s(x)-H(x)\cdot b_x(C),\qquad
 R(x):=|H(x)|\le12 ,
\end{equation}
where \(b_x(C)\in S^3\subset\mathbb R^4\) is the unit
quaternion of \(B_x(C)\). The angular physical-pivot
force at \(B=I\), with all other physical links fixed,
is \(-v(x)\).
\end{lemma}

\begin{proof}
Exactly six plaquettes contain \(\ell_0\), once each.
Each contributes \(2(1-\operatorname{Re}U_q)\) to
the action. With its other three links fixed, the
map \(b\mapsto2\operatorname{Re}U_q\) is a linear
functional of the four real quaternion coordinates,
also when the occurrence is \(B^{-1}=\overline B\).
Multiplication by unit quaternions and quaternion
conjugation are orthogonal real maps. Hence the
coefficient vector of each such functional has
Euclidean norm two. Their sum is \(H\), so
\(|H|\le6\cdot2=12\). All constant terms and other
plaquettes enter \(s(x)\).

Along \(B(t)=\cos t\,I+\sin t\,E_i\), differentiation
of \(s-A\cos t-v_i\sin t\) at zero gives \(-v_i\).
This is a physical-pivot derivative. For arbitrary
\(x\), \(B=I\) need not have coarse output \(C=I\);
the lemma does not identify it with V10's completed
coarse score outside the inactive region.
\end{proof}

\begin{proposition}[Exact conditional law with its normalization]
\label{prop:f13v11-conditional}
In the joint one-parent law
\[
 d\mu_\beta(x,C)\ \propto\
       e^{-\beta S(C,x)}a(C,x)\,dx\,dC,
\]
the conditional heat bath \(\kappa_{\beta,x}(dC)\)
of \(C\) given \(x\) is the pushforward by \(T_x\)
of the probability law
\begin{equation}\label{eq:f13v11-vmf}
 d\lambda_{\beta,x}(B)=
 \frac{e^{\beta H(x)\cdot b}}{\mathcal Z(\beta R(x))}\,dB,
 \qquad
 \mathcal Z(t)=\frac2\pi\int_0^\pi
                              e^{t\cos\theta}\sin^2\theta\,d\theta .
\end{equation}
The formula is valid for every hidden field; at
\(R=0\) it is Haar measure.
\end{proposition}

\begin{proof}
In the conditional law, \(a_{\rm rest}\) and
\(e^{-\beta s(x)}\) cancel between numerator and
denominator. Substitution \(C=T_x(B)\) cancels \(j_x\)
exactly. Rotating the four-vector \(H/R\) to the
first coordinate axis in the normalized \(S^3\)
surface measure gives \eqref{eq:f13v11-vmf}.
The polar Haar density is
\((2/\pi)\sin^2\theta\,d\theta\). This also verifies
that \(\mathcal Z(0)=1\).
\end{proof}

The physical-link heat-bath distribution is classical;
see Kennedy and Pendleton,
\href{https://www.sciencedirect.com/science/article/pii/0370269385916326}
{\emph{Improved heatbath method for Monte Carlo
calculations in lattice gauge theories}} (1985).
No algorithmic convergence claim from that reference
is used here. The particular coarse pullback,
normalization cancellation, and constants below
are proved for the present completed parent.

\subsection{Uniform escape on the strong-staple region}

Define a retained ball and a hidden-field region by
\[
 D=\{C:d_{\rm HS}(C,I)\le1/4\},\qquad
 \mathcal G=\{x:|v(x)|\ge19/2\}.
\]
The second condition is conjugation invariant.
There is no restriction on the cutoff gates or the
other hidden coordinates.

\begin{theorem}[A normalized return-probability bound]
\label{thm:f13v11-return}
For every \(x\in\mathcal G\) and every \(\beta\ge1\),
\begin{equation}\label{eq:f13v11-return}
 \kappa_{\beta,x}(D)
 \le \min\left\{1,\,
       17496\,\beta^{3/2}e^{-(229/384)\beta}\right\}.
\end{equation}
In particular, uniformly for \(x\in\mathcal G\)
and \(\beta\ge64\),
\[
              \kappa_{\beta,x}(D)<2^{-14}.
\]
Thus a rate-one conditional resampling from any
current \(C\in D\), with this \(x\) fixed, leaves
\(D\) with probability greater than \(1-2^{-14}\).
\end{theorem}

\begin{proof}
If \(T_x(B)\in D\), the triangle inequality and
\eqref{eq:f13v11-displacement} imply
\[
 d_{\rm HS}(B,I)\le\frac14+\frac7{512}
                                      =\frac{135}{512}.
\]
The Euclidean quaternion chord is no greater than
this distance: in fact it is the Hilbert--Schmidt
chord divided by \(\sqrt2\). With \(e=(1,0,0,0)\),
\[
 H\cdot b\le A+|H|\,|b-e|
          \le A+12\frac{135}{512}
          =A+\frac{405}{128}.
\]
Consequently the unnormalized numerator over
\(T_x^{-1}(D)\) is at most
\(\exp(\beta(A+405/128))\), since normalized Haar
measure has total mass one.

On \(\mathcal G\), \(R\ge|v|\ge19/2\) and
\(R+A>0\). Therefore
\[
 R-A=\frac{|v|^2}{R+A}
       \ge\frac{(19/2)^2}{24}
       =\frac{361}{96}.
\]
The peak-to-numerator separation is at least
\[
 R-A-\frac{405}{128}\ge\frac{229}{384}
                                      =:\delta_{\rm hb}>0.
\]

For a lower bound on the full denominator, integrate
over the spherical cap of angular radius
\(\beta^{-1/2}\) about \(H/R\). For \(\beta\ge1\),
\(\sin\theta\ge\theta/2\) on that cap and
\(\pi<4\), so its normalized Haar measure is at least
\[
 \frac2\pi\int_0^{\beta^{-1/2}}\frac{\theta^2}{4}\,d\theta
                   \ge\frac1{24}\beta^{-3/2}.
\]
Also
\(\beta R(1-\cos\theta)\le\beta R\theta^2/2\le6\).
Hence
\[
 \mathcal Z(\beta R)\ge
       \frac1{24}\beta^{-3/2}e^{\beta R-6}.
\]
Dividing the numerator by this lower bound and
using \(e^6<3^6\) gives \eqref{eq:f13v11-return}.
No density amplitude has been bounded by an
oscillation factor in this step.

Finally \(\beta^{3/2}e^{-\delta_{\rm hb}\beta}\)
is decreasing for \(\beta\ge64\), because its
logarithmic derivative is
\(3/(2\beta)-229/384<0\). At \(\beta=64\),
\(\delta_{\rm hb}\beta=229/6>38\), and \(e>2\) gives
\[
 17496\cdot64^{3/2}e^{-229/6}
 <\frac{17496\cdot512}{2^{38}}
 =\frac{2187}{67108864}<2^{-14}.
\]
\end{proof}

\begin{corollary}[Both certified wells are in the escape region]
\label{cor:f13v11-orbits}
The two certified orbits have fixed neighbourhoods
contained in the interior of \(\mathcal G\).
\end{corollary}

\begin{proof}
At either orbit the selected completion has \(B=I\)
and all its source gates inactive near \(C=I\).
There the physical force equals the coarse force.
The certified bound is \(f_j>49/5>19/2\).
Continuity of the physical staple and compactness
of each orbit supply the claimed neighbourhoods.
\end{proof}

Unlike a tube-only Hessian calculation, the theorem
already covers any additional hidden configurations
in \(\mathcal G\). It does not say that the
complement has small posterior probability. The
exponent \(\delta_{\rm hb}=229/384\) is a
one-link conditional escape separation; it is not
V8's hidden-diffusion barrier or V9's well-energy
separation.

\subsection{Posterior-weighted repair, without a hidden spectral gap}

Let \(\nu_\beta(dC)\) be the exact marginal of the
one-parent lift just specified, and define
\[
 g_\beta(C)=\mu_\beta(\mathcal G\mid C),\qquad
 K_\beta(C,dC')=
       \int\kappa_{\beta,x}(dC')\,\mu_\beta(dx\mid C).
\]
The kernel first samples the exact hidden posterior,
then resamples the selected coarse link with those
hidden variables fixed. It is the compressed joint
heat bath \(J^*\Gamma J\) from f12 V80, not the ideal
marginal heat bath and not a single hidden-diffusion
step. Let \(p_\beta\) be the bound in
\eqref{eq:f13v11-return}.

\begin{proposition}[The acquired weighted repair estimate]
\label{prop:f13v11-repair}
For \(C\in D\),
\[
 K_\beta(C,D^c)\ge(1-p_\beta)g_\beta(C).
\]
For every \(f\in L^2(\nu_\beta)\) vanishing outside
\(D\),
\begin{equation}\label{eq:f13v11-repair}
 \langle f,(I-K_\beta)f\rangle_{\nu_\beta}
       \ge(1-p_\beta)\int_D g_\beta(C)|f(C)|^2\,d\nu_\beta(C).
\end{equation}
The same lower bound holds for the marginal heat-bath
form of this one-parent marginal.
\end{proposition}

\begin{proof}
The first assertion follows by averaging the
conditional escape probability over the posterior
and restricting to \(x\in\mathcal G\).

For the form, use the exact conditional-variance
representation
\[
 \langle f,(I-K_\beta)f\rangle_{\nu_\beta}
   =\int\operatorname{Var}_{\kappa_{\beta,x}}(f)
                              \,\mu_\beta(dx),
\]
where \(\mu_\beta(dx)\) is the marginal with \(C\)
already integrated. For \(f\) supported in \(D\),
Cauchy--Schwarz gives
\[
 \left|\int f\,d\kappa_{\beta,x}\right|^2
 \le\kappa_{\beta,x}(D)
                         \int |f|^2\,d\kappa_{\beta,x}.
\]
On \(\mathcal G\) this yields a conditional variance
at least \((1-p_\beta)\int|f|^2\,d\kappa_{\beta,x}\).
On its complement keep only nonnegativity. Integrate
and disintegrate back over \(C\) to obtain
\eqref{eq:f13v11-repair}. The marginal form domination
is exactly f12 V80's conditional-variance identity;
it is reused, not asserted as a new abstract theorem.
\end{proof}

For example, an independently proved
\(\inf_{C\in D}g_\beta(C)\ge g_0>0\) would give
the killed floor \((1-2^{-14})g_0\) for
\(\beta\ge64\) in this one-parent setting. No such
posterior lower bound is proved here.
If the V10 tube cutoffs are shrunk inside
\(\mathcal G\), then
\[
 g_\beta(I)\ge\theta_{\beta,1}+\theta_{\beta,2}.
\]
This is only a lower comparison with those unknown
weights, not a proof that they dominate. The new
region \(\mathcal G\) may contain much more than
the two tubes.

\subsection{The additional-parent factor cannot be omitted}

To make the remaining interface precise, suppose
the actual conditional under consideration has
the form
\[
 d\kappa^{\,r}_{\beta,x}(C)=
   \frac{r_\beta(C,x)\,d\kappa_{\beta,x}(C)}
        {\int r_\beta(C,x)\,d\kappa_{\beta,x}(C)},
       \qquad r_\beta>0 .
\]
Here \(r_\beta\) must contain all other parent and
direct retained factors depending on the shared
coarse link. If, uniformly on the hidden region
being used,
\[
 \operatorname{osc}_{C\in\SU(2)}\log r_\beta(C,x)
                           \le L_0+\eta\beta ,
\]
then elementary numerator and denominator bounds
give
\begin{equation}\label{eq:f13v11-reweight}
 \kappa^{\,r}_{\beta,x}(D)
 \le17496e^{L_0}\beta^{3/2}
                  e^{-(229/384-\eta)\beta},
                    \qquad x\in\mathcal G .
\end{equation}
Thus a verified \(\eta<229/384\) would retain
exponential escape. This is a sufficient criterion,
not a bound already established for the other
parents. If the oscillation is larger, the estimate
is simply inconclusive.

The need for this qualification is substantive.
For a hidden field on either certified orbit,
\(T_x(I)=I\). In physical coordinates, multiplication
by the positive factor
\[
 r_\beta(T_x(B),x)
             =\exp\{-\beta H(x)\cdot b+\beta Lb_0\},
                         \qquad L>0,
\]
changes the conditional density exactly to
\(e^{\beta Lb_0}\,dB\), after normalization. Its
mass tends to one in every neighbourhood of \(B=I\):
outside a fixed angular neighbourhood the exponent
has a strict loss, whereas integration over a
smaller cap gives a positive-volume lower bound.
Continuity and \(T_x(I)=I\) then imply
\(\kappa^{\,r}_{\beta,x}(D)\to1\). This artificial
factor is not claimed to arise from another YM
parent. It demonstrates why positivity alone of
uncontrolled factors does not preserve the theorem.

Accordingly the next task is twofold: determine
the actual shared-coarse-link factors in the
whole retained repair law, and control the
posterior weak-staple region \(\mathcal G^c\)
or replace it by another certified repair region.
The strong-staple conditional normalization is
now explicit; a hidden Poincare inequality is
not one of its assumptions.

\subsection{Exact checks and preservation record}

The checker
\path{scripts/verify_ym_f13_v11_pivot_heatbath.py}
verifies the 944-group/96-pivot geometry, absence
of parent pivots in every nonroot source family,
and the six single-occurrence plaquettes of
\(\ell_0\). A scan of all coarse chronological
paths finds this physical pivot only in its own
two duplicate root paths; no exterior coarse
source depends on it. The checker verifies the Wilson linear-functional
identity on 256 exact noncommutative rational
quaternion cases. Those cases are regression
checks, not the universal proof; that proof is
Lemma~\ref{lem:f13v11-staple}. The displacement,
separation exponent, Haar-cap prefactor, and
\(\beta=64\) comparison are checked with rational
arithmetic. The global diffeomorphism and Haar
change of variables are inherited analytic
inputs, not established by finite tests.

The checker's SHA--256 is
\begin{center}\scriptsize\ttfamily
99A1784E06D1F49FF5A52EA4A278DE34CBE1D8BC955AE12F1001DAF9B9B471E1.
\end{center}
Predecessor V10 contains 185,526 bytes and
SHA--256
\begin{center}\scriptsize\ttfamily
82B2940B04B1C0B4D5FAC04072DF6371D7F0349EA16C9CBD268D1A37D864CED4.
\end{center}
Its full body is retained apart from the title
version. Only active V10 is replaced by active
V11. All archives and companion notes are
unchanged. The next companion audit remains
V15 and direction review V20.

\begin{firewall}
This version proves a global-in-hidden-field
conditional identity and a quantitative escape
bound on the explicitly defined strong-staple
region, together with its posterior-weighted
form consequence. It does not prove that this
region has a uniformly positive posterior
weight, that other-parent factors satisfy
\eqref{eq:f13v11-reweight} with a useful exponent,
or that the full retained law has a gap.
The ultraviolet and infinite-volume continuum
construction and physical Yang--Mills mass gap
remain unproved. The full goal stays active.
\end{firewall}
\section{V12: fixed-coarse interface descent and an admissible repair floor}
\label{sec:f13v12}

\subsection{State-space audit: coarse inputs are not fine-interface coordinates}

The upstream audit found a distinction that must precede
any comparison of additional parent factors. In f12 V21
and V28, the complete coarse field \(C\) is fixed.
The block variables \(x_b=Y_{I_b}\) and the remaining
fine-interface variables \(v=Y_{\mathcal S}\) are
coordinates \emph{within that fixed coarse fibre}.
In particular f12 (32.2)--(32.5) read
\[
 V(Y;C)=V_{\mathcal S}(v;C)+\sum_bU_b(x_b;v,C),
                      \qquad C\ \hbox{fixed}.
\]
The repair law in f12 V80 updates components of \(v\),
not components of \(C\).

V9--V11 explicitly varied a coarse input and proved
statements for the corresponding partial coarse
integrals or enlarged one-parent law. Those statements
remain statements about the laws defined there.
However, their coarse-link resampling does not preserve
the original fixed-\(C\) fibre. V10's reference to the
f12 V80 factorization supplies an abstract conditional
variance identity, not an identification of these
different variables or update mechanisms. V11's
coarse-link escape bound cannot be inserted as the
missing fixed-\(C\) interface repair floor.
This section corrects that handoff explicitly.

There is a further bookkeeping reason for care: f12
V87 identifies the full Wilson action with the regional
parent action up to a constant under \emph{hidden}
variation at fixed exterior. Such a constant can depend
on a newly varied exterior or coarse input. Here the
full Wilson action and the full inverse-pivot amplitude
are used throughout; no such constant is dropped.

We now return to a genuine free fine-interface group,
keep every coarse link equal to \(I\), and acquire a
rate-one heat-bath estimate on an explicit joint
region. This both avoids the incorrect state-space
handoff and uses the noncentral wells already certified.

\subsection{A coordinate in the canonical interface}

Use the actual fine torus \((\mathbb Z/6\mathbb Z)^4\)
and parent \(P=\{0,1\}^4\). There are 3645 global
independent fine groups after tree fixing and pivot
elimination. The parent contains 944 of them, leaving
2701 independent exterior groups. Select
\[
 \ell=((0,0,1,5),2),\qquad y=Y_\ell,\qquad y_*=-I .
\]
This is a nonpivot boundary link, not a tree link.
It is outside \(I_P\). Its coarse edge and that
edge's pivot are
\[
 e_\ell=((0,0,0,2),2),\qquad
 p_\ell=((0,0,1,4),2).
\]
Both endpoints of \(e_\ell\) are outside \(P\).
The value 5 is the torus coordinate; the selected
link is immediately across the periodic interface
from part of the parent collar.

The link \(\ell\) occurs in exactly two chronological
source paths, both over \(e_\ell\). None of the
fourteen nonroot sources over that edge contains a
parent hidden link or a pivot. At the central exterior
all fourteen source products equal \(-I\). Exactly
two of them contain \(\ell\).

At least one plaquette contains both \(\ell\) and a
parent hidden group. Thus \(\ell\) is a genuine
exterior input of this parent, not a remote
spectator coordinate. In the one-parent packing it
is an interface variable. If this local stencil is
placed in a compatible separated packing satisfying
f12 V28's input-closure condition
\(J_{P_b}^{\partial}\cap\bigcup_c I_c=\varnothing\),
the same exterior-input status puts \(\ell\) in
the interface. No arbitrary packing is presumed
to contain the certified finite stencil.

Write \(x\in\SU(2)^{944}\) for the parent hidden
groups and \(\bar v\) for the other 2700 independent
exterior groups. For the uniform local result below
allow
\begin{equation}\label{eq:f13v12-protected}
 d_{\rm HS}(\bar v_i,\bar v_i^*)\le1/16
          \quad\hbox{for all these exterior groups},
\end{equation}
where \(\bar v^*\) is the fixed central exterior.
The coarse field remains \(C_e=I\) on every edge.
Only a finite source collar is needed for the
clearance argument; the displayed condition is a
convenient stronger finite-torus condition.

\begin{lemma}[An exact inactive coordinate ball]
\label{lem:f13v12-inactive}
For every hidden field \(x\), every exterior
\(\bar v\) in \eqref{eq:f13v12-protected}, and every
\(y\) with \(d_{\rm HS}(y,-I)\le1/4\), the completed
pivot at \(e_\ell\) equals \(I\) and its inverse
Haar Jacobian equals one. Changing \(y\) inside
this ball changes no other pivot. Thus only the
physical link \(\ell\) changes, while the complete
coarse field stays fixed.
\end{lemma}

\begin{proof}
The fourteen nonroot source products are independent
of \(x\). Their central value is \(-I\). For a path
containing \(\ell\), bi-invariance and the product
triangle inequality bound its displacement from
\(-I\) by \(1/4+1/16=5/16\). For the other paths
the bound is at most \(2/16\). Tree links stay fixed.
Since \(d_{\rm HS}(I,-I)=\pi\sqrt2>4\), every such
product has distance at least \(4-5/16>\rho\)
from \(I\), where \(\rho=1/64\).

At the candidate pivot \(B=I\) all cutoff terms
vanish in an open neighbourhood. Its forward map
there is \(B\mapsto B\), with Haar Jacobian one.
Global uniqueness of completion identifies this
candidate with the actual pivot. The global
chronological-path enumeration shows that no
other pivot has \(\ell\) in a source. Every other
pivot and its inverse-Jacobian factor is therefore
independent of \(y\).
\end{proof}

The lemma is uniform over the full hidden compact
domain. It does not restrict \(x\) to a commuting
chart or to either known well.

\subsection{Certified force at both hidden minima}

At fixed \((x,\bar v)\), use the angular right
variation \(y(t)=-\exp(tE_i)\), and define
\[
 g_i(x,\bar v)=
   \left.\frac{d}{dt}S(x,\bar v,-\exp(tE_i))\right|_{t=0}.
\]
Here \(S\) is the full completed Wilson action at
\(C=I\), not a minimized action. The vector \(g\)
is in the angular normalization; an \(E_i\)
generator has Hilbert--Schmidt norm \(\sqrt2\).

\begin{proposition}[An admissible interface force at the two orbits]
\label{prop:f13v12-force}
At the commuting representative of each certified
orbit, with \(\bar v=\bar v^*\), the third component
satisfies the following outward rational enclosures:
\[
\begin{array}{c|cc}
 &\text{lower bound for }g_3&\text{upper bound for }g_3\\ \hline
\mathcal O_1&-499560279/250000000&-399648223/200000000\\
\mathcal O_2&-124890067/62500000&-1998241071/1000000000
\end{array}
\]
In both cases
\[
                  -2<g_3<-199/100 .
\]
After simultaneous conjugation the force vector
rotates and has unchanged norm. In particular both
hidden orbits are regular for the enlarged
\emph{fixed-coarse} fine-variable problem.
\end{proposition}

\begin{proof}
The checker constructs the full completed field
on the inherited point enclosures. For the second
orbit it reruns V8's scalar inclusion before using
the radius \(2/10^{12}\) enclosure. The six oriented
plaquettes at \(\ell\) give the angular physical force
\[
        2\sum_{q\ni\ell}\sigma_{q,\ell}\sin\phi_q .
\]
Lemma~\ref{lem:f13v12-inactive} proves that the
selected pivot response is zero, so this is also
the actual completed free-coordinate derivative.
Outward integer interval multiplication gives
the displayed enclosures. No floating-point
minimum is substituted for an enclosed critical
point. Conjugation equivariance proves the
statement on each orbit.
\end{proof}

The checker also audits all exterior free-coordinate
forces, including nonzero external-pivot responses.
For a nonparent pivot whose central sources have
\(N_+\) positive products, its longitudinal response
to an exterior free angle is
\[
 b'_e=-\frac{\sum_{\text{positive paths}}p'_j}{16-N_+}.
\]
Those terms are added to the physical forces.
The largest scanned force uses such a response;
omitting it would be incorrect. The chosen
\(\ell\) instead has all fourteen sources negative
and zero response, allowing the exact finite ball
in Lemma~\ref{lem:f13v12-inactive}. The full scan
is not a coverage theorem for other configurations.

\subsection{A finite descent that stays on the coarse fibre}

Because only one physical link changes in the
inactive ball, the ordinary single-link Wilson
identity gives, for \(|n|=1\),
\begin{equation}\label{eq:f13v12-orbit-action}
\begin{split}
 &S(x,\bar v,-\cos t\,I-\sin t\,n\cdot E)
       -S(x,\bar v,-I)\\
 &\hspace{8mm}=A(x,\bar v)(1-\cos t)
                +g(x,\bar v)\cdot n\,\sin t,\qquad |A|\le12 .
\end{split}
\end{equation}
Indeed each of the six plaquettes is a linear
functional of the four real components of this
physical link, of coefficient norm two.
The combined coefficient norm is at most twelve.
In particular this action, on the inactive ball,
is 12-Lipschitz in Hilbert--Schmidt distance.

Define the hidden force region at exterior \(\bar v\)
by
\[
 \mathcal A_{\bar v}
        =\{x:|g(x,\bar v)|\ge199/100\}.
\]
It contains fixed neighbourhoods of both certified
orbits at \(\bar v^*\), and neighbourhoods of the
corresponding joint points \((\mathcal O_j,\bar v^*)\).
This last statement uses continuity and the strict
interval margins, not a numerical radius for those
neighbourhoods.

\begin{proposition}[Rational endpoint with strict action decrease]
\label{prop:f13v12-descent}
For every \(x\in\mathcal A_{\bar v}\), take
\(n=-g/|g|\) and
\[
 y_1(x,\bar v)=-\frac{255}{257}I
                         -\frac{32}{257}\,n\cdot E .
\]
This is a unit quaternion. It lies at
Hilbert--Schmidt distance less than \(3/16\) from
\(-I\), and
\begin{equation}\label{eq:f13v12-descent}
 S(x,\bar v,-I)-S(x,\bar v,y_1)
       \ge\frac{992}{6425}>\frac3{20}.
\end{equation}
The entire short rotation to \(y_1\) keeps \(C=I\).
\end{proposition}

\begin{proof}
The squared rational coefficients sum to one.
The angle is \(t=2\arctan(1/16)<1/8\).
Using \(\sqrt2<3/2\) gives
\(d_{\rm HS}(-I,y_1)=\sqrt2t<3/16<1/4\),
so Lemma~\ref{lem:f13v12-inactive} applies along
the full short arc. Equation
\eqref{eq:f13v12-orbit-action} and the choice of
axis give a decrease at least
\[
 \frac{199}{100}\frac{32}{257}
                   -12\left(1-\frac{255}{257}\right)
        =\frac{992}{6425}.
\]
\end{proof}

At the displayed orbit representatives the axis
is \(n=(0,0,1)\). Only the interface fine link is
changed, not a coarse input. This is the key
difference from the rational coarse endpoint in
f12 V92 and V9.
\subsection{Equal-Haar-ball comparison for the actual heat bath}

Fix \(x,\bar v,C=I\). Let
\(\kappa_{\beta,x,\bar v}\) be the canonical
single-free-group heat bath for \(y\) in the
original full fixed-\(C\) density, with every
other free group held fixed. It is integrated
over the full compact group, including regions
where the pivot cutoff is active.
Set
\[
 \varepsilon=1/1000,\qquad
 D=\{y:d_{\rm HS}(y,-I)\le\varepsilon\},\qquad
 D_1(x,\bar v)=\{y:d_{\rm HS}(y,y_1)\le\varepsilon\}.
\]

\begin{theorem}[Fixed-coarse conditional escape]
\label{thm:f13v12-return}
For every protected \(\bar v\), every
\(x\in\mathcal A_{\bar v}\), and every \(\beta>0\),
\begin{equation}\label{eq:f13v12-return}
 \kappa_{\beta,x,\bar v}(D)
 \le e^{-(4189/32125)\beta}
 \le e^{-\beta/8}.
\end{equation}
For \(\beta\ge8\), return is less than \(1/2\);
a rate-one update therefore exits \(D\) with
probability greater than \(1/2\).
\end{theorem}

\begin{proof}
Both balls lie inside the inactive ball:
\(3/16+1/1000<1/4\).
Their Haar volumes are exactly equal by
translation invariance. On their union the
selected inverse-pivot Jacobian is one, and
every other amplitude factor is the same
positive constant in \(y\). Thus the full
amplitude cancels in a comparison of their
unnormalized masses.

By the 12-Lipschitz bound and
\eqref{eq:f13v12-descent},
\[
 \inf_{y\in D}S(x,\bar v,y)
       -\sup_{z\in D_1}S(x,\bar v,z)
   \ge\frac{992}{6425}-24\varepsilon
   =\frac{4189}{32125}>\frac18.
\]
Bound the numerator over \(D\) using its
infimum, and the denominator of the full heat
bath from below by the integral over \(D_1\),
using its supremum. The equal Haar volumes
and identical amplitudes cancel exactly.
All other denominator contributions are
nonnegative; none is presumed small or
discarded from the actual conditional law.
This proves \eqref{eq:f13v12-return}.
At \(\beta\ge8\), \(e^{-\beta/8}\le e^{-1}<1/2\).
\end{proof}

There is no \(\beta^{1415}\) hidden-volume factor,
no one-link Laplace prefactor, and no altered
coarse-input distribution in this estimate.
The conditional comparison uses two balls of
the same physical free group. Its small radius
is fixed, and the estimate is uniform over all
hidden configurations satisfying the force
condition. It is not a claim that this condition
holds with large posterior probability.

\subsection{An original-clock joint floor and its marginal limitation}

Consider the full fixed-\(C\) fine-coordinate
heat-bath generator with rate one for the
selected interface group, as in f12 V80.
Define the joint region
\[
 \mathcal R=
 \{(x,\bar v,y):\bar v\text{ satisfies }
   \eqref{eq:f13v12-protected},\
                 x\in\mathcal A_{\bar v},\ y\in D\}.
\]

\begin{corollary}[A killed floor in the actual fibre]
\label{cor:f13v12-killed}
The full fine-coordinate heat bath, killed on
leaving \(\mathcal R\), has Dirichlet floor
at least \(1-e^{-\beta/8}\). In particular its
floor is greater than \(1/2\) for \(\beta\ge8\).
\end{corollary}

\begin{proof}
Let \(F\) vanish outside \(\mathcal R\).
Condition on all groups except \(y\).
If \((x,\bar v)\) is outside the stated
protected force region, \(F\) is zero for
all \(y\). Otherwise \(F\) is supported in
\(D\), and Cauchy--Schwarz yields
\[
 \left|\int F\,d\kappa_{\beta,x,\bar v}\right|^2
 \le\kappa_{\beta,x,\bar v}(D)
                         \int|F|^2\,d\kappa_{\beta,x,\bar v}.
\]
Theorem~\ref{thm:f13v12-return} therefore bounds
the conditional variance from below by
\((1-e^{-\beta/8})\int|F|^2\).
Integrate the conditioning variables. The
other heat-bath contributions are nonnegative,
so the full form has the same lower bound.
Taking the Rayleigh infimum proves the claim.
\end{proof}

This region is not a visible cylinder: it
restricts hidden configurations through
\(\mathcal A_{\bar v}\). The hidden condition
cannot be removed when applying f12 V80's
visible-cylinder transfer.

For the single-parent packing, fix a protected
\(\bar v\) and let \(\mu_\beta^{\bar v}(dx,dy)\)
be the full fixed-\(C\) law conditioned on it.
Let \(\nu_\beta^{\bar v}(dy)\) be its marginal,
and put
\[
 a_\beta(y,\bar v)
    =\mu_\beta^{\bar v}(\mathcal A_{\bar v}\mid y).
\]
The compressed joint update, and hence the
ideal marginal heat-bath form, obey for
\(f\) supported in \(D\)
\begin{equation}\label{eq:f13v12-weighted}
 \mathcal E(f)\ge
 (1-e^{-\beta/8})
       \int_D a_\beta(y,\bar v)|f(y)|^2\,d\nu_\beta^{\bar v}(y).
\end{equation}
This follows by the same conditional-variance
calculation, now integrating only over
\(x\in\mathcal A_{\bar v}\), and by f12 V80's
exact marginal form domination. The weight
is the actual fixed-coarse hidden posterior,
not the posterior of V11's variable-coarse law.

An independently proved lower bound on this
weight throughout \(D\) would give a visible
repair floor. Such a bound is still missing.
For a many-parent packing, all additional
hidden inputs must likewise remain in the
conditioning or be integrated with their
true posterior. This section proves the
finite fixed-parent stencil and its full
fine-fibre update; it does not assert an
all-packing or all-exterior marginal estimate.

\subsection{What changes in the acquisition order}

The unknown additional coarse-parent factor
from V11 is not used as the next fixed-fibre
repair parameter. First match the coordinates
and update law, as done here. The force region
\(\mathcal A_{\bar v}\) is now the appropriate
one for the selected admissible interface
update. The immediate missing input is control
of its complementary weak-interface-force
region under the original hidden posterior,
or a further collection of admissible updates
covering that region.

Both known hidden minima are already inside
this acquired joint repair region for central
exterior data. Re-establishing their coarse
force or local hidden Poincare constant would
not address the remaining coverage problem.
The next work should examine the actual
interface score on configurations outside
these regions, preserving \(C\) and every
normalizing factor.

\subsection{Verification and append record}

The exact checker is
\path{scripts/verify_ym_f13_v12_interface_force.py}.
It verifies the global free-coordinate count,
the selected link's canonical exterior status,
all its chronological source occurrences,
the fourteen negative source products at the
central exterior, the source exclusion of
hidden variables, and the actual interval
forces at both enclosed critical points.
The rational endpoint, equal-ball radius,
and energy-separation constants are checked
with integer/Fraction arithmetic.
The analytic proofs above supply the
uniform hidden-field extension and the
conditional probability and form statements.

The checker's SHA--256 is
\begin{center}\scriptsize\ttfamily
587CEB39AFE10A9D191AB0E85F597ADA258F499A79AA4809F97541D30262D8A2.
\end{center}
The predecessor V11 contains 201,680 bytes
and SHA--256
\begin{center}\scriptsize\ttfamily
B317AD788D8DED3ACB036575DAD3D9F54B46EDAC42D1076C7E70A88B7F6629E1.
\end{center}
Its full body is preserved, with only the
title updated. The coordinate-handoff
qualification is appended here rather than
silently rewriting earlier claims. Only
active V11 is replaced by V12. Archives and
companion files remain unchanged; the next
companion audit is V15.

\begin{firewall}
V12 corrects the coarse-versus-interface
handoff and proves a finite descent, a
normalized conditional escape bound, and
an original-clock joint killed floor within
the actual fixed-coarse fibre. The marginal
consequence retains its unknown posterior
weight. There is still no unrestricted
retained gap, uniform cofinal clustering,
continuum construction, or physical
Yang--Mills mass gap. The full goal remains
unproved and active.
\end{firewall}
\section{V13: an all-hidden interface shell and a marginal repair floor}
\label{sec:f13v13}

\subsection{The local posterior bottleneck is bypassed}

V12's single-link repair required a lower bound on
the hidden posterior of a force region. For the
present central-interface neighbourhood this version
removes that requirement. A collective direction on
133 actual fine-interface groups has negative action
curvature uniformly over \emph{every} hidden field.
The 84 plaquettes depending on hidden variables are
bounded in their worst possible direction, rather
than excluded probabilistically.

The new result is a positive, large-\(\beta\)-uniform
killed floor for the original interface heat bath
on a fixed visible neighbourhood. It uses the
fixed-\(C\) state space corrected in V12 and the
pre-integration clock comparison of f12 V77--V80.
Unlike that archive's cleared-sheet application,
the direction here meets actual hidden-dependent
plaquettes. Their contribution is included in the
curvature budget.

All statements in this version concern the declared
fine torus of side six, its one parent
\(P=\{0,1\}^4\), and \(C=I\). They do not assert
a volume-uniform extension of this wrapping shell,
an arbitrary-interface repair theorem, or a global
spectral gap.

\subsection{The 133-group shell and its inactive pivots}

Let \(\mathcal S\) be the 2701 fine-interface groups
outside the 944-group parent, and \(v^*\) their
central values. Define
\begin{equation}\label{eq:f13v13-shell}
 T=\{((a,b,1,c),2):a,b,c\in\{0,\ldots,5\},\
       Y^*_{((a,b,1,c),2)}=-I,\
                 ((a,b,1,c),2)\notin I_P\}.
\end{equation}
Then \(|T|=133\). Every member is an independent
nonpivot boundary group. These groups occupy
the seven nonroot boundary positions over each
of 19 direction-2 coarse edges outside \(E(P)\).

Each group in \(T\) occurs in exactly two
chronological paths, both over its own coarse
edge. Over each of the 19 edges all fourteen
nonroot source products have central value
\(-I\), contain no hidden parent group, and
contain no pivot. The sources of every other
pivot do not contain groups in \(T\).
Consequently, on a product neighbourhood of
the central interface in which the groups in
\(T\) are within \(1/4\) of \(-I\) and the
other interface groups within \(1/16\) of
their central values, all 19 pivots equal
\(I\) and their inverse Haar Jacobians are one.
The proof is precisely the source-distance
argument of Lemma~\ref{lem:f13v12-inactive},
applied to these 19 disjoint source families.
It is uniform over all parent hidden fields.

It follows that varying the groups in \(T\)
within this neighbourhood changes only those
physical links. Every coarse link remains
fixed. The complete inverse-pivot amplitude
is independent of the groups in \(T\) there;
it may still depend on the hidden field and
on the other interface coordinates.

\subsection{An exact plaquette budget independent of hidden fields}

At the central interface, keep \(x\in\SU(2)^{944}\)
arbitrary. Fix an imaginary unit quaternion \(E\)
and rotate all selected groups together:
\[
             Y_\ell(t)=-\exp(tE),\quad \ell\in T .
\]
All unselected interface groups stay central,
and the hidden groups are fixed. For \(|t|\le1/8\)
the preceding inactive-pivot statement applies.

The exact plaquette enumeration is:
\begin{center}
\begin{tabular}{lrr}
\toprule
Type&Number&Contribution to \(S''(0)\)\\
\midrule
Two selected links, fixed positive plaquette&306&0\\
One selected link, fixed negative plaquette&102&\(-2\)\\
One selected link, hidden-dependent plaquette&84&at most \(2\)\\
\bottomrule
\end{tabular}
\end{center}
There are no remaining plaquettes involving \(T\).
The incidence identity
\[
                  2\cdot306+102+84=6\cdot133
\]
checks all six incidences per selected link.

\begin{theorem}[A non-cleared all-hidden negative direction]
\label{thm:f13v13-curvature}
For every hidden field \(x\),
\begin{equation}\label{eq:f13v13-curvature}
 \left.\frac{d^2}{dt^2}S(x,Y(t))\right|_{t=0}
                  \le-36.
\end{equation}
The unnormalized direction has squared
Hilbert--Schmidt norm \(266\). Its unit
normalization therefore satisfies
\[
                    \operatorname{Hess}S(a,a)
                            \le-\frac{18}{133}
\]
at the central interface, uniformly over \(x\).
\end{theorem}

\begin{proof}
For a two-selected-link plaquette the occurrences
have opposite orientation and the other links
are central and independent of the parent.
The two common rotations therefore cancel
exactly, not merely to second order.

A fixed negative one-selected-link plaquette
has action \(2+2\cos t\), so contributes \(-2\)
to the second derivative. A hidden-dependent
one-selected-link plaquette is a real linear
functional of \(\cos t\) and \(\sin t\), plus
its constant action term. Multiplication by
the other unit quaternions preserves norm,
so its cosine coefficient in the action
difference has absolute value at most two.
This estimate is valid even when the hidden
field is noncommutative or far from either
known orbit.

Summing the exact three classes gives
\(-2\cdot102+2\cdot84=-36\).
There are 133 generators of squared
Hilbert--Schmidt norm two. Dividing the
second derivative by \(266\) gives the
unit-direction bound. Along the displayed
curve the product-group acceleration is
zero, so its second derivative is the
intrinsic Hessian in that direction.
\end{proof}

There is also a useful finite, fixed-coarse
descent statement. Along this same curve,
\[
 S(x,Y(t))-S(x,v^*)
                  =A(x)(1-\cos t)+B(x)\sin t,
                         \qquad A(x)\le-36 .
\]
Here \(B\) can have either sign. Choose the
sign of \(t=2\arctan(1/16)\) so that
\(B\sin t\le0\). Then
\[
 S(x,v^*)-S(x,Y(t))\ge
              36\left(1-\frac{255}{257}\right)
                      =\frac{72}{257}>\frac14 .
\]
This holds for every \(x\), without knowing
which hidden configurations dominate.
It is a geometric descent, not a claim
that the original chain executes a 133-group
simultaneous update. The original-clock
conclusion is established separately below.

\subsection{A visible cylinder with negative measured curvature}

Use product Hilbert--Schmidt orthonormal
normal coordinates \(z\) for all 2701
interface groups about \(v^*\). Leave all
hidden groups uncharted and unrestricted.
Let \(a\) be the constant Euclidean unit
vector having coordinate \(1/\sqrt{133}\)
in the \(E/\sqrt2\) direction of each
selected group, and zero elsewhere.
At \(z=0\) the corresponding physical
curve has parameter \(t=s/\sqrt{266}\).

Joint smoothness of the global completion
and compactness of the entire hidden domain
make \(\partial_a^2S(x,z)\) uniformly
continuous in \(z\). By
Theorem~\ref{thm:f13v13-curvature}, a fixed
open product cylinder \(\mathcal D^+\)
about \(z=0\) can be chosen so that
\begin{equation}\label{eq:f13v13-uniform-chart}
                  \partial_a^2S(x,z)\le-9/133
                   \quad(x\in M,\ z\in\mathcal D^+).
\end{equation}
Choose it inside the inactive-pivot
neighbourhood and small enough that its
inverse-metric lower bound is \(c_g=1/2\).
The radius exists; no numerical radius
is certified in this version.

Relative to \(dx\,dz\), where \(dx\) is
hidden product Haar, the joint measured
potential is
\[
 \Phi_\beta(x,z)=\beta S(x,z)-\log j(x,z)
                     -\sum_{i\in\mathcal S}\log j_G(z_i).
\]
The full inverse-pivot factor \(j\) is
independent of the selected group variables
in this neighbourhood. Thus
\(\partial_a^2\log j=0\). The remaining
coordinate-Haar derivative is bounded in
absolute value by a fixed \(C_{\rm Haar}\).
After a fixed large-\(\beta\) threshold,
\begin{equation}\label{eq:f13v13-measured}
                 \partial_a^2\Phi_\beta
                   \le-\eta\beta,\qquad \eta=9/266,
\end{equation}
uniformly over all hidden configurations.
For example the Haar contribution is
absorbed once \(\beta\ge266C_{\rm Haar}/9\).

The actual interface marginal in these
coordinates is defined by
\[
 \Psi_\beta(z)=-\log\int_M e^{-\Phi_\beta(x,z)}\,dx .
\]
Differentiating the finite positive integral
gives
\[
 \partial_a^2\Psi_\beta
   =\mathbb E_z[\partial_a^2\Phi_\beta]
                    -\operatorname{Var}_z(\partial_a\Phi_\beta)
                  \le-\eta\beta .
\]
This marginal curvature bound contains no
posterior coverage hypothesis. The
covariance has the favourable sign.
For the heat-bath clock, however, we work
before integration so as to retain the
linear-in-\(\beta\) mixed-Hessian bound.

\subsection{Transfer to the original interface heat bath}

Let \(\mu_\beta(dx,dv)\) be the actual
fixed-\(C=I\) full fine-fibre law on this
torus, and \(\nu_\beta(dv)\) its exact
interface marginal after integrating
the 944 hidden groups. Use the original
rate-one full-group heat baths
\[
 H_\mu=\sum_{\text{all free groups }i}(I-\Gamma_i),
 \qquad
 H_\nu=\sum_{i\in\mathcal S}(I-P_i),
\]
with the canonical conditional laws of
these two measures. No collective update
has been added to either generator.

Choose a smaller visible product cylinder
\(\mathcal D\) with closure in
\(\mathcal D^+\), and a smooth visible
cutoff \(\chi\) equal to one on
\(\mathcal D\), supported in
\(\mathcal D^+\). Set
\(M_\chi=\sup|\nabla\chi|^2<\infty\).
Both cylinders restrict only interface
coordinates; the joint cylinder is
\(M\times\mathcal D\).

\begin{theorem}[A genuine all-hidden marginal repair floor]
\label{thm:f13v13-floor}
There are fixed \(C_{\rm f}>0\) and
\(\beta_0<\infty\) such that, with
\(c_0=9/1064\),
\begin{equation}\label{eq:f13v13-floor}
 \kappa_{\mathcal D}^{H_\nu}
 \ge\frac{c_0\beta}
       {2\{c_0\beta+8C_{\rm f}(\beta+1)\}}
 \ge\frac{c_0}{2(c_0+16C_{\rm f})}>0,
                         \qquad\beta\ge\beta_0 .
\end{equation}
This is the bottom of the killed
interface heat bath, not the gap of a
reflected chain on \(\mathcal D\).
No hidden Poincare inequality or lower
bound on a hidden force-region weight
is assumed.
\end{theorem}

\begin{proof}
We spell out the transfer to distinguish
the new stencil input from the inherited
operator argument. Let \(L\) be the
global weighted diffusion for \(\mu_\beta\).
Ground-state conjugation in the constant
visible chart direction gives, for \(h\)
compactly supported in the chart variables
of \(M\times\mathcal D^+\),
\[
 E_L(h)\ge\lambda_\beta\|h\|^2,\qquad
 \lambda_\beta=\frac{c_g\eta\beta}{2}
                        =\frac9{1064}\beta.
\]
Indeed, writing \(u=h e^{-\Phi_\beta/2}\),
integration by parts gives the potential
term
\(|\partial_a\Phi_\beta|^2/4
 -\partial_a^2\Phi_\beta/2\ge\eta\beta/2\);
the metric lower bound is \(c_g=1/2\).
Hidden variables are integrated over
their entire compact domain in this
argument.

The global unintegrated Haar potential
\(W_\beta=\beta S-\log j\) is a sum of
fixed smooth finite-stencil terms.
Its mixed intrinsic Hessian row therefore
has a bound \(K_\beta\le C_{\rm f}(\beta+1)\).
The square estimate and global resolvent
comparison of f12 (81.2)--(81.3) apply:
\[
                         H_\mu\succeq L(L+K_\beta)^{-1}.
\]
The comparison uses global operators
for this same law. It does not replace
their resolvent by one with Dirichlet
boundary conditions.

Apply the local diffusion inequality to
\(\chi u\), for a global diffusion-form
vector \(u\). The product rule gives
\[
 \|1_{M\times\mathcal D}u\|^2
 \le \frac2{\lambda_\beta}E_L(u)
                   +\frac{2M_\chi}{\lambda_\beta}\|u\|^2.
\]
For the spectral projector
\(\Pi_t=1_{[0,t]}(L)\), this implies
\(\|1_{M\times\mathcal D}\Pi_t\|^2
\le2(t+M_\chi)/\lambda_\beta\).
Taking adjoints yields the same bound
for the low-spectral mass of a vector
supported in the cylinder. If
\(\lambda_\beta\ge8M_\chi\), take
\(t=\lambda_\beta/8\). Such a vector
has at least half its squared norm
above \(t\). The resolvent comparison
therefore gives
\[
 \kappa_{M\times\mathcal D}^{H_\mu}
       \ge\frac{\lambda_\beta}
                        {2(\lambda_\beta+8K_\beta)}.
\]

Finally the isometric lift \(Jf(x,v)=f(v)\)
has the same squared norm and support in
\(M\times\mathcal D\) when \(f\) is supported
in \(\mathcal D\). Conditional variance
decomposition gives
\(E_{H_\nu}(f)\ge E_{H_\mu}(Jf)\).
Hidden-coordinate updates leave \(Jf\)
unchanged. The joint cylinder contains
all hidden configurations, so there is
no missing weight in this transfer.
Substitution of \(\lambda_\beta=c_0\beta\)
and \(K_\beta\le C_{\rm f}(\beta+1)\)
proves the first inequality. Enlarge
\(\beta_0\) to include \(\beta\ge1\);
then \(\beta+1\le2\beta\) proves the
second one.
\end{proof}

The finite constants \(C_{\rm f}\),
\(M_\chi\), the cylinder radius, and
the large-\(\beta\) threshold have not
been numerically enclosed. The proof
establishes their existence and the
stated \(\beta\)-uniform positive
floor for this fixed geometry.
It is not a numerically specified
mixing-time bound at a chosen coupling.

\subsection{Verification, scope, and the next missing region}

The exact checker is
\begin{center}\small
\path{scripts/verify_ym_f13_v13_interface_shell.py}.
\end{center}
It verifies all 133 free-coordinate indices,
the 19 independent inactive source
families, and every plaquette in the
curvature budget. In particular it
checks that every two-selected-link
plaquette has no parent-dependent
physical link and has opposite selected
orientations. The worst-case bound for
the 84 remaining hidden-dependent
plaquettes is analytic and valid for
arbitrary unit quaternions. The checker
also verifies the Hilbert--Schmidt
normalization, rational descent, and
the coefficient \(c_0\).

Its SHA--256 is
\begin{center}\scriptsize\ttfamily
4398F8522BE06AF9D083B20CD67C08294BEE24C94F579B063791F112A95DC120.
\end{center}
No new critical-point search was needed:
the new estimate is independent of
which hidden field is present. In
particular it includes both V8 wells
and every weak-force configuration
that V12 left out.

The local posterior obstacle is thus
removed for the present visible central
neighbourhood. It has not been removed
for all interface configurations.
The next issue is coverage by such
certified visible repair cylinders,
including deformations of the central
interface and the effect of other
packing geometries. This periodic
six-site shell cannot be declared a
non-wrapping, volume-uniform stencil
without checking its changed incidences.
Likewise a killed floor on one cylinder
does not bound the unrestricted
interface Poincare constant.

Predecessor V12 contains 219,745 bytes
and SHA--256
\begin{center}\scriptsize\ttfamily
C7DFB012367EFD34DF85B535F3D3B4FA5F4ACBDF139AE95F52BD6BBE9BE6A938.
\end{center}
Its complete body is retained apart
from the title version. Active V13
replaces V12, leaving all archives
and companion files unchanged. The
next companion checkpoint remains V15.

\begin{firewall}
This version establishes an all-hidden,
non-cleared local interface repair
floor for the specified original
fixed-coarse law and original
single-group heat-bath clock.
It closes the posterior-weight
requirement for this particular
central-interface cylinder, not for
the entire compact interface.
No global cofinal gap, continuum
construction, or physical Yang--Mills
mass gap is claimed. The full goal
remains active and unproved.
\end{firewall}
\section{V14: a non-wrapping shell with a volume-uniform local repair floor}
\label{sec:f13v14}

\subsection{What survives when the periodic boundary is removed}

V13's 133-group shell used wraparound on a fine torus
of side six. That certificate cannot be inserted
unchanged into larger volumes. This version constructs
the stencil on the integer lattice, counts its true
outer boundary, and then embeds the complete input
closure into larger tori.

The same constant rotation on three- and four-cell
outer cubes does not pass the worst-case curvature
test. A five-cell outer cube does: its 819 selected
groups have Hilbert--Schmidt unit curvature at most
\(-4/21\), uniformly over every uncontrolled field.
A finite, explicitly constructed interface collar
suffices to prescribe the central data used in that
statement. Neither the rest of the fine volume nor
the parent hidden variables are restricted.

The resulting killed interface heat-bath floor is
uniform in the surrounding volume for the declared
embedded central pattern and visible-collar condition.
This is still a local repair theorem, not an
unrestricted interface gap or a continuum result.

\subsection{Integer-lattice geometry and the actual boundary cost}

All coordinates in this subsection are integers,
with no modular arithmetic. Keep
\(P=\{0,1\}^4\) and the canonical 944-group set
\(I_P\), including the boundary variables on
every coarse edge incident to \(P\). Its 96
pivots form the physical set \(B_P\). Set
\[
 \epsilon_\ell=
 \begin{cases}
 -1,&\ell=(x,2),\ x_2=1,\
                      \text{some }x_a\text{ is odd},\ a\in\{0,1,3\},\\
 +1,&\text{otherwise}.
 \end{cases}
\]
These signs prescribe the same local central
one-direction sheet as before, now without
identifying its opposite boundaries.

For \(q\ge3\), use the transverse coarse cells
\[
 Q_q=\{-1,0,\ldots,q-2\}^3\setminus\{0,1\}^3.
\]
For each \((a,b,c)\in Q_q\), select the seven
links
\[
 ((2a+u,2b+v,1,2c+w),2),\qquad
 (u,v,w)\in\{0,1\}^3\setminus\{(0,0,0)\}.
\]
Their union is \(T_q\), with
\(|T_q|=7(q^3-8)\). Every selected group is
free, has central value \(-I\), and is
outside \(I_P\). It belongs to a direction-2
coarse edge with both endpoint cells outside
\(P\). All fourteen nonroot source paths of
that edge avoid \(I_P\), have central product
\(-I\), and contain no pivot.

Enumerate the six plaquettes incident to each
selected link and remove duplicates. Classify
a plaquette as uncontrolled if it contains
a physical link in \(I_P\cup B_P\); otherwise
its central value is prescribed. The exact
counts are:
\begin{center}\small
\begin{tabular}{rrrrrrr}
\toprule
\(q\)&\(|T_q|\)&two selected&fixed negative&fixed positive&uncontrolled&
\(S''\) upper\\
\midrule
3&133&261&102&126&48&144\\
4&392&828&324&288&84&96\\
5&819&1845&690&450&84&\(-312\)\\
\bottomrule
\end{tabular}
\end{center}
Every two-selected plaquette has opposite
selected orientations, no parent-dependent
physical link, and positive central product.
Each of the other columns has one selected
link per plaquette. In particular, at \(q=5\),
\[
 2\cdot1845+690+450+84=6\cdot819.
\]
The positive upper bounds for \(q=3,4\)
mean that this particular sufficient
worst-case test fails; they do not exclude
other negative directions or sharper
configuration-dependent estimates.

For the rest of V14 fix \(q=5\) and write
\(T=T_5\). There are 117 associated inactive
pivot edges. A common small rotation of all
selected links changes no other physical
link when their source collars have the
prescribed central values.

\subsection{A finite visible input closure}

The central prescription must be imposed on
free coordinates, not on fictitiously
independent pivots. Here is the exact closure
used by the checker.

For a finite set of physical links \(A\),
let \(\mathfrak f(A)\) be obtained by
discarding tree links, retaining every
nonpivot free link, and replacing each pivot
by the free links in its fourteen nonroot
source paths. No iteration is needed:
nonroot source paths contain no pivots.

Let \(\mathcal P_{\rm fix}\) consist of the
two-selected and fixed one-selected
plaquettes in the \(q=5\) enumeration.
Let \(R_T\) be the free input groups of the
117 source families belonging to \(T\).
Define
\begin{equation}\label{eq:f13v14-collar}
 F=T\cup R_T\cup
       \mathfrak f\!\left(
            \bigcup_{p\in\mathcal P_{\rm fix}}\{\text{physical links of }p\}
                    \right).
\end{equation}
The exact finite enumeration gives
\[
                         |F|=10034,\qquad F\cap I_P=\varnothing .
\]
This is a sufficient collar, not a claim
of a minimal one. It is large but fixed
independently of ambient volume.

At \(Y_\ell=\epsilon_\ell I\) for
\(\ell\in F\), every required nonparent
pivot in a fixed plaquette equals \(I\):
its source products are central, so all
truncated logarithms vanish at \(B=I\),
and global uniqueness identifies the
completed pivot. Thus all fixed plaquette
claims hold independently of the values
of every free coordinate outside \(F\).
The 84 uncontrolled plaquettes are not
assigned central values.

The closure of \emph{all} plaquette inputs,
including the uncontrolled ones, contains
10431 free groups. The vertices needed for
these inputs and their defining paths have
coordinate bounds
\[
 [-3,9]\times[-3,9]\times[0,3]\times[-3,9].
\]
Translate by the even fine vector
\((8,8,8,8)\). All these coordinates then
lie strictly inside a fine torus of any
even side length \(L\ge32\), with no
wraparound or coincident stencil vertices.
Even translation preserves the parity-tree
and root-pivot convention. The incident
plaquettes and source paths are therefore
exactly the integer-lattice ones in every
such embedding.

For an interface marginal, require the
computed \(F\) to remain among the retained
fine-interface groups. This holds when only
the central parent is integrated. It also
holds for separated parent packings whose
other hidden sets avoid \(F\). We state this
buffer condition explicitly; no previously
chosen packing is silently changed. The
hidden sets outside this collar may be
arbitrarily large, and their variables
need not be near the section.

\subsection{The all-uncontrolled-field curvature certificate}

\begin{theorem}[Non-wrapping negative curvature]
\label{thm:f13v14-curvature}
Keep the coarse field \(C=I\). Prescribe the
central values on \(F\) and leave all other
free groups arbitrary. Under the common
rotation
\[
                     Y_\ell(t)=-e^{tE},\quad \ell\in T,
                    \qquad \|E\|_{\rm HS}=\sqrt2,
\]
the full completed action obeys
\[
                  \left.\frac{d^2}{dt^2}S(Y(t))\right|_{t=0}
                                    \le-312 .
\]
The squared norm of this direction is
\(2\cdot819=1638\). Hence its unit
normalization \(a\) satisfies
\begin{equation}\label{eq:f13v14-curvature}
                        \operatorname{Hess}S(a,a)\le-4/21.
\end{equation}
The estimate is uniform in \(L\ge32\)
and in every free group outside \(F\).
\end{theorem}

\begin{proof}
Each boundary group occurs in exactly two
chronological paths over its own coarse
edge, and no other source family.
The 117 associated source families are
independent of \(I_P\); their central
nonroot products are \(-I\). For a small
common rotation the cutoff terms at
candidate pivot \(I\) stay zero. Thus
those pivots are exactly \(I\), and every
other pivot is independent of \(t\).
Their inverse-Jacobian factors are
likewise independent of the selected
coordinates in this neighbourhood.

The two-selected fixed plaquettes cancel
the common rotation exactly. A fixed
negative one-selected plaquette contributes
\(-2\) to the second derivative, and a
fixed positive one contributes \(+2\).
Each uncontrolled one-selected plaquette
has second derivative at most \(2\), by
the unit-quaternion linear-functional
bound used in V13. No probability or
stationarity assumption is made about
the fields in that last class.
Consequently
\[
 S''(0)\le 2(-690+450+84)=-312 .
\]
Division by 1638 gives \(-4/21\).
All inputs and incidences lie in the
embedded finite stencil, proving the
claimed volume and exterior uniformity.
\end{proof}

The same enumeration also gives
\[
 S(Y(t))-S(Y(0))=A(1-\cos t)+B\sin t,\qquad A\le-312
\]
for \(|t|\le1/8\), with
\(F\setminus T\) held central.
Indeed each rotated source stays
within \(\sqrt2/8\) of \(-I\),
well outside the cutoff support.
Set
\[
                 t_* = 2\arctan(1/16)<1/8.
\]
Choose \(t\in\{t_*,-t_*\}\) so that
\(B\sin t\le0\). This gives a decrease
at least \(624/257\). As in V13, this
geometric descent does not replace the
original heat-bath clock by a collective
rotation kernel.

\subsection{Uniform local cylinders and the original-clock floor}

Use Hilbert--Schmidt orthonormal normal
coordinates \(z_F\) about the central
values on \(F\), leaving all other
variables uncharted. Normalize the
constant selected direction \(a\) to
Euclidean norm one. The action and
completion derivatives in this direction
involve a fixed finite set of compact
input groups. Their bounds are independent
of ambient volume. Uniform continuity
and \eqref{eq:f13v14-curvature} therefore
give a fixed product chart
\(\mathcal D_F^+\) on which
\[
                      \partial_a^2S\le-2/21
\]
for every value of the uncharted groups.
This assertion uses compactness only
of the fixed finite input closure,
not a compactness claim for an
infinite-dimensional product manifold.

Choose the chart inside the inactive
neighbourhood of the 117 source families
and with inverse-metric lower bound
\(c_g=1/2\). The full inverse-pivot
amplitude is independent of the
selected coordinates there. The
coordinate-Haar contribution has a
fixed bounded directional Hessian.
Thus, for sufficiently large \(\beta\),
the joint measured potential satisfies
\[
                    \partial_a^2\Phi_\beta\le-\beta/21 .
\]
The local diffusion coefficient in the
ground-state conjugation estimate is
therefore
\[
                   c_0=\frac{c_g}{2\cdot21}=\frac1{84}.
\]
The chart radius and the threshold are
existential uniform constants; the
integer checker does not numerically
enclose them.

Let \(\mu_{\beta,L}\) be the full
fixed-\(C=I\) fine-fibre law, and let
\(\nu_{\beta,L}\) be any of its exact
interface marginals for which
\(F\subset\mathcal S\). In particular
the central parent may be integrated,
together with compatible other parent
sets disjoint from \(F\). Take a fixed
smaller cylinder \(\mathcal D_F\) and
a cutoff \(\chi(z_F)\) supported in
\(\mathcal D_F^+\), equal to one on
\(\mathcal D_F\). Put
\(M_\chi=\sup|\nabla\chi|^2\).
The visible event \(D_L\) constrains
only the coordinates in \(F\); every
other retained coordinate is unrestricted.

\begin{theorem}[Volume-uniform local interface repair]
\label{thm:f13v14-floor}
There exist \(C_{\rm f}>0\) and
\(\beta_0<\infty\), independent of
\(L\ge32\), such that the original
rate-one interface heat bath satisfies
\begin{equation}\label{eq:f13v14-floor}
 \kappa_{D_L}^{H_{\nu_{\beta,L}}}
 \ge\frac{c_0\beta}
        {2\{c_0\beta+8C_{\rm f}(\beta+1)\}}
 \ge\frac{c_0}{2(c_0+16C_{\rm f})}>0,
           \qquad c_0=1/84,\quad \beta\ge\beta_0 .
\end{equation}
The condition \(F\subset\mathcal S\)
is part of the theorem. No hidden
configuration is excluded.
\end{theorem}

\begin{proof}
The local measured curvature just
proved gives a diffusion floor
\(\lambda_\beta=c_0\beta\) for
functions supported in the larger
chart cylinder, uniformly in all
uncharted variables.

The global unintegrated potential is
a sum of smooth fixed-stencil Wilson
and inverse-pivot terms. The number
of terms touching a fixed free group,
their compact derivative bounds,
and their source ranges are uniform
in \(L\). Consequently the mixed
Hessian row bound in f12 (81.2) can
be chosen
\(K_\beta\le C_{\rm f}(\beta+1)\)
with \(C_{\rm f}\) independent of
ambient volume.

The localization and global resolvent
argument in the proof of
Theorem~\ref{thm:f13v13-floor} now
applies without alteration:
if \(c_0\beta\ge8M_\chi\), it gives
the first bound in
\eqref{eq:f13v14-floor} for the
full fine-coordinate heat bath
killed outside the cylinder over
\(F\). The cutoff depends on a
fixed number of coordinates, so
\(M_\chi\) is volume independent.
The underlying diffusion and its
resolvent are global operators on
the same law throughout.

The event constrains only visible
coordinates because \(F\subset
\mathcal S\). Lifting a function
of the interface leaves every
hidden group unrestricted.
Exact conditional-variance
domination transfers the fine-law
floor to the actual interface
heat bath with no posterior or
dimension factor. Finally
\(\beta+1\le2\beta\) for
\(\beta\ge1\) yields the second
inequality.
\end{proof}

This improves the geometric scope of
V13, not its status as a local theorem.
The event still prescribes a central
pattern on a large finite collar.
No assertion is made that these
cylinders cover all defects, have
large equilibrium probability, or
already imply an unrestricted
Poincare inequality.

\subsection{Verification and next checkpoint}

The exact verifier is
\begin{center}\small
\path{scripts/verify_ym_f13_v14_nonwrapping_shell.py}.
\end{center}
It uses integer-lattice shifts only,
enumerates all distinct incident
plaquettes, checks source exclusion
of hidden groups, constructs the
free-coordinate input closure, and
verifies the stated embedding bounds.
All curvature budgets and metric
normalizations are rational. The
probabilistic and operator consequences
are proved above; they are not inferred
from tests at sampled hidden fields.

The checker's SHA--256 is
\begin{center}\scriptsize\ttfamily
C29F61721DE134C5CF6C927F56D6D7CD654E3294E6D2B76B4CDB280CF5BBAED4.
\end{center}
Predecessor V13 contains 234,980 bytes
and SHA--256
\begin{center}\scriptsize\ttfamily
13D8839FDE37EE725D40820E464E1913DB850206C6E403B194FD9635683D6B69.
\end{center}
Its full body is preserved, with
only the title version changed.
Only active V13 is replaced by V14;
archives and companion files remain
unchanged.

V15 is the required companion-file
checkpoint. The mathematical task
after this extension is coverage
beyond the prescribed central collar,
or a verified local-to-global
assembly using the acquired repair
cylinders. Counting a periodic
boundary as absent is no longer
an available shortcut: its positive
contribution has been included
explicitly in the non-wrapping
certificate.

\begin{firewall}
V14 proves a volume-uniform local
killed interface heat-bath floor
for an explicitly embedded stencil
at fixed coarse field \(C=I\), with
an explicit visibility condition
on its finite collar. It does not
prove the full compact-interface
gap, a cofinal continuum theory,
or the physical Yang--Mills mass
gap. The full goal remains
unproved and active.
\end{firewall}
\section{V15: companion checkpoint and locally curved coarse backgrounds}
\label{sec:f13v15}

\subsection{Checkpoint and the new acquisition}

The active companion files were re-inventoried, their
terminal conclusions reread, and their hashes checked.
They are unchanged from V10:
\begin{center}\small
\begin{tabular}{lrr}
\toprule
Programme&Active version and bytes&SHA--256 prefix\\
\midrule
SM--Einstein&V72,\ 607372&\texttt{F44F9745D43379E1}\\
NS stability&V26,\ 278283&\texttt{82C887F8C2C69E4F}\\
Parallel YM&V5,\ 54174&\texttt{306F1BB84CFA8A8D}\\
\bottomrule
\end{tabular}
\end{center}
SM V72's fixed-mode many-copy limit is not an
ultraviolet YM construction. NS V26 supplies no
compact-interface repair estimate. Parallel YM
V5's extensive-charge obstruction remains a
reason not to impose global sector tightness
where only local control is needed. Its internal
stopping language does not override this task.
No companion result closes the global YM goal.

The new result here extends the specific V14
non-cleared repair cylinder from \(C=I\) to
coarse fields with small plaquette curvature
on a finite patch. It is an application of
local path reordering and the programme's
gauge equivariance, not a claim of a new
general gauge-fixing principle. The explicit
patch and its path bound are checked, and
the original fixed-coarse heat bath is
transported exactly.

The gauge is chosen from \(C\) alone and is
then held fixed. No fine-field-dependent
gauge fixing, change of the update clock,
or coarse resampling is used. Global
holonomies need not be small or trivial.

\subsection{The finite coarse patch}

Use V14's unshifted integer coordinates,
then the same even translation when
embedding into a torus. Let
\[
                  \mathcal B=\{-2,-1,\ldots,5\}^4
\]
be a coarse vertex box, of seven edges
per side. The relevant coarse inputs
consist of the 117 selected source-edge
outputs and the coarse inputs of every
pivot appearing in a plaquette that
touches the 819-link shell. Their union
has 550 coarse edges. The exact checker
verifies that both endpoints of every
one of them lie in \(\mathcal B\).
Denote this finite input set by \(K\).

For the V14 embedding translated by
\((8,8,8,8)\) fine units, the coarse box
is translated by \((4,4,4,4)\). It lies
inside every coarse torus of side
\(L/2\ge16\), without identifications.
Only plaquettes entirely in this
contractible box enter the following
gauge estimate.

For an oriented coarse edge \(e=(v,\mu)\),
write \(C_e\in\SU(2)\), and let \(Q_p(C)\)
be the oriented elementary plaquette
holonomy. Assume
\begin{equation}\label{eq:f13v15-curvature-input}
             \max_{p\subset\mathcal B}d_{\rm HS}(I,Q_p(C))
                           \le\varepsilon .
\end{equation}
No bound on the sum over all coarse
plaquettes is imposed, and no condition
on plaquettes outside the box is used.

\begin{lemma}[A finite-box link bound]
\label{lem:f13v15-gauge}
There is a coarse gauge \(h_v\), depending
only on the links of \(C\) in \(\mathcal B\),
such that
\begin{equation}\label{eq:f13v15-link-bound}
 C^h_{v,\mu}=h_vC_{v,\mu}h_{v+e_\mu}^{-1},
 \qquad
             d_{\rm HS}(I,C^h_{v,\mu})\le21\varepsilon
\end{equation}
for every edge inside \(\mathcal B\).
When \(\varepsilon=0\), all these
transformed edges equal \(I\) exactly.
\end{lemma}

\begin{proof}
Translate the lower corner of the box
to zero, so its vertices are
\(\{0,\ldots,7\}^4\). For a vertex
\(v\), take the monotone path
\(p(v)\) that makes all direction-0
steps first, then direction 1, then
2, then 3. Let \(P(v)\) be its ordered
link product and set \(h_v=P(v)\).

For an edge \((v,\mu)\), compare
the transport on \(p(v)\) followed
by this edge with the transport
on \(p(v+e_\mu)\). Move the final
direction-\(\mu\) step leftward
past the later-direction steps.
Exactly
\[
                      \sum_{\nu>\mu}v_\nu\le3\cdot7=21
\]
adjacent path swaps are needed.
Each swaps two sides of one
elementary plaquette for the other
two sides. Both intermediate paths
stay inside the box.

The distance between the two path
products at a single swap equals
the distance of that plaquette
holonomy, or its inverse, from
identity. Common prefixes and
suffixes do not change this
distance, by bi-invariance.
Triangle inequality therefore gives
\[
 d_{\rm HS}(P(v)C_{v,\mu},P(v+e_\mu))
                                    \le21\varepsilon .
\]
Right multiplication by
\(P(v+e_\mu)^{-1}\) proves
\eqref{eq:f13v15-link-bound}.
No commutativity or infinitesimal
curvature expansion is used.
\end{proof}

Extend \(h\) arbitrarily to the
remaining coarse vertices, for
example by identity outside the
box. This gives a legitimate
global lattice gauge. It need
not make coarse links outside
the box close to identity.

\subsection{Exact gauge transport of the fixed-coarse fibre}

For a fine link from \(x\) to
\(x+e_\mu\), let
\(r=\lfloor x/2\rfloor\) and
\(s=\lfloor(x+e_\mu)/2\rfloor\)
be its endpoint coarse cells.
Use the cellwise constant fine
gauge
\begin{equation}\label{eq:f13v15-fine-gauge}
                         W^h_{x,\mu}=h_rW_{x,\mu}h_s^{-1}.
\end{equation}
An internal tree link has \(r=s\),
so its value \(I\) is preserved.
Thus this transformation preserves
the programme's tree-fixed free
coordinate carrier.

\begin{proposition}[Gauge covariance without a new Jacobian]
\label{prop:f13v15-conjugacy}
For each fixed coarse field \(C\),
\eqref{eq:f13v15-fine-gauge} is a
product Haar-preserving isometry
of the independent fine groups,
taking the completed fibre over
\(C\) to the completed fibre over
\(C^h\). It transports the full
conditional measure exactly:
\[
                         (\mathcal T_h)_*\pi_{\beta,C}
                                          =\pi_{\beta,C^h}.
\]
It also conjugates their original
single-free-group heat baths and
any exact interface marginals
with the same index partition.
\end{proposition}

\begin{proof}
Each chronological path over a
coarse edge from \(r\) to \(s\)
transforms by \(P_j\mapsto
h_rP_jh_s^{-1}\); the root pivot
transforms the same way. Hence
\[
 (B^h)^{-1}P_j^h
                 =h_s(B^{-1}P_j)h_s^{-1}.
\]
The truncated logarithm is a
conjugation-equivariant radial
function. Its exponential and
the full forward pivot map
therefore transform as
\[
 T_{P^h}(h_rBh_s^{-1})
                        =h_rT_P(B)h_s^{-1}.
\]
Uniqueness of the completed
pivot gives covariance of its
inverse. Left and right group
translations preserve Haar and
the Hilbert--Schmidt metric, so
the corresponding inverse Haar
Jacobian is unchanged. All
Wilson plaquette traces are
also unchanged.

On the independent free groups
the transformation is simply
the product of these fixed left
and right multiplications.
The gauge \(h\) depends on the
fixed parameter \(C\), not on
the integrated fine coordinates.
Consequently there is no extra
field-dependent determinant.
The full positive density and
its normalization are carried
to the density at \(C^h\).

Each transformed free coordinate
depends only on that same free
coordinate and on fixed matrices.
The conditioning sigma fields
for a single-coordinate heat
bath are therefore transported
exactly. The same holds after
integrating any declared hidden
index set. Norms, Dirichlet forms,
and the rate-one clock are
preserved.
\end{proof}

This is an identification between
two gauge-equivalent fixed-coarse
fibres for purposes of proof.
The chain in the original
variables continues to keep its
original \(C\) fixed. In
particular it does not use
V11's variable-coarse update.

\subsection{Stability of the particular repair certificate}

We first state the local parameter
step, before using the gauge.
The action derivative in V14's
selected direction depends only
on its finite fine input closure
and on the coarse edges in \(K\).
On that finite compact parameter
space, the completed action is
jointly smooth. The same is true
of its required derivatives.

At central \(F\)-data and
\(C_e=I\) for \(e\in K\), the
V14 unit-direction bound is
\(-4/21\), uniformly over all
uncontrolled fine groups and
all irrelevant coarse data.
Uniform continuity gives fixed
\(\delta_*>0\) and \(r_*>0\)
such that
\begin{equation}\label{eq:f13v15-open}
 \partial_a^2S\le-2/21
\end{equation}
whenever
\[
 \max_{e\in K}d_{\rm HS}(I,C_e)\le\delta_*,
 \qquad
 \max_{\ell\in F}|z_\ell|\le r_*,
\]
where \(z_\ell\) are the
Hilbert--Schmidt normal
coordinates about
\(\epsilon_\ell I\). Choose
the radii strictly inside an
open neighbourhood on which
this inequality holds.

These constants are independent
of volume because the entire
parameter closure is fixed.
The coarse field outside \(K\)
does not enter this directional
action derivative. No global
closeness of \(C\) to identity
has been assumed in this step.

Shrink the radii if necessary.
The 117 selected source families
then stay near their negative
central products, while \(C_e\)
is near \(I\). All their
cutoffs are inactive at the
candidate pivot \(B=C_e\);
the strict distance margin to
the cutoff support persists.
Thus the actual pivots are
\(B=C_e\), with inverse Haar
Jacobian one, and are independent
of the selected free coordinates.
No other pivot contains these
boundary groups in a source.
The amplitude cancellation in
the selected direction therefore
persists for these curved
coarse inputs.

The coordinate-Haar derivative
is uniformly bounded on the
fixed chart. For large \(\beta\),
\eqref{eq:f13v15-open} yields
\(\partial_a^2\Phi_\beta\le
-\beta/21\), with inverse-metric
lower bound \(1/2\). Hence the
same diffusion coefficient
\(c_0=1/84\) and the same
type of volume-uniform
original-clock comparison as
in V14 apply. Neither
\(\delta_*\) nor \(r_*\) is
numerically enclosed here.

\subsection{The locally curved-background repair theorem}

Put \(\varepsilon_*=\delta_*/21\),
decreasing it if necessary to
meet the strict inactive-source
margins just described. For
each \(C\) satisfying
\eqref{eq:f13v15-curvature-input}
with \(\varepsilon\le
\varepsilon_*\), fix the
gauge \(h(C)\) of
Lemma~\ref{lem:f13v15-gauge}.

Define the visible repair
cylinder \(D_C\) by requiring
the transformed free groups
\(Y_\ell^h\), \(\ell\in F\),
to lie in a fixed smaller
product chart about
\(\epsilon_\ell I\).
No other fine group is
restricted. The cylinder
depends on the fixed coarse
input through \(h(C)\).
In original variables its
central free-link values are
\[
                   Y_\ell=\epsilon_\ell h_r^{-1}h_s,
\]
with endpoint cells as in
\eqref{eq:f13v15-fine-gauge}.
This prescription need not
be globally near \(I\).

\begin{theorem}[Local repair at small coarse curvature]
\label{thm:f13v15-floor}
Assume the non-wrapping embedding
of V14, \(F\subset\mathcal S\),
and the local coarse curvature
bound with
\(\varepsilon\le\varepsilon_*\).
For the exact fixed-\(C\)
interface marginal, there are
\(C_{\rm f}>0\) and
\(\beta_0<\infty\), uniform in
ambient volume and in all
coarse inputs meeting this
local condition, such that
\[
 \kappa_{D_C}^{H_{\nu_{\beta,C}}}
 \ge\frac{c_0\beta}
        {2\{c_0\beta+8C_{\rm f}(\beta+1)\}}
 \ge\frac{c_0}{2(c_0+16C_{\rm f})}>0,
 \qquad c_0=1/84,\quad\beta\ge\beta_0 .
\]
All hidden variables and all
unprescribed fine-interface
variables remain unrestricted.
\end{theorem}

\begin{proof}
The box gauge gives
\(\max_{e\in K}d(I,C_e^h)
\le21\varepsilon\le\delta_*\).
The stable measured curvature
estimate therefore holds in
the fixed transformed chart,
uniformly over all fine data
outside it. Apply the V14
localization, global diffusion
resolvent comparison, and
exact marginal form domination
to the transformed law.

The global mixed-Hessian row
constant can be chosen uniformly
in \(C\): each local completion
is smooth on a fixed compact
set of group inputs, and each
free coordinate meets a bounded
number of such terms. The
coarse values outside the patch
may be large but remain compact
group parameters. The chart
and cutoff constants are fixed
in the transformed variables.

Finally
Proposition~\ref{prop:f13v15-conjugacy}
transports the measure, visible
cylinder, and the original
heat-bath operators back to
the original fixed-\(C\)
fibre isometrically. The
stated killed floor and
its constants are unchanged.
\end{proof}

If the coarse field is flat
on the patch, the gauge makes
its patch links exactly
identity. This does not
force noncontractible
holonomies of the surrounding
torus to vanish. More
generally the theorem needs
only local plaquette control,
not a volume-dependent sum
of plaquette errors or a
global link-norm bound.

\subsection{Verification and what is still missing}

The checker is
\begin{center}\small
\path{scripts/verify_ym_f13_v15_local_gauge.py}.
\end{center}
It verifies containment of
all 550 required coarse inputs
in the patch; exhaustively
checks the path reordering
for all 14336 box edges,
with maximum 21 plaquette
swaps; and tests 96 exact
noncommutative plaquette
covariance identities.
It also checks exact local
trivialization of rational
pure-gauge fields. Those
finite algebraic regressions
do not replace the general
gauge-covariance proof above.

Its SHA--256 is
\begin{center}\scriptsize\ttfamily
76B508692D5AE1641279FB95CFE619F8592AEDFF740FD11F937B6F05D0C8D263.
\end{center}
The predecessor V14 has
249,765 bytes and SHA--256
\begin{center}\scriptsize\ttfamily
CBC7DC8857049BD001B568625D2E0D91EE442BE2DEB0D0527602784EA9858DF7.
\end{center}
Its complete body is retained
apart from the title version.
Active V15 replaces V14;
archives and companion notes
are unchanged. The next
companion checkpoint and
direction review are V20.

\begin{firewall}
The repair theorem now applies
to a local gauge orbit of
central-interface patterns
over a small-curvature coarse
patch, with constants uniform
in surrounding volume.
It remains conditional on
the specified visible
fine-interface pattern.
There is no covering theorem
for all bad interface data,
no unrestricted spectral
gap, and no continuum
Yang--Mills construction or
physical mass-gap proof.
The full goal remains active.
\end{firewall}
\section{V16: a noncentral visible-curvature certificate on the fixed shell}
\label{sec:f13v16}

\subsection{Direction and nonduplication}

V14 fixed a central pattern and V15 transported a
neighbourhood of it to mildly curved coarse backgrounds.
Here we replace that prescribed pattern by a test of
the actual visible plaquettes. The test is a symmetric
\(3\times3\) matrix together with explicit inactivity
conditions for the selected pivots. It applies to
arbitrary compact coarse data whenever those tests pass.
It does not assert that they pass on every bad field.

The full-fibre zero-action uniqueness question was
checked upstream: f12 V13 already proves the corresponding
section statement for its compact block, with extensions
and corrections tracked in that lineage. It is not a
new acquisition here. Nor do we reprove the abstract
diffusion-to-heat-bath comparison of f12 V77--V80.
The new ingredients are the noncentral shell matrix,
its all-hidden-field implication, and a fixed-patch
compact covering argument for the entire certified event.
The latter is not a spatial union over growing volume,
the operation obstructed in f12 V79.

\subsection{Visible inputs and a strict source margin}

Keep V14's non-wrapping shell \(T\), its collar \(F\),
and the same embedding into even fine side length
\(L\ge32\). Thus \(|T|=819\), \(|F|=10034\), and
the coarse input set \(K\) of V15 has 550 edges.
The coarse field \(C\) remains fixed by every update.
The hidden index set is arbitrary subject to
\(F\subset\mathcal S\), where \(\mathcal S\) is the
retained fine-interface index set.

Let \(\mathcal P_1\) be the 1140 fixed one-selected
plaquettes and \(\mathcal P_2\) the 1845 two-selected
plaquettes of V14. Here \emph{fixed} means determined
by \((C_K,Y_F)\), not central or constant as \(Y_F\)
varies. The remaining 84 incident plaquettes have one
selected link each and will be bounded without
restricting any of their unprescribed inputs.
The definition of \(F\) includes the free input
closure of every plaquette in \(\mathcal P_1\cup
\mathcal P_2\); no completed pivot is treated as an
independent visible coordinate.

Write \(\rho=1/64\) for the chordal cutoff radius.
For a fixed number \(\sigma>0\), impose
\begin{equation}\label{eq:f13v16-margin}
 d_{\rm HS}(I,C_e^{-1}P_j(Y_F))\ge\rho+\sigma
 \quad\text{for all fourteen nonroot sources of
 each of the 117 selected edges }e.
\end{equation}
All these source paths use only the stated visible
collar. At candidate pivot \(B=C_e\), every truncated
logarithm vanishes on an open neighbourhood of these
data. Uniqueness of the actual completed pivot gives
\(B=C_e\), with inverse Haar Jacobian one, throughout
that neighbourhood.

Every selected boundary group belongs only to the
two chronological nonroot paths of its own coarse
edge. Hence, while \eqref{eq:f13v16-margin} persists,
moving the selected groups changes no pivot and no
inverse-pivot Jacobian. This statement is uniform
over all free groups outside \(F\).
It requires no smallness of \(C\) or of the unselected
fine groups.

\subsection{An exact three-dimensional curvature matrix}

Identify \(\SU(2)\) with unit quaternions
\((q_0,\mathbf q)\), with
\(\operatorname{Tr}q=2q_0\).
For a unit vector \(n\in\mathbb R^3\), write
\(E_n=(0,n)\) and \(R_n(t)=\exp(tE_n)\).
Left multiply every selected free link by the same
\(R_n(t)\), leaving all other free groups fixed.
The squared Hilbert--Schmidt speed is
\begin{equation}\label{eq:f13v16-speed}
                         N_T=2|T|=1638.
\end{equation}
Use the inherited action normalization
\(s_p=2-\operatorname{Tr}Q_p=2-2(Q_p)_0\).

For \(p\in\mathcal P_1\), let \(q_{p,0}\) be the
scalar part of its actual completed holonomy.
For \(p\in\mathcal P_2\), cyclically start the
oriented plaquette word at its positive selected
link. It has the form
\[
                    U A V^{-1}B=X B,
                    \qquad X=UAV^{-1},
\]
where \(U,V\) are the two selected positive link
variables and \(A,B\) are the intervening oriented
unselected links. Cyclic rotation does not change
the trace. Put \(X=(x_0,x)\) and \(B=(b_0,b)\), and set
\begin{equation}\label{eq:f13v16-matrix}
 M_{\rm vis}(C_K,Y_F)
 =2\sum_{p\in\mathcal P_1}q_{p,0}I_3
   +8\sum_{p\in\mathcal P_2}
       \left\{\frac{x_pb_p^{\mathsf T}
                         +b_px_p^{\mathsf T}}2
                    -(x_p\cdot b_p)I_3\right\}.
\end{equation}
This is a smooth function of the finite compact
input carrier. It is a criterion in the fixed tree
coordinates; no invariance under independent
cell gauges of a common axis \(n\) is asserted.

\begin{lemma}[Exact noncentral curvature]
\label{lem:f13v16-matrix}
Under \eqref{eq:f13v16-margin}, for every unit \(n\)
and every choice of all unprescribed fine fields,
\begin{equation}\label{eq:f13v16-upper}
       \left.\frac{d^2}{dt^2}S(Y(t))\right|_{t=0}
                       \le n^{\mathsf T}M_{\rm vis}n+168.
\end{equation}
In particular, if
\begin{equation}\label{eq:f13v16-test}
                 \lambda_{\min}(M_{\rm vis})
                                  \le-168-\eta,\qquad\eta>0,
\end{equation}
there is a direction of unit Hilbert--Schmidt
norm with action Hessian at most \(-\eta/1638\),
uniformly over all unprescribed groups.
\end{lemma}

\begin{proof}
The source margin makes every other physical link
constant under the selected motion. In a plaquette
with one selected link, that link or its inverse has
second derivative equal to its negative. Thus the
plaquette action has second derivative \(2q_{p,0}\),
independently of \(n\). The same reasoning bounds
each of the 84 uncontrolled contributions by 2.

For a two-selected plaquette the moving trace is
\(\operatorname{Tr}(R_n X R_n^{-1}B)\).
Quaternion multiplication gives
\[
 [E_n,[E_n,X]]
            =(0,4\{n(n\cdot x)-x\}).
\]
Consequently its action second derivative is
\[
 -2\bigl([E_n,[E_n,X]]B\bigr)_0
       =8\{(n\cdot x)(n\cdot b)-x\cdot b\},
\]
which is the quadratic form in \eqref{eq:f13v16-matrix}.
Adding the contributions proves
\eqref{eq:f13v16-upper}. Choose a unit eigenvector
for the least eigenvalue and divide by
\eqref{eq:f13v16-speed}. The product of left
one-parameter subgroup curves is a product-metric
geodesic, so the normalized second derivative is
the Riemannian Hessian in that direction.
\end{proof}

A scalar sufficient test, avoiding an eigensolver, is
\[
 \frac{\operatorname{Tr}M_{\rm vis}}3
 =2\sum_{\mathcal P_1}q_{p,0}
       -\frac{16}{3}\sum_{\mathcal P_2}x_p\cdot b_p
                         \le-168-\eta.
\]
It is weaker than the matrix test, not equivalent to it.
At V14's central collar and \(C_K=I\), the pair
matrices vanish and
\[
 M_{\rm vis}=2(450-690)I_3=-480I_3,\qquad
                       M_{\rm vis}+168I_3=-312I_3.
\]
Thus the event is nonempty for \(0<\eta<312\)
and sufficiently small positive \(\sigma\).
Strict inequalities persist under noncentral
perturbations. No commuting-field assumption
was used in the general matrix identity.

\subsection{One fixed-patch covering, including all certified fields}

For fixed \(\sigma,\eta>0\), define \(D_C^{\sigma,\eta}\)
to be the visible cylinder specified by
\eqref{eq:f13v16-margin} and \eqref{eq:f13v16-test},
with no restrictions on any other free groups.
Its parameter set
\[
 \mathcal A_{\sigma,\eta}\subset
                     \SU(2)^K\times\SU(2)^F
\]
is compact, since its defining inequalities are closed.
It may be empty for some \(C\); the following
form statement then remains valid trivially.

\begin{lemma}[Uniform localization of the certified event]
\label{lem:f13v16-localize}
Let \(L_{\beta,C}\) be the full joint reversible
product-metric diffusion, with Dirichlet form
\(\mathcal E_L\) and actual completed-fibre law \(\mu\).
There are \(M<\infty\) and a finite lower threshold
for \(\beta\), independent of ambient volume and
of all coarse data, such that
\begin{equation}\label{eq:f13v16-localize}
 \|1_{D_C^{\sigma,\eta}}u\|_{L^2(\mu)}^2
 \le\frac{2}{c_\eta\beta}\mathcal E_L(u,u)
             +\frac{2M}{c_\eta\beta}\|u\|_{L^2(\mu)}^2,
             \qquad c_\eta=\frac{\eta}{26208}.
\end{equation}
Constants may depend on \(\sigma,\eta\) and on this
fixed shell.
\end{lemma}

\begin{proof}
At each point of \(\mathcal A_{\sigma,\eta}\), select
one least-eigenvalue axis, without assuming a continuous
choice across eigenvalue crossings. In orthonormal
product normal coordinates about its \(Y_F\)-component,
freeze the corresponding Euclidean unit direction.
At the centre its action second derivative is at most
\(-\eta/1638\) for all unprescribed inputs.

The action derivatives involve a fixed finite compact
input closure, including the uncontrolled plaquettes.
Continuity and the strict source margin give an open
neighbourhood in the \(C_K,Y_F\) carrier where the
frozen directional second derivative is at most
\(-\eta/(2\cdot1638)\), for all those inputs.
Shrink the coordinate charts so the inverse metric is
at least \(1/2\). All inverse-pivot amplitudes are
independent of the selected coordinates on these
neighbourhoods. The coordinate-Haar directional
Hessians are uniformly bounded on smaller charts.
For sufficiently large \(\beta\), the measured
potential therefore has second derivative at most
\(-\beta\eta/(4\cdot1638)\).

The local ground-state conjugation estimate used
in V14 gives, for any function \(v\) supported in
one such chart cylinder,
\[
                   \mathcal E_L(v,v)
                              \ge c_\eta\beta\|v\|^2,
       \qquad c_\eta=\frac{\eta}{16\cdot1638}.
\]
The factor \(1/16\) includes the two continuity
and amplitude allowances, the inverse-metric factor
\(1/2\), and the conjugation factor \(1/2\).

Compactness gives finitely many such product
neighbourhoods, with smooth nonnegative cutoffs
\(\psi_j(C_K,Y_F)\) supported in them and with
\(T=\sum_j\psi_j^2\ge1\) on
\(\mathcal A_{\sigma,\eta}\). Choose a smooth
nonnegative function \(g\) equal to 1 for \(T\le1/4\)
and equal to 0 for \(T\ge1/2\), and put
\[
                  \chi_j=\frac{\psi_j}{\sqrt{T+g(T)^2}}.
\]
The denominator is everywhere positive. Thus
\(\sum_j\chi_j^2\le1\) everywhere and equals 1
on the certified set. Their derivatives with respect
to the fine variables have the finite uniform bound
\[
                        M=\sup_{C_K,Y_F}\sum_j|\nabla\chi_j|^2.
\]
No coarse derivative enters the diffusion: \(C\)
is fixed. This construction avoids differentiating
any selected eigenvector.

Apply the local estimate to \(\chi_j u\), sum,
and use the product rule and Cauchy--Schwarz:
\[
 c_\eta\beta\|1_Du\|^2
 \le\sum_j\mathcal E_L(\chi_j u,\chi_j u)
 \le2\mathcal E_L(u,u)+2M\|u\|^2.
\]
The inequality first holds on a smooth form core
and extends by form closure. All supports constrain
only \(F\). Finiteness of the cover and the input
closure, not compactness in a growing number of
coordinates, proves the stated uniformity.
\end{proof}

\subsection{Original-clock repair for the whole certified event}

\begin{theorem}[Matrix-certified interface repair]
\label{thm:f13v16-repair}
For the exact fixed-\(C\) interface marginal \(\nu\),
with \(F\subset\mathcal S\), there are
\(C_{\rm f}<\infty\) and \(\beta_0<\infty\),
uniform in even \(L\ge32\), arbitrary compact
coarse data, and compatible hidden index sets, such that
\[
 \kappa_{D_C^{\sigma,\eta}}^{H_\nu}
 \ge\frac{c_\eta\beta}
           {2\{c_\eta\beta+8C_{\rm f}(\beta+1)\}}
 \ge\frac{c_\eta}{2(c_\eta+16C_{\rm f})}>0,
 \qquad \beta\ge\beta_0 .
\]
The heat bath updates each retained fine group at
rate one, with its actual marginal conditional law.
\end{theorem}

\begin{proof}
Write \(L=L_{\beta,C}\). The inherited global
comparison on the full joint fibre is
\[
 H_\mu\succeq L(L+K_\beta)^{-1},
                  \qquad K_\beta=C_{\rm f}(\beta+1).
\]
The mixed-Hessian row constant is uniform in all
coarse fields by smoothness on finite compact
completion stencils and bounded incidence.
No global diffusion gap is being assumed.

Put \(\lambda=c_\eta\beta\), and choose
\(\beta_0\) large enough that \(\lambda\ge8M\).
For the spectral projection \(\Pi_t\) of \(L\)
onto \([0,t]\), \eqref{eq:f13v16-localize} gives
\[
            \|1_D\Pi_t\|^2\le2(t+M)/\lambda.
\]
Taking adjoints gives the same bound for
\(\Pi_t1_D\); no commutation is asserted.
With \(t=\lambda/8\), a function supported in
\(D\) has at least half its squared norm above
that spectral interval. The global comparison yields
\[
 \mathcal E_{H_\mu}(f,f)
                  \ge\frac{\lambda}{2(\lambda+8K_\beta)}
                                      \|f\|_{L^2(\mu)}^2
                  \quad\text{if }f=1_Df .
\]
This is a global resolvent, not a resolvent of a
killed or chart-restricted diffusion.

Finally lift a marginal function by \(Jf=f(Y_{\mathcal S})\).
Norms are preserved. Hidden-coordinate updates leave
\(Jf\) unchanged, and conditional variance gives
\(\mathcal E_{H_\nu}(f,f)\ge
\mathcal E_{H_\mu}(Jf,Jf)\), as in V14.
The lifted support is precisely a cylinder with
all hidden variables allowed. Apply the joint
bound and use \(\beta+1\le2\beta\) for \(\beta\ge1\).
\end{proof}

\subsection{Verification and the remaining coverage question}

The new checker is
\begin{center}\small
\path{scripts/verify_ym_f13_v16_visible_curvature.py}.
\end{center}
It performs 864 exact rational noncommutative
second-derivative checks, including mixed axes,
both single-link orientations, and the two-link
commutator formula. It reruns the V14 shell audit
and recovers the matrices \(-480I_3\) and \(168I_3\)
and the speed 1638. It does not numerically enclose
the compact-cover radii, \(M\), or \(\beta_0\).
The finite checks supplement, not replace, the
algebraic and analytic proofs.

Its SHA--256 is
\begin{center}\scriptsize\ttfamily
F6C4FE4BFAAA76304FF0DCC12384BDA7D57F543017AAEABFA39AAA978F9A3AED.
\end{center}
The predecessor V15 has 263,760 bytes and SHA--256
\begin{center}\scriptsize\ttfamily
006599263CD9A4BD642E0F7C2ECB0EBCB0EE50FAA54AD87DE4523E140E5FA9C8.
\end{center}
Its body is retained apart from the title version.
Only active V15 is replaced; archives and companions
are unchanged. The next companion checkpoint remains V20.

\begin{firewall}
The certified event now contains all configurations
passing an actual noncentral visible matrix test on
this fixed shell, rather than only a preselected
central chart. The theorem controls their union by
one finite-dimensional covering and leaves all hidden
variables unrestricted. It does not cover fields
with near-active selected sources or an inconclusive
matrix test. It does not justify a growing spatial
union, unrestricted interface gap, cofinal iteration,
continuum construction, or physical mass gap.
The next substantive coverage obligation is to
control those complementary source and curvature
cases, not to equate this sufficient test with a
classification of all bad fields.
\end{firewall}
\section{V17: repair certificates with genuinely active pivot sources}
\label{sec:f13v17}

\subsection{The complementary case being acquired}

V16's source margin excludes any selected source in the
inner logarithm region. We now exhibit a family with such
sources and include the exact moving-pivot response in
the proof. The response and its acceleration are not
discarded. Nor is the actual inverse-Jacobian amplitude
replaced by a constant. The construction remains on one
fixed non-wrapping patch, with all hidden fields free.

Use \(C=I\), the 819-link shell \(T\), and the
10034-group visible collar \(F\) from V14.
Choose a subset \(R\subset T\), and reverse the central
signs on exactly those links. Thus selected links in
\(R\) equal \(+I\), selected links outside \(R\) equal
\(-I\), and all other groups in \(F\) retain their V14
central values. Let \(k_e\in\{0,\ldots,7\}\) count the
reversed links on a selected coarse edge \(e\).
Only \(F\) is prescribed; all other free coordinates
are arbitrary throughout.

This is not a reclassification of the hidden stationary
orbits of V8, and no global ground-state assertion follows.
The new acquisition is an explicit set of active-source
interface patterns that V16's condition does not admit.

\subsection{The actual pivot motion is exactly solvable}

\begin{lemma}[Active-source speed and acceleration]
\label{lem:f13v17-pivot}
Left multiply every link in \(T\) by
\(R_n(t)=\exp(tE_n)\), for any unit \(n\in\mathbb R^3\).
For sufficiently small \(|t|\), the completed pivot on
each selected coarse edge is exactly
\begin{equation}\label{eq:f13v17-pivot}
 B_e(t)=\exp(-a_e tE_n),\qquad
          a_e=\frac{k_e}{8-k_e}.
\end{equation}
In particular \(B'_e(0)=-a_eE_n\) and
\(B''_e(0)=-a_e^2I\). Every other pivot is independent
of this motion. These statements hold for all
unprescribed fine fields.
\end{lemma}

\begin{proof}
Each of the seven selected boundary links occurs
in two nonroot paths of its own coarse edge, and
in no other source family. Their accompanying internal
links have central value \(I\) in the prescribed collar.
Hence the fourteen nonroot paths consist of \(2k_e\)
copies of \(R_n(t)\) and \(14-2k_e\) copies of
\(-R_n(t)\). They contain no hidden inputs.

At \(t=0\), every truncated logarithm vanishes at
\(B=I\): a positive source has logarithm zero, while
a negative source is outside the support. Uniqueness
identifies the actual pivot with \(I\).
For small \(|t|\), at the candidate
\(B=\exp(-a_etE_n)\), positive relative sources
are \(\exp((1+a_e)tE_n)\) in the inner region
where the cutoff is identically one; negative
relative sources remain outside the support.
The forward pivot equation therefore reduces to
\[
 -a_etE_n+\frac{2k_e}{16}(1+a_e)tE_n=0.
\]
It is satisfied precisely by \(a_e=k_e/(8-k_e)\).
All exponents commute in this calculation and the
small logarithm is unambiguous. Global uniqueness
of the completed inverse then proves
\eqref{eq:f13v17-pivot}. There are only eight
occupancies, so a common positive interval in \(t\)
can be chosen. No selected group enters another
coarse edge's source family.
\end{proof}

The inner sources are active in the sense that their
cutoff is one and their first derivative contributes.
Their logarithm happens to vanish at the central point;
that does not make their response zero.

\subsection{A uniform debit for the pivot terms}

For \(a\ge0\) define
\[
 d(a)=\max\{2a^2+4a,\ 4a^2\},\qquad
 \ell(k)=24k+6d\!\left(\frac{k}{8-k}\right),\qquad
                         \mathcal L(R)=\sum_e\ell(k_e).
\]
The exact rational values are
\[
\begin{array}{c|rrrrrrrr}
 k&0&1&2&3&4&5&6&7\\ \hline
 \ell(k)&0&1356/49&172/3&2268/25&132&580/3&360&1344 .
\end{array}
\]

\begin{theorem}[Active-shell curvature certificate]
\label{thm:f13v17-curvature}
At the prescribed \(R\)-pattern, for every unit
axis \(n\) and every choice of all unprescribed groups,
the actual fixed-coarse action satisfies
\begin{equation}\label{eq:f13v17-curvature}
             S''(0)\le-312+\mathcal L(R).
\end{equation}
The squared speed in the independent free-group
product metric remains \(1638\); pivot speeds are
not additional independent coordinates.
\end{theorem}

\begin{proof}
First compute an auxiliary physical-link variation
in which the selected links rotate but every pivot
is artificially held at its central value. This
auxiliary variation is used only to compare derivatives;
it is not asserted to stay in the actual fibre.

All fixed plaquettes of V14 still have central
physical links at \(t=0\). Their nonselected pivots
remain \(I\), because their prescribed sources
are central and both signs give zero truncated
logarithms at \(B=I\). The two-selected fixed
plaquettes cancel the common rotation even after
sign reversals. A one-selected fixed plaquette
changes its second derivative by at most 4 if its
selected link is reversed. Each reversed link
meets six plaquettes. The 84 uncontrolled
one-selected plaquettes still contribute at most
2 each, without comparing their individual signs.
The auxiliary second derivative is therefore at most
\[
                           -312+24|R|.
\]

It remains to add every term caused by moving
pivots. All moving links are parallel (direction 2),
so any plaquette contains at most two of them.
Use unit-quaternion norm for this product-rule
bound: the action is \(2-2q_0\), so a term in its
second derivative has absolute value at most twice
the quaternion norm of the corresponding product.
A pivot of speed magnitude \(a\) has first-jet
norm \(a\) and second-jet norm \(a^2\), including
the acceleration in Lemma~\ref{lem:f13v17-pivot}.
An inverse oriented factor has the same norms.

If there is one moving pivot and one selected
free link, all the new terms have total absolute
value at most \(2a^2+4a\). If the other parallel
link does not move, the bound is \(2a^2\).
If both parallel links are moving pivots, of
magnitudes \(a,b\), there is no selected free
link in that plaquette, and its new terms are bounded by
\[
 2a^2+2b^2+4ab
               \le4a^2+4b^2\le d(a)+d(b).
\]
These estimates use only unitarity of the other
factors, so they hold for uncontrolled plaquettes
and noncommuting unprescribed fields as well.
Each moving pivot has six incident plaquettes.
Summing gives the additional upper bound
\(6\sum_e d(a_e)\), with no omitted cross term.
Together with \(|R|=\sum_e k_e\), this proves
\eqref{eq:f13v17-curvature}.
\end{proof}

For eleven sign reversals on distinct coarse edges,
\[
 S''(0)\le-312+11\frac{1356}{49}
                         =-\frac{372}{49}<0.
\]
For five sign reversals on one edge,
\[
 S''(0)\le-312+\frac{580}{3}
                         =-\frac{356}{3}<0.
\]
The latter has ten inner nonroot sources on that
edge. Occupancies six and seven fail this sufficient
debit test on their own; failure of the bound does
not prove absence of a different descent direction.

\subsection{Noncentral neighbourhoods with the actual amplitude}

Fix \(\eta>0\), and let \(\mathfrak R_\eta\) be the
finite collection of nonempty subsets \(R\subset T\)
with \(\mathcal L(R)\le312-\eta\). It is unnecessary
to enumerate this collection to prove the theorem:
the preceding estimates apply to every subset.
For example, \(\eta=372/49\) includes both examples above.

At each such centre, normalize one fixed axis
direction in the independent free-group metric.
Its action Hessian is at most \(-\eta/1638\).
Continuity, with all unprescribed input groups
kept in their finite compact carrier, supplies
a product chart on \(F\) where the frozen
coordinate-direction second derivative is at most
\(-\eta/(2\cdot1638)\). These neighbourhoods
contain noncentral configurations.

The source pattern is stable after shrinking:
the \(2k_e\) positive relative sources stay in
the inner cutoff region of the actual moving
pivot, and the others stay outside its support.
No source is near a cutoff transition.
This supplies open active-source regions, not
just isolated central points.

Unlike V14--V16, the inverse-pivot amplitude is
\emph{not} constant in the selected direction.
The actual measured potential in a chart is
\[
 \Phi_\beta=\beta S-\log j_{\rm inv}-\log h_{\rm Haar}.
\]
Only the 117 selected-edge Jacobian factors can
depend on a selected free coordinate. The completed
map is smooth with positive inverse Haar Jacobian
on its compact input carrier. Their logarithmic
directional derivatives through second order are
therefore bounded by a fixed finite constant.
The same holds for coordinate-Haar terms on
smaller charts. These bounds are independent of
ambient volume and of unprescribed fields.
For sufficiently large \(\beta\), it follows that
\[
                   \partial_a^2\Phi_\beta
                           \le-\frac{\beta\eta}{4\cdot1638}.
\]
Thus the amplitude is absorbed as a bounded
order-one term, not deleted from the law.

There are at most 936 possibly moving physical
links, including the 117 pivots, and 3156 incident
plaquettes. The checker finds that their full free
input closure lies within the unshifted vertex bounds
\[
 [-4,9]\times[-4,9]\times[0,3]\times[-3,9].
\]
The translation by \((8,8,8,8)\) embeds it without
wraparound for every even \(L\ge32\).
This closure is used for uniform derivative bounds;
it does not impose additional visible restrictions
beyond \(F\). The proof above already allows
every other input to be arbitrary.

\begin{theorem}[Active-source repair in the original clock]
\label{thm:f13v17-repair}
Assume \(C=I\) and \(F\subset\mathcal S\).
For each \(\eta>0\), choose sufficiently small
closed inner chart cylinders about the patterns
in \(\mathfrak R_\eta\), and let \(D_\eta^{\rm act}\)
be their union. There are \(C_{\rm f}<\infty\) and
\(\beta_0<\infty\), uniform in even \(L\ge32\)
and compatible hidden index sets, such that
the exact interface marginal satisfies
\[
 \kappa_{D_\eta^{\rm act}}^{H_\nu}
 \ge\frac{c_\eta\beta}
           {2\{c_\eta\beta+8C_{\rm f}(\beta+1)\}}
 \ge\frac{c_\eta}{2(c_\eta+16C_{\rm f})}>0,
 \qquad c_\eta=\frac{\eta}{26208},\quad\beta\ge\beta_0.
\]
For \(\eta=372/49\), \(c_\eta=31/107016\).
The cylinders may be chosen entirely outside
V16's all-sources-inactive event.
\end{theorem}

\begin{proof}
On every outer chart, choose inverse-metric lower
bound \(1/2\). The measured curvature estimate
above gives local diffusion floor \(c_\eta\beta\)
by the same ground-state conjugation as V16.
There are finitely many centres on the fixed
carrier \(F\). The smooth normalized cutoffs of
Lemma~\ref{lem:f13v16-localize}, now with \(C=I\),
therefore give
\[
 \|1_{D_\eta^{\rm act}}u\|^2
       \le\frac{2}{c_\eta\beta}\mathcal E_L(u,u)
               +\frac{2M}{c_\eta\beta}\|u\|^2
\]
with a finite volume-independent \(M\).
This step does not require constant Jacobian;
the actual measured local floor was just proved.
Apply the global resolvent comparison, low-spectral
projection estimate, and exact marginal form
domination in Theorem~\ref{thm:f13v16-repair}.
All hidden variables remain unrestricted by the
lifted cylinder, so no posterior-weight factor
is introduced.

Every centre in \(\mathfrak R_\eta\) has at least
one inner source with \(C_e^{-1}P_j=I\).
Shrink its closed inner chart so that this
source remains at chordal distance less than
\(\rho\) from \(C_e\). It then fails V16's
margin condition for every \(\sigma>0\).
Finiteness permits this shrinkage at every centre.
\end{proof}

\subsection{Verification record and remaining cases}

The exact checker
\begin{center}\small
\path{scripts/verify_ym_f13_v17_active_shell.py}
\end{center}
checks all \(117\cdot128=14976\) local source
occupancies, the twofold chronological multiplicity,
the exact inverse-pivot speed equation, every pair
of possible pivot magnitudes in the debit inequality,
and the non-wrapping input closure. It also reruns
the V14 baseline enumeration. Its SHA--256 is
\begin{center}\scriptsize\ttfamily
FB7CF7119B09D73CAA6F1AAAA14A98B15214740D4DC8DC406B5117A0DD5A8F67.
\end{center}
It does not certify numerical chart radii or
the threshold \(\beta_0\); these are existential
constants in the analytic proof.

The predecessor V16 has 278,605 bytes and SHA--256
\begin{center}\scriptsize\ttfamily
776BF03503C30F65F879A57B961245CE33F7CD72C1D8C2C0A7C8A203C4C28081.
\end{center}
Its body is preserved apart from the title version.
Active V17 replaces V16. Archives and companion
notes are unchanged; the next companion review is V20.

\begin{firewall}
We have acquired explicit active-source repair
regions excluded by V16, with the actual pivot
speed, acceleration, and Jacobian retained.
The certified family still lies near specified
finite-collar sign patterns at identity coarse
data. General transition-cutoff configurations,
large occupancy loads, and other noncentral
patterns remain uncovered. This is neither a
classification of the complement nor an unrestricted
interface spectral gap. The cofinal continuum
construction and physical Yang--Mills mass gap
remain unproved; the full goal stays active.
\end{firewall}
\end{document}
